\chapter{分布式 MPC 与多机编队}
\label{ch:mr-dmpc}

% ════════════════════════════════════════════════════════════════
%  全吸收章：重渲染 Tutorial 笔记
%   - 04_移动机器人规控/50_多机器人协作/50_分布式MPC多足编队.md
%     （集中→分布动机 / 编队三范式入 MPC / 对偶分解 / consensus-ADMM /
%      互联 SRB 网络 + 耦合力 / 分层 MPC-WBC / 拓扑-稳定性）
%  记号归一：R in SO(3)、T in SE(3)（preamble 宏 \SO \SE \so \se）；
%            MPC/优化/图论自然记号保留。
%  源由 AI 生成，作者归属/年份/数值/前沿论文编号存疑处就地 \pz/\rebuilt，
%  并入 claims 台账，待联网核对后二次校验。
%  教学层遵循 v5.0 理论教学规范（动机→反面→历史→理论→陷阱→练习
%   + 误解汇总/小结/速查/故障排查）。
% ════════════════════════════════════════════════════════════════

% —— 本章局部数学算子（章文件 \input 入正文，故用 \providecommand + \operatorname 兜底）——
\providecommand{\argminop}{\operatorname*{arg\,min}}
\providecommand{\Real}{\operatorname{Re}}

\rebuilt{本章重渲染自 Tutorial 笔记《分布式 MPC 与多足编队》，其源稿由 AI 撰写，论文归属、年份、卷期、引用量与前沿工作（Kim--Fawcett--Hamed、De Vincenti--Coros、Turrisi、Fawcett、Imran 等）的编号可能有误，部分实验数值（如 $15$\,ms/$6$\,ms、$<3\%$ gap、$5$\,kg 负载、变量规模表）系源稿给出、未经独立复算。本章忠实重述、不擅自“修正”，逐处 \texttt{\textbackslash cite} 标源；存疑断言已登记 claims 台账，待联网核对后二次校验。}

\begin{practice}[前置自测]
开始之前，先试答下列问题。若已能流畅作答，可只速读本章表格与速查卡；若卡住（答不出 $\ge 2$ 题），标 (1) 的回 \cref{ch:mr-consensus} 图论/Laplacian、标 (2) 的回 \cref{ch:mr-consensus} 的 ADMM 一节、标 (3) 的回单体四足 MPC（足式控制部）、标 (4) 的回凸优化对偶/KKT，再来读本章。本章会大量复用这些前置，但不重新教它们。
\begin{enumerate}
  \item （(2)）标准 ADMM 求解 $\min_{x,z}f(x)+g(z)$ s.t.\ $Ax+Bz=c$ 的三步迭代是哪三步？对偶变量 $y$ 的物理含义是什么？罚参数 $\rho$ 太大、太小分别会怎样？
  \item （(1)）通信图的 Laplacian $\mathbf{L}=\mathbf{D}-\mathbf{A}$ 的零特征值对应哪个特征向量？为什么“$\mathbf{L}$ 有且仅有一个零特征值”等价于“图连通”？
  \item （(3)）单刚体（SRB）模型把四足机器人近似成什么？状态有几维？摩擦锥约束为什么常被近似成金字塔而非直接用二阶锥？
  \item （(4)）什么叫拉格朗日对偶？对偶函数 $g(\boldsymbol{\lambda})=\min_x\mathcal{L}(x,\boldsymbol{\lambda})$ 有何性质（凹性、对偶间隙）？强对偶在什么条件下成立（Slater）？
  \item （(1)）为什么把 $N$ 台机器人的轨迹优化拼成一个联合 QP，其求解成本会随 $N$ 三次方甚至更高地膨胀？这正是本章 DMPC 要回避的“集中式爆炸”。
\end{enumerate}
\end{practice}

\section*{本章目标}
学完本章后，你应当能够：
\begin{enumerate}
  \item \textbf{论证}为什么多机编队的联合最优控制问题必须分布式求解——从集中式 MPC 的 $O(N^3)$ 计算爆炸、单点故障、通信瓶颈三个可量化角度，而非凭口号；
  \item \textbf{从零建模}编队约束：把“保持队形”写成机器人间的相对位置/相对质心等式或软约束，区分 position/displacement/distance 三范式在 MPC 代价与约束中的不同写法，并说明它们如何使联合优化在机器人间\emph{耦合}；
  \item \textbf{推导对偶分解}：把带耦合约束的联合 MPC 通过拉格朗日松弛拆成 $N$ 个本地子问题 $+$ 一个价格协调步，解释对偶变量为何叫“影子价格”、对偶子梯度法为何收敛慢、它与 ADMM 的本质区别；
  \item \textbf{从增广拉格朗日推导} consensus-ADMM 分布式 MPC：写出 $\mathbf{X}$-更新（本地 MPC $+$ 二次惩罚）、$\mathbf{z}$-更新（共识平均）、$\mathbf{y}$-更新（对偶上升）三步，分析 $\rho$ 对收敛的影响与残差平衡，并理解为何非凸接触约束下 ADMM 无收敛保证；
  \item \textbf{建立互联 SRB 网络模型}：把 $N$ 个四足 $+$ 共享负载建模为互联单刚体网络，推导刚性耦合约束产生的约束力 $\boldsymbol{\lambda}$（拉格朗日乘子），并把 CoM 协调与接触约束（摩擦锥、单边接触、踏步时序）统一写进每个 agent 的本地 MPC；
  \item \textbf{掌握分层架构}：复现“高层分布式 MPC $\to$ 中层步态调度 $\to$ 低层分布式 WBC”三层控制流，说明每层输入输出、为何 WBC 可纯本地无通信、一步通信延迟与一致性协议如何嵌入；
  \item \textbf{分析通信拓扑与稳定性}：把共识收敛理论接到 DMPC 上——说明 ADMM 协调步本质是共识，通信拓扑（$\lambda_2$、直径、时延、丢包）如何影响收敛轮数与编队稳定性，并理解 DMPC 的递归可行性与终端代价两个稳定性支柱；
  \item \textbf{做架构选型}：按 $N$、通信可靠性、对最优性的需求，在集中/分散/分布/分层之间判断该用哪一种，并说清各自代价。
\end{enumerate}

\subsection*{本章知识导航}
\textbf{一句话定位}：本章把“多个会走路的机器人保持队形协同前进”这件事，从“如果有一台超级计算机统一指挥该怎么解”（集中式 MPC），一步步拆解成“每台机器人自己解一小块、只和邻居交换预测轨迹”（分布式 MPC），并把这套方法落到多足机器人特有的难点上——\emph{质心（CoM）必须协调、每条腿的接触力必须满足摩擦锥、踏步时序还在不停切换}。

\begin{center}
\small
\begin{tikzpicture}[
  node distance=5mm and 7mm, >=Stealth,
  blk/.style={draw, rounded corners, align=center, font=\small, inner sep=3pt, fill=main!6, draw=main!60!black},
  grp/.style={draw, rounded corners, align=center, font=\small, inner sep=3pt, fill=second!6, draw=second!60!black},
  fin/.style={draw, rounded corners, align=center, font=\small, inner sep=3pt, fill=third!8, draw=third!60!black},
  ln/.style={->, thick, gray!60!black}]
\node[blk] (mot) {为何分布式\\ $O(N^3)$、单点、通信};
\node[grp, right=of mot] (form) {编队建模\\ 三范式 $\to$ 耦合};
\node[blk, right=of form] (dual) {对偶分解\\ 影子价格、慢};
\node[grp, below=9mm of mot] (admm) {consensus-ADMM\\ $\mathbf{X}/\mathbf{z}/\mathbf{y}$、$\rho$};
\node[blk, below=9mm of form, fill=main!8, draw=main!60!black] (srb) {SRB 网络\\ 耦合力 $\boldsymbol{\lambda}$、接触};
\node[grp, below=9mm of dual] (lay) {分层控制\\ MPC$\to$步态$\to$WBC};
\node[fin, below=9mm of admm] (top) {拓扑与稳定性\\ 共识接口、递归可行};
\draw[ln] (mot)--(form); \draw[ln] (form)--(dual);
\draw[ln] (mot.south)--(admm.north);
\draw[ln] (dual.south)--(lay.north);
\draw[ln] (admm.east)--(srb.west);
\draw[ln] (srb.east)--(lay.west);
\draw[ln] (admm.south)--(top.north);
\draw[ln] (lay.south) |- (top.east);
\end{tikzpicture}
\end{center}

\noindent\textbf{推荐路径}：\textbf{理论线}（理解数学）\cref{sec:dmpc-why}$\to$\cref{sec:dmpc-dual}$\to$\cref{sec:dmpc-admm}$\to$\cref{sec:dmpc-topo}，重点是对偶分解推导、ADMM 的增广拉格朗日来源、收敛性分析；\textbf{工程线}（会用就行）\cref{sec:dmpc-why}$\to$\cref{sec:dmpc-formation}$\to$\cref{sec:dmpc-admm}$\to$\cref{sec:dmpc-srb}$\to$\cref{sec:dmpc-layered}，重点是 SRB 网络建模、ADMM 骨架、分层控制流。\cref{sec:dmpc-dual,sec:dmpc-admm,sec:dmpc-srb,sec:dmpc-topo} 是四个硬核小节，值得花最多时间。

\subsection*{前置知识桥接}
本章紧接 \cref{ch:mr-consensus}（共识/ADMM/编队三范式）并依赖单体四足 MPC，这里把最关键的前置点重新激活——你不必翻回去就能跟上：
\begin{itemize}
  \item \textbf{集中式架构的 $O(N^3)$ 爆炸}（\cref{ch:mr-consensus}）。把 $N$ 台机器人的规划拼成一个联合优化，求解成本随 $N$ 三次方膨胀，且依赖一个中央节点（单点故障）。本章 \cref{sec:dmpc-why} 把它\emph{量化}到多足 MPC——一个双四足联合 QP 约 $960$ 个变量、$N=4$ 时近 $2000$ 变量，内点法每步 $O(N^3)$ 的 KKT 求解在高频控制循环里就吃不消了。
  \item \textbf{Laplacian $\mathbf{L}$ 与连通性}（\cref{ch:mr-consensus}）。$\mathbf{L}=\mathbf{D}-\mathbf{A}$、$\mathbf{L}\mathbf{1}=\mathbf{0}$、$\mathbf{L}$ 半正定、$\lambda_2>0\iff$ 图连通。本章 \cref{sec:dmpc-topo} 会让 $\mathbf{L}$ 的这些性质\emph{直接决定 ADMM 协调步的收敛速度}。
  \item \textbf{ADMM 三步与 $\rho$}（\cref{sec:cs-admm}）。从增广拉格朗日导出的 $\mathbf{x}/\mathbf{z}/\mathbf{y}$ 三步，罚参数 $\rho$ 控制“本地自主 vs 全局一致”。本章 \cref{sec:dmpc-admm} 把这套原封不动地搬到多足 MPC——唯一区别是：那里的 $\mathbf{x}$ 是标量或小向量，这里是一个机器人在整个预测时域上的状态-控制轨迹（数百维），$\mathbf{z}$ 是邻居间要一致的接口变量。算法骨架完全一样，规模和物理含义变了。
  \item \textbf{编队三范式}（\cref{sec:cs-formation}）。position/displacement/distance 三范式，且 displacement-based 控制律本质就是共识协议作用在“相对位置误差”上。本章 \cref{sec:dmpc-formation} 把这三范式\emph{升级进 MPC}——不再是静态控制律，而是预测时域上的代价函数与约束，编队误差被滚动优化。
  \item \textbf{单体四足 MPC：SRB $+$ 接触力 QP}（足式控制部）。把一台四足近似成单刚体（$12$ 维状态），在预测时域上优化质心轨迹与四条腿的接触力，约束含摩擦锥、单边接触、力矩限制。本章 \cref{sec:dmpc-srb} 把 $N$ 台 SRB \emph{用刚性耦合约束连成网络}——单体 MPC 是本章的“原子”，本章把 $N$ 个原子协调起来。
\end{itemize}

\noindent\textbf{前向预告}：本章的编队保持的是“机器人之间的相对位姿/质心”；当机器人要\emph{共同搬运一个负载}时，会多出一层——它们对负载施加的力必须协调，否则会把负载“撕裂”或内耗。后续协同搬运章会引入 grasp matrix 与内力分配，把本章的“运动协调”扩展到“力协调”。现在只需记住：\textbf{本章管“一起走到哪”，下一章管“一起使多大劲”}——本章 \cref{sec:dmpc-srb} 的刚性耦合力 $\boldsymbol{\lambda}$ 就是通往内力的桥梁。

\subsection*{如果跳过本章会怎样}
\begin{itemize}
  \item \textbf{集中式 MPC 实时性崩了}：两台 Go2 保持 $1$ 米间距同步行走，图省事把状态和控制堆成一个联合 MPC 扔给 OSQP。$N=2$、预测 $20$ 步勉强能跑；加到 $N=4$，求解时间飙升、控制频率掉到危险区、机器人摇晃失稳。你知道“集中式不可扩展”，却不知道\emph{具体在哪一步爆炸}（内点法 KKT 求解的 $O(N^3)$），更不知道怎样把联合 QP 拆开。\cref{sec:dmpc-why,sec:dmpc-admm} 正是讲这个拆解。
  \item \textbf{ADMM 写出来了却几百轮不收敛}：照着骨架把联合 MPC 用 ADMM 切分，$\rho$ 随手取 $1.0$，结果要么剧烈震荡、要么 $500$ 轮才勉强一致。你在代码里到处找 bug——但这不是 bug。\cref{sec:dmpc-admm} 会告诉你：$\rho$ 的量级直接决定收敛速度，且对非凸接触约束 ADMM 根本没有收敛保证，工程上靠 warm-start $+$ 固定迭代轮数 $+$ 残差平衡来“够用就停”。
  \item \textbf{耦合力凭空多出来不知怎么处理}：给“双四足抬一根刚性杆”建 MPC，照搬单体 SRB，却发现两台机器人通过杆\emph{互相施力}——这力既不是接触力也不是重力。\cref{sec:dmpc-srb} 会告诉你：刚性连接产生\emph{约束力 $\boldsymbol{\lambda}$（拉格朗日乘子）}，由“附着点相对位置固定”这个约束微分两次后解出，必须显式进入每个 agent 的 SRB 动力学，并在 ADMM 协调步达成一致。
\end{itemize}

\begin{table}[H]\centering\small
\caption{本章阅读方式建议}
\label{tab:dmpc-reading}
\begin{tabularx}{\linewidth}{LlL}
\toprule
阅读方式 & 时间 & 适合谁 \\
\midrule
精读（推完对偶分解与 ADMM 增广拉格朗日、推导 SRB 耦合力、复现分层控制流） & $22$--$30$ 小时 & 第一次系统学分布式 MPC 与多足编队 \\
速读（读懂集中$\to$分布动机、ADMM 三步在 MPC 里的形式、SRB 网络与分层架构，跳过推导细节） & $6$--$8$ 小时 & 有 ADMM 与单体 MPC 背景、想建立全局图景 \\
速查（集中 vs 分布对比表 $+$ ADMM-MPC 三步 $+$ SRB 网络方程 $+$ 分层数据流 $+$ 故障排查） & $50$--$70$ 分钟 & 实现时回查 \\
\bottomrule
\end{tabularx}
\end{table}

\begin{note}
\textbf{本章记号与约定.}\;
旋转矩阵记 $\mathbf{R}\in\SO(3)$、位姿 $\mathbf{T}\in\SE(3)$（preamble 宏 $\SO,\SE,\so,\se$，与全书一致，\nameref{ch:notation}）；MPC/优化/图论领域的自然记号保留。
$N$ 为 agent 数、$T$ 为预测时域步数、$i,j$ 为 agent 索引（$j\in\mathcal{N}_i$ 表 $j$ 是 $i$ 的通信邻居）。agent $i$ 的状态 $\mathbf{x}_i$（SRB 为 $12$ 维 $[\mathbf{p}_i,\boldsymbol{\Theta}_i,\dot{\mathbf{p}}_i,\boldsymbol{\omega}_i]$）、控制 $\mathbf{u}_i$（足端接触力堆叠）；整个时域上的轨迹记 $\mathbf{X}_i,\mathbf{U}_i$，本地代价 $J_i(\cdot)$。质心位置 $\mathbf{p}_i\in\mathbb{R}^3$，期望相对位移 $\mathbf{d}_{ij}$，摩擦系数 $\mu$，附着点 $\mathbf{p}_i^{\text{att}}$。
\textbf{字母隔离声明}：本章 $\boldsymbol{\lambda}$ 在 \cref{sec:dmpc-dual} 指对偶变量、在 \cref{sec:dmpc-srb} 指刚性耦合约束力（二者数学上同为拉格朗日乘子，物理对应见正文）；$\rho$ 指 ADMM 罚参数；$\mathbf{z}$ 指 ADMM 全局共识变量、$\mathbf{y}_i$ 指对偶变量、$\mathbf{u}_i^{\text{s}}$（或正文写的缩放 $\mathbf{u}_i$）指缩放对偶变量；$r^k,s^k$ 指第 $k$ 轮原始/对偶残差。上标星 $(\cdot)^*$ 表最优。
\end{note}

\subsection*{科研发展脉络}
在钻进具体方法前，先把这条研究线的来龙去脉理清——知道每个方法“从哪来、解决了前人什么痛点、又留下什么给后人”，比孤立地记名字有用得多。\rebuilt{下表的年份、venue、引用量与作者列表系 Tutorial 源稿给出，多处存疑（见各行 \texttt{\textbackslash pz}），引用量尤其粗略，下游须联网核验。}

\begin{table}[H]\centering\small
\caption{分布式 MPC 与多足协同的研究脉络}
\label{tab:dmpc-lineage}
\begin{tabular}{p{0.06\linewidth}p{0.40\linewidth}p{0.14\linewidth}p{0.30\linewidth}}
\toprule
年份 & 论文 & Venue & 核心贡献 \\
\midrule
2006 & Dunbar \& Murray, Distributed Receding Horizon Control for Multi-Vehicle Formation Stabilization\cite{dunbar2006distributed} & Automatica & \textbf{奠基}：每车解本地 RHC $+$ 交换预测轨迹；“假设轨迹” $+$ 一致性惩罚保稳定 \\
2008 & Keviczky, Borrelli 等, Decentralized Receding Horizon Control\ldots\cite{keviczky2008decentralized} & Automatica/T-CST & DMPC 编队系统化：预测交换 $+$ 迭代协商；动力学解耦、约束耦合的标准范式\pz{venue 应为 IEEE T-CST 16(1):19--33, 2008（非 Automatica），作者全表为 Keviczky, Borrelli, Fregene, Godbole, Balas。另：标题含“dynamically decoupled”的是同作者另一篇 Keviczky--Borrelli--Balas, Automatica 42(12):2105--2115, 2006，勿与本键（T-CST 2008）混用——见 \cref{sec:dmpc-spectrum} 陷阱处的同一标注。} \\
2010s & Summers \& Lygeros; Conte 等\cite{conte2016distributed} & CDC/ECC/Automatica & \textbf{ADMM/对偶分解求解 DMPC}（凸假设）；合作式 DMPC 稳定性综合 \\
2022 & Kim, Hamed, Cooperative Locomotion via Supervisory Predictive Control\cite{kim2022cooperative} & ACC\rebuilt{venue 错——该文实为 ASME J.\ Dyn.\ Sys.\ Meas.\ Control 144(3):031005, 2022（期刊，非 ACC 会议）；DOI 10.1115/1.4052917。论文/作者/年份/实质均正确，仅 venue 误标“ACC”，待修正、暂不就地改。} & 监督式预测控制 $+$ 分布式非线性控制器，多足协同早期模板 \\
2023 & \textbf{Kim, Fawcett, Kamidi, Ames, Hamed, Layered Control for Cooperative Locomotion of Two Quadrupedal Robots}\cite{kim2023layered} & \textbf{T-RO} & \textbf{多足协同标杆}：互联 SRB 网络 $+$ 集中/分布 MPC（一步通信延迟 $+$ agreement protocol）$+$ 低层分布式控制器 \\
2023 & De Vincenti, Coros, Centralized MPC for Collaborative Loco-Manipulation\cite{devincenti2023centralized} & RSS & 集中式 SQP：robot/payload/gait-agnostic；$\SO(3)$ 李群四元数积分 \\
2024 & Turrisi 等, PACC: Passive-Arm Approach for High-Payload Collaborative Carrying\cite{turrisi2024pacc} & IROS & 被动臂方案——每个四足加被动臂托住负载，\emph{免去刚性耦合建模} \\
2024 & Fawcett, Amanzadeh, Hamed 等, Distributed Data-Driven Predictive Control\ldots\cite{fawcett2024datadriven} & — & 数据驱动模板模型 $+$ 分布式预测控制，绕开精确动力学建模\pz{年份/venue 误：该文实为 Proc.\ IEEE ICRA 2023, pp.\,9924--9930（arXiv:2211.06917，2022-11 投稿；获 ICRA 2023 Outstanding Paper），非 2024、亦非“—”；键宜改 fawcett2023distributed。又：“数据驱动模板模型（data-driven template model）”一语出自姊妹论文 Fawcett 等 RA-L 7(3):7636--7643, 2022，本文用的是 behavioral/DeePC 式数据驱动预测控制，措辞被借用，需区分。} \\
2025/26 & Imran, Hamed 等, Obstacle-Aware Distributed NMPC\cite{imran2025obstacle} / ADMM-Based DMPC with CBF\cite{zeng2026admmcbf} & J.\ Dyn.\ Sys./arXiv & 障碍环境分布式 NMPC；ADMM $+$ 控制屏障函数（CBF）的安全多足分布式 MPC \\
\bottomrule
\end{tabular}
\end{table}

\pz{实验室“脉络”与引用量经核多处不实，作为个人/团队谱系是张冠李戴。Kaveh Akbari Hamed 的实际履历（VT 官方 CV）为：Univ.\ of Tabriz $\to$ Sharif Univ.\ of Technology（博 2011）$\to$ 博后 Univ.\ of Michigan（2012--2014）$\to$ San Diego State Univ.（2014--2018）$\to$ Virginia Tech（2018--），与 Caltech/UPenn/UC~Berkeley/ETH 均无师承；与 Caltech (Ames) 仅为 NSF 课题合作，非谱系。引用量方面：ARCOL/Hamed 无任何论文达 $\sim$$1500$ 或 $\sim$$500$（其个人总被引约 $1.7$k、单篇最高约 $262$）——$1500$/$500$ 量级属上游他人经典（如 ETH/Hutter、Ames CBF）；$30$/$50$/$\sim$$80$ 可分别对应 ARCOL 真实论文。故“机构谱系收束于 ARCOL”系虚构、引用量不可按此归属，整段仅作粗略背景、勿引以为据。}

看这条线，有一条清晰的主线：\textbf{决策从“一个中央大脑解全局”走向“每个 agent 解本地 $+$ 邻居协商”}，而协商工具从早期的“交换预测轨迹 $+$ 一致性惩罚”（Dunbar/Keviczky），演进到“对偶分解/ADMM 严格分解联合优化”（Summers/Conte），再落地到“多足特有的 SRB 网络 $+$ 分层 MPC/WBC”（Hamed 团队）。本章以 Kim--Fawcett--Hamed T-RO 2023\cite{kim2023layered} 为核心案例，因为它最完整地把“分布式 MPC $+$ 多足 $+$ 编队/协同”三者打通。

\begin{insight}[分布式 MPC 史 $=$ 一部“安全地切开紧耦合全局优化”的历史]\label{ins:dmpc-history}
分布式 MPC 的研究史，本质是一部“如何把一个紧耦合的全局优化，安全地切成可并行的局部优化”的历史。早期方法（交换预测轨迹）切得“软”——靠惩罚项鼓励一致，不保证严格满足耦合约束；ADMM 切得“硬”——通过对偶变量强制收敛到满足耦合约束的解。多足把难度又抬高一层：耦合约束不再是线性的相对位置，而是带摩擦锥、单边接触的非凸约束，于是连 ADMM 的收敛保证都失效了——这正是 $2024$–$2026$ 年这个方向仍在活跃发论文的原因。
\end{insight}

% ════════════════════════════════════════════════════════════════
\section{从集中式到分布式：为什么联合 MPC 必须拆开}
\label{sec:dmpc-why}

\paragraph{动机：两台机器人好解，四台就崩。}
先把场景定下来。$N$ 台四足机器人要保持一个队形（比如一字排开、间距 $1$ 米）向前行走。每台用单刚体（SRB）模型描述，状态 $12$ 维。我们希望在每个控制周期（典型 $1$\,kHz，即每 $1$\,ms 一次）求解一个 MPC：在未来 $T$ 步（比如 $T=20$，对应 $0.4$ 秒预测）内，优化每台机器人的质心轨迹与四条腿的接触力，使得（a）每台机器人自身平衡、踏步可行；（b）机器人之间保持期望队形。最直接的想法是\textbf{集中式}：把所有机器人的状态和控制堆成一个大向量，写一个联合优化问题，交给一台中央计算机求解，再把结果分发回各机器人：
\begin{align}
  \min_{\mathbf{U}_1,\dots,\mathbf{U}_N}\ & \sum_{i=1}^N J_i(\mathbf{X}_i,\mathbf{U}_i) + J_{\text{form}}(\mathbf{X}_1,\dots,\mathbf{X}_N) \label{eq:dmpc-central-obj}\\
  \text{s.t.}\quad & \mathbf{x}_i(t{+}1)=f_i(\mathbf{x}_i(t),\mathbf{u}_i(t)),\quad \forall i,t \quad(\text{各自动力学}),\notag\\
  & \text{formation}(\mathbf{X}_1,\dots,\mathbf{X}_N)=\mathbf{0}\quad(\text{编队耦合}),\notag\\
  & \text{friction\_cone}(\mathbf{u}_i)\le\mathbf{0},\quad \forall i\quad(\text{接触约束}).\notag
\end{align}
这个问题在 $N=2$ 时确实能解。问题是：\textbf{它解不动 $N=4$，更解不动 $N=10$。} 为什么？这正是本节要拆开讲的，也是 \cref{sec:dmpc-dual,sec:dmpc-admm} 所有分解方法的动机地基。

\paragraph{反面：如果坚持用集中式会怎样。}
集中式 MPC 的失败不是单一原因，而是三个瓶颈叠加。

\textbf{瓶颈一：计算爆炸（$O(N^3)$）。} 联合 MPC 的决策变量是所有机器人在整个时域上的控制（与被消去前的状态），变量数 $\approx N\cdot(n_x+n_u)\cdot T$。对 SRB（$n_x=12$，每条腿 $3$ 维力、$4$ 条腿 $n_u=12$），$T=20$，得到 \cref{tab:dmpc-blowup}。

\begin{table}[H]\centering\small
\caption{集中式联合 MPC 随 $N$ 的计算膨胀（源稿数值，未独立复算）}
\label{tab:dmpc-blowup}
\begin{tabular}{ccccc}
\toprule
$N$ & 联合变量数（量级） & 内点法每步 KKT $O(\text{vars}^3)$ 相对成本 & $1$\,kHz 下能否实时 \\
\midrule
$1$（单体） & $\sim480$ & $1\times$ & 轻松 \\
$2$ & $\sim960$ & $\sim8\times$ & 勉强（$\sim15$\,ms） \\
$4$ & $\sim1920$ & $\sim64\times$ & 掉到 $\sim16$\,Hz \\
$10$ & $\sim4800$ & $\sim1000\times$ & 完全不可能 \\
\bottomrule
\end{tabular}
\end{table}

\noindent 内点法每步要解一个 KKT 线性系统，稠密求解成本为变量数的立方 $O(\text{vars}^3)$。$N$ 从 $2$ 到 $4$，变量翻倍、KKT 求解成本翻约 $8$ 倍（$2^3$），控制频率从勉强可用直接掉到危险区。MPC 的价值在于高频滚动重规划，频率一掉，预测就跟不上扰动，机器人失稳。\pz{存疑（数值有据但结论夸大、与源不符）：算术骨架（每机器人 $480$ 变量、$480N$、倍率 $N^3$）自洽可复算；但 (i) $N^3$ 仅是\emph{稠密} QP 最坏情形，真实集中式 MPC 用稀疏/结构化求解器（OSQP/HPIPM/Riccati/condensing）成本约随 stage 数\emph{线性}，不应把 $N^3$ 当固有缩放；(ii) “$N=10\approx4800$ 变量＝不可能”被本表所依据的源论文 De Vincenti--Coros (RSS 2023) 自身反驳——其脚注载 Grandia 等以 $\sim$$5000$ 决策变量、$3$ 机器人达 $100$\,Hz；(iii) 时序锚点互不自洽：“$N{=}2\to15$\,ms”（$\times8$）隐含基线 $\sim$$1.9$\,ms，“$N{=}4\to16$\,Hz”（$\times64$）隐含基线 $\sim$$1.0$\,ms，两者矛盾，$15$\,ms 系拍脑袋；(iv) 维度与所引 CLM 论文不符（该文 $n_x{=}13$ 含四元数、$T{=}15$、每机器人输入 $\sim$$30$，非本表 $n_x{=}12,T{=}20,n_u{=}12$——后者对“单台四足 SRB”标准正确，但用于该集中式 CLM 情景与原文不符）。本表仅作“朴素稠密集中式”示意，具体时延/可行性须以稀疏求解器实测并核对原论文维度/频率为准。}

\textbf{瓶颈二：通信瓶颈与单点故障。} 集中式需要一个中央节点：所有机器人每周期把状态发给它、算完再把控制发回。于是（a）中央节点上下行带宽 $\propto N$，机器人一多就堵；（b）中央节点宕机或与某台机器人断联，\emph{整个编队瘫痪}——这是单点故障。野外多机器人作业（搜救、巡检）中，可靠的中央通信往往根本不存在。

\textbf{瓶颈三：隐私与模块化。} 集中式要求中央节点知道每台机器人的完整模型（动力学参数、约束）。异构编队里不同机型模型可能由不同厂商提供、不愿暴露；每加一台机器人就要重写整个联合问题，毫无模块化可言。

\begin{insight}[三个瓶颈是同一根源的三个侧面]\label{ins:dmpc-three-bottleneck}
集中式 MPC 的三个瓶颈不是独立的工程缺陷，而是同一个根源的三个侧面——\textbf{它把“协调”建模成了“一个全知全能的中心”}。计算爆炸是因为中心要解全局；单点故障是因为中心不可替代；不可模块化是因为中心要懂全部。分布式 MPC 的全部努力，就是\textbf{取消这个中心}，让协调从“中心统一裁决”变成“邻居之间协商”。
\end{insight}

\paragraph{反面对照：那能不能干脆“各管各的”（完全去中心、零协调）？}
走到另一个极端：每台机器人完全独立解自己的单体 MPC，\emph{互不通信}。这当然没有计算爆炸、没有单点故障。但它\textbf{根本保持不了队形}——每台机器人只优化“我自己走得好不好”，没有任何项告诉它“我离邻居该多远”。机器人之间的相对位置会随各自扰动、各自局部最优而漂移，队形分崩离析。更糟的是，如果它们共享一个负载，零协调会让它们对负载施加\emph{互相冲突的力}，把负载内耗甚至撕裂。

\begin{insight}[不是把问题切成 $N$ 块各自解，而是切成 $N$ 块协同收敛]\label{ins:dmpc-not-split}
分布式 MPC 不是“把集中式问题简单地切成 $N$ 块各自解”（那是完全去中心，会丢掉协调），\textbf{而是“切成 $N$ 块、但通过邻居间的信息交换让这些块协同收敛到（接近）联合最优解”}。关键词是“协同收敛”——既要并行（解决计算爆炸），又要一致（保住编队）。这两个目标的张力，正是 \cref{sec:dmpc-dual,sec:dmpc-admm} 所有分解算法要平衡的。
\end{insight}

\subsection{分布式 MPC 的基本范式：预测交换}
\label{sec:dmpc-prediction-exchange}

集中式（全知中心）和完全去中心（零协调）是两个极端。分布式 MPC 落在中间：\textbf{每个 agent 解一个本地 MPC，但在求解时，把邻居的“预测轨迹”作为已知信息纳入考虑。} 这个思想最早由 Dunbar \& Murray（$2006$）\cite{dunbar2006distributed} 和 Keviczky--Borrelli（$2008$）\cite{keviczky2008decentralized} 系统化，核心机制叫\textbf{预测交换（prediction exchange）}或\textbf{假设轨迹（assumed trajectory）}：
\begin{enumerate}
  \item \textbf{每个 agent 维护并广播自己的预测轨迹}。agent $i$ 解完本地 MPC 后，得到自己未来 $T$ 步的状态预测 $\hat{\mathbf{X}}_i=[\hat{\mathbf{x}}_i(1),\dots,\hat{\mathbf{x}}_i(T)]$，把它发给所有邻居 $j\in\mathcal{N}_i$。
  \item \textbf{每个 agent 用邻居的预测构造本地的编队项}。agent $i$ 解本地 MPC 时，编队代价里邻居 $j$ 的状态用\emph{上一轮收到的} $\hat{\mathbf{X}}_j$ 代入（当作已知常量），于是耦合的 $J_{\text{form}}(\mathbf{X}_1,\dots,\mathbf{X}_N)$ 退化成只含 $\mathbf{X}_i$ 的本地项。
  \item \textbf{迭代/滚动}。各 agent 用新的预测再算一轮（同步协商），或直接进入下一个控制周期（异步、靠滚动收敛）。
\end{enumerate}
形式上，agent $i$ 的本地问题变成
\begin{align}
  \min_{\mathbf{U}_i}\ & J_i(\mathbf{X}_i,\mathbf{U}_i) + \sum_{j\in\mathcal{N}_i} J_{ij}(\mathbf{X}_i, \hat{\mathbf{X}}_j) \label{eq:dmpc-pred-exchange}\\
  \text{s.t.}\quad & \mathbf{x}_i(t{+}1)=f_i(\mathbf{x}_i(t),\mathbf{u}_i(t)),\quad \text{friction\_cone}(\mathbf{u}_i)\le\mathbf{0}.\notag
\end{align}
现在每个 agent 解的问题规模是\textbf{单体大小}（变量数 $\approx(n_x+n_u)T$，与 $N$ 无关），$N$ 个 agent 并行求解——计算爆炸解决了。

\begin{insight}[预测交换像列队前进时的余光（标边界）]\label{ins:dmpc-peripheral}
预测交换很像\textbf{一群人排队列队前进时各自的余光}。每个人主要盯着自己脚下走稳（本地 MPC），同时用余光瞥一眼左右邻居的位置和趋势（邻居的预测轨迹），据此微调步伐保持队形。\emph{像}的地方：没有教官统一指挥（无中心），每个人只看邻居（局部信息），整体队形却能维持。\emph{不像}的地方：人的余光是连续实时的，而预测交换是\textbf{离散轮次}的——agent $i$ 用的是邻居\emph{上一轮}算出的预测（有一步信息延迟），这个延迟正是 \cref{sec:dmpc-topo} 要分析的稳定性隐患；且人凭直觉调整，预测交换是在解一个严格的优化问题。
\end{insight}

但预测交换有个根本弱点：\textbf{它对编队约束的满足是“软”的}。编队项 $J_{ij}$ 是惩罚项（鼓励 $\mathbf{X}_i$ 靠近期望相对位置），不是硬约束。如果两个 agent 各自基于对方的旧预测优化，它们可能\emph{永远差一点点}而不完全一致——因为没有任何机制强制它们的预测互相吻合。这在“必须严格满足耦合约束”的场景（如刚性连接的负载，附着点位置必须精确一致）下不够用。这正是引出\textbf{对偶分解}（\cref{sec:dmpc-dual}）和 \textbf{ADMM}（\cref{sec:dmpc-admm}）的动机：它们用拉格朗日乘子/对偶变量，把“软鼓励”升级为“强制收敛到满足耦合约束的解”。Dunbar/Keviczky 的预测交换是分布式 MPC 的“第一代”，对偶分解/ADMM 是“第二代”——后者在前者的并行结构上，补上了\emph{严格满足耦合约束的保证}（凸情形下）。

\subsection{系统性分类：MPC 控制架构的完整谱系}
\label{sec:dmpc-spectrum}

为避免把“集中 vs 分布”看成非黑即白的二分，这里给一个\textbf{系统性的架构谱系}。按“谁掌握全局信息、谁做决策、需要多少通信”分类，多机 MPC 架构落在一条连续谱上（\cref{tab:dmpc-spectrum}）。

\begin{table}[H]\centering\small
\caption{多机 MPC 控制架构谱系}
\label{tab:dmpc-spectrum}
\begin{tabular}{p{0.12\linewidth}p{0.11\linewidth}p{0.10\linewidth}p{0.12\linewidth}p{0.11\linewidth}p{0.10\linewidth}p{0.13\linewidth}}
\toprule
架构 & 全局信息 & 决策主体 & 通信 & 最优性 & 鲁棒性 & 典型代表 \\
\midrule
集中式 & 中央全知 & 中央节点 & 全部状态上下行 & 全局最优 & 单点故障 & De Vincenti RSS 2023\cite{devincenti2023centralized} \\
分散式 & 各自本地 & 各 agent & 零通信 & 无协调（队形崩） & 无单点 & 完全去中心 \\
分布式 & 本地 $+$ 邻居 & 各 agent 协商 & 仅邻居、接口变量 & 近最优（有 gap） & 无单点 & Kim--Hamed T-RO 2023\cite{kim2023layered} \\
分层 & 分层 & 分层 & 高层通信、低层本地 & 近最优 & 高层弱中心 & 监督式预测控制\cite{kim2022cooperative} \\
\bottomrule
\end{tabular}
\end{table}

\textbf{关键区分（消除常见混淆）}：
\begin{itemize}
  \item \textbf{分散式（decentralized）$\neq$ 分布式（distributed）}。这是最常被混淆的一对。分散式 $=$ \emph{零通信}、各自为政（队形崩）；分布式 $=$ \emph{邻居通信}、协商一致（本章主角，队形保持）。中文里两词都常译“分布式”，但英文严格区分——读论文务必看清是 decentralized（无协调）还是 distributed（有协调）。
  \item \textbf{分布式 vs 分层不是对立的}。Kim--Fawcett--Hamed 同时是分布式（高层 MPC 邻居协商）和分层（高层 MPC $+$ 低层 WBC）——分层是“垂直”切分（按频率/抽象层），分布式是“水平”切分（按 agent）。一个真实系统两个维度都切（\cref{sec:dmpc-layered}）。
\end{itemize}

\begin{insight}[四种架构是“协调强度”轴上的采样点]\label{ins:dmpc-coordination-axis}
这四种架构不是四个孤立的选项，而是\textbf{“协调强度”这条轴上的采样点}——从“中央全知全控”（集中式，协调最强但最脆）到“各自为政”（分散式，最鲁棒但无协调），分布式和分层落在中间，用“邻居通信”换取“协调 $+$ 鲁棒”的平衡。选架构的本质是回答：\emph{我需要多强的协调？我能容忍多大的通信/计算成本？我的中央节点有多可靠？} 多足协同搬运需要强协调（耦合约束严格）$+$ 高鲁棒（野外无可靠中央）$+$ 实时（毫秒级），这三个需求的交集恰好把它推向“分布式 $+$ 分层”——这就是本章主线方法是这个组合、而非任意选的原因。拿到一篇多机 MPC 论文，先把它放进谱系表的某一格（看通信模式和决策主体），你就立刻知道它的最优性/鲁棒性/计算的大致权衡，不必读完全文。
\end{insight}

\begin{pitfall}[概念误区：以为分布式 MPC 一定比集中式“更好”]
认为既然分布式解决了计算爆炸，就应无脑用分布式替代集中式。\emph{后果}：在 $N=2$、有可靠中央计算机的场景下用分布式，反而引入 ADMM 迭代的额外延迟和“不收敛风险”，轨迹质量还比集中式差（有最优性 gap）。\emph{根因}：分布式是\textbf{用最优性换可扩展性和鲁棒性}的权衡，不是免费午餐。\emph{正确做法}：$N$ 小（$\le2$–$3$）、中央通信可靠 $\to$ 集中式；$N$ 大、需鲁棒性/模块化 $\to$ 分布式。\rebuilt{原句所附数值（集中约 $15$\,ms、分布约 $6$\,ms、$<3\%$ gap）经核与 \cite{kim2023layered} 源文矛盾、系杜撰：源文集中式 $1.38$\,ms、分布式 $0.41$\,ms（降约 $70\%$），且无“$<3\%$ gap”（仅“性能相近”）。已从正文删去这些数，原句见下方注释行待修。}
% —— 原源稿句（数值杜撰、停用待修）：正确做法：$N$ 小（≤2–3）、中央通信可靠 → 集中式（源稿称 Kim--Hamed 集中式 MPC 计算约 15 ms、轨迹最优）；$N$ 大、需鲁棒性/模块化 → 分布式（源稿称各 agent 并行约 6 ms、轨迹 <3% gap）。——
\end{pitfall}

\begin{pitfall}[思维陷阱：把“动力学解耦”和“约束解耦”混为一谈]
以为既然每台机器人有自己的动力学方程（动力学解耦），整个问题就天然可分。\emph{后果}：写本地 MPC 时漏掉了编队耦合的处理，结果各 agent 各走各的，队形崩溃（退化成完全去中心）。\emph{根因}：多足编队是\textbf{动力学解耦但约束（和代价）耦合}的系统——每台机器人的动力学 $f_i$ 确实只依赖自己的 $\mathbf{x}_i,\mathbf{u}_i$，但编队约束/代价 $J_{\text{form}}$ 把不同 agent 的状态绑在一起。Keviczky--Borrelli 2008\cite{keviczky2008decentralized} 标题里的“dynamically decoupled”说的正是前者，但耦合恰恰来自后者。\pz{标注误植：标题含“dynamically decoupled”的是同作者另一篇——Keviczky, Borrelli, Balas, “Decentralized receding horizon control for large scale dynamically decoupled systems”, Automatica 42(12):2105--2115, 2006；本键 keviczky2008decentralized 指的 T-CST 2008 论文标题并无该短语。此处概念正确（动力学解耦、约束/代价耦合），但勿把该短语归于本键所引论文。}\emph{正确做法}：动力学解耦让“本地求解”成为可能；约束/代价耦合是“必须协调”的根源（\cref{sec:dmpc-dual,sec:dmpc-admm} 的分解就是处理这一层）。
\end{pitfall}

\begin{pitfall}[概念误区：以为预测交换里“用邻居的预测”是用当前轮的实时值]
实现时假设 agent $i$ 能拿到邻居 $j$ \emph{本轮}正在算的最新轨迹。\emph{后果}：在异步/有通信延迟的真实系统里这个假设不成立——你拿到的总是邻居上一轮的结果；若代码逻辑依赖“实时同步”，会在真机上出现难以复现的震荡。\emph{根因}：分布式系统中信息天然有延迟，Kim--Fawcett--Hamed\cite{kim2023layered} 明确建模为“one-step communication delay（一步通信延迟）”。\emph{正确做法}：设计算法时显式假设“邻居信息滞后一步”，warm-start 用滞后的邻居预测，不要假装能拿到实时值。
\end{pitfall}

\begin{exercise}
\begin{enumerate}
  \item \textbf{（计算瓶颈量化）} 设内点法求解 QP 成本 $C(\text{vars})=\alpha\cdot\text{vars}^3$，单体 MPC（$480$ 变量）实测 $10$\,ms。估算 $N=2,4,6,8$ 集中式联合 MPC 求解时间（只看变量立方）。在 $1$\,kHz（预算 $1$\,ms）和 $200$\,Hz（预算 $5$\,ms）下各自最多支持几台机器人集中式实时？
  \item \textbf{（预测交换 vs 完全去中心）} 用 Python 写最小例子：$N=3$ 个双积分器（$\ddot{\mathbf{p}}_i=\mathbf{u}_i$，2D）保持一字队形（相邻间距 $1$\,m）向 $+x$ 移动。分别实现（a）完全去中心、（b）预测交换（加“与邻居预测保持期望相对位置”的惩罚，用邻居上一轮预测）。对比 $30$ 个控制周期后的队形误差。预期：(a) 漂移，(b) 大致维持但有残差。
  \item \textbf{（开放设计）} 阅读 Dunbar \& Murray 2006\cite{dunbar2006distributed} 与 Keviczky--Borrelli 2008\cite{keviczky2008decentralized} 摘要，用一段话说明：为什么“交换预测轨迹”还不足以保证编队约束被严格满足？描述两个 agent 都用对方旧预测、编队项只是惩罚而非硬约束时“永远差一点点”的具体情形。
\end{enumerate}
\end{exercise}

% ════════════════════════════════════════════════════════════════
\section{编队约束建模：把“保持队形”写进 MPC}
\label{sec:dmpc-formation}

\paragraph{动机：队形不是一个数，是一组关系。}
“保持一字队形，间距 $1$ 米”——这句话怎么变成优化问题里的数学对象？初学者的第一反应往往是：给每台机器人指定一个目标位置（agent $1$ 去 $(0,0)$、agent $2$ 去 $(1,0)$、agent $3$ 去 $(2,0)$），让每台机器人 MPC 跟踪自己的目标点。但这立刻暴露一个问题：\textbf{如果整个编队要往前走、转弯、绕障，这些“绝对目标点”每个周期都得有人重新算并下发}——这又退回到了需要中央节点的集中式。更本质的观察是：\textbf{队形描述的是机器人之间的“相对关系”，而不是每台机器人的“绝对位置”。} “间距 $1$ 米”约束的是 agent $i$ 和 agent $j$ 的相对位移 $\mathbf{p}_j-\mathbf{p}_i$，而整个队形可以自由平移、旋转。这决定了：编队约束应写成\textbf{相对量}，这样每台机器人只需和邻居比较、无需知道全局坐标——这正是分布式的前提。\cref{sec:cs-formation} 把 displacement-based 控制律刻画为共识协议作用在“相对位置误差”上；本节把那个\emph{静态控制律}升级——不再是“当前误差 $\to$ 当前控制”，而是“在未来 $T$ 步预测时域上让相对误差的累积代价最小”，让编队对扰动和障碍有了预测性的应对。

\paragraph{反面：如果用绝对位置建模会怎样。}
假设硬要用绝对目标点 $\mathbf{p}_i^{\text{ref}}$，编队代价写成 $\sum_i\|\mathbf{p}_i-\mathbf{p}_i^{\text{ref}}\|^2$，会出三个问题：
\begin{enumerate}
  \item \textbf{需要全局坐标系与中央规划器}。谁来决定 $\mathbf{p}_i^{\text{ref}}$？编队前进时这些目标点持续变化，必须有个中央节点统一计算队形整体运动并下发——单点故障回来了。
  \item \textbf{过约束}。绝对位置把队形的平移/旋转自由度也锁死了。但很多任务只关心“保持相对队形”，整体怎么移动应由更高层（或领航者）灵活决定。锁死全局位姿会让编队无法自然避障或跟随。
  \item \textbf{不可分布式}。$\|\mathbf{p}_i-\mathbf{p}_i^{\text{ref}}\|^2$ 看似可分（每项只含一个 agent），但 $\mathbf{p}_i^{\text{ref}}$ 的生成依赖全局信息，本质仍是中心化的。
\end{enumerate}

\begin{insight}[约束相对几何，让整体位姿成为自由自由度]\label{ins:dmpc-relative}
编队建模的核心不是“给每个机器人一个绝对目标”，\textbf{而是“约束机器人之间的相对几何，让队形的整体位姿成为自由（或由领航者决定）的自由度”}。前者把队形钉死在地图上，后者让队形像一个可平移可旋转的“刚性图案”漂浮在空间里。
\end{insight}

\paragraph{历史：从图刚性到编队三范式。}
编队控制的数学源头之一是\textbf{图刚性理论（graph rigidity）}。把每台机器人看作图的顶点、机器人间的“相对位置约束”看作边，“队形可保持”等价于“约束图是刚性的”——满足这些边约束的构型整体只剩平移和旋转自由度。Tanner、Pappas、Kumar 在 $2000$ 年代初把图刚性引入多机编队，奠定了“用相对量描述队形”的范式。\pz{归属误：Tanner--Pappas--Kumar 2000s 的工作是 formation graphs 上的 input-to-state / leader-to-formation 稳定性（CDC 2002、IEEE T-RA 2004），\emph{非}图刚性。把图刚性引入分布式编队镇定的奠基是 Olfati-Saber \& Murray, “Graph rigidity and distributed formation stabilization of multi-vehicle systems”, CDC 2002；后续 Eren/Anderson/Whiteley 2004、Krick/Broucke/Francis 2009，数学根在 Laman、Asimow--Roth、Tay--Whiteley。本句应把 Tanner/Pappas/Kumar 换成 Olfati-Saber \& Murray 2002 等。本行主引 Oh--Park--Ahn 2015 综述本身正确。} 后来，编队控制按“用什么相对量约束、需要什么传感/通信”分成三个层次清晰的范式（Oh--Park--Ahn 2015 综述\cite{oh2015formation} 系统化），见 \cref{tab:dmpc-three-paradigm}。

\begin{table}[H]\centering\small
\caption{编队三范式（依 Oh--Park--Ahn 2015\cite{oh2015formation}）}
\label{tab:dmpc-three-paradigm}
\begin{tabular}{p{0.14\linewidth}p{0.16\linewidth}p{0.18\linewidth}p{0.18\linewidth}p{0.20\linewidth}}
\toprule
范式 & 约束量 & 需要的信息 & 队形保留的自由度 & 通信/传感要求 \\
\midrule
Position-based & 绝对位置 $\mathbf{p}_i$ & 全局坐标 & 无（完全锁定） & 全局定位（GPS/动捕） \\
Displacement-based & 相对位移 $\mathbf{p}_j-\mathbf{p}_i$ & 共同朝向 $+$ 相对位置 & 平移 & 邻居相对位置 $+$ 朝向一致 \\
Distance-based & 相对距离 $\|\mathbf{p}_j-\mathbf{p}_i\|$ & 仅相对位置 & 平移 $+$ 旋转 & 邻居相对距离 \\
\bottomrule
\end{tabular}
\end{table}

\noindent 越往下，需要的信息越少（越分布式友好），但控制律越非线性、分析越难。\textbf{Displacement-based 是分布式 MPC 编队最常用的折中}——只需邻居相对位置和一个共同朝向，约束是线性的（相对位移是状态的线性函数），写进 MPC 后保持凸性。本章后续以 displacement-based 为主。

\subsection{三范式在 MPC 中的代价与约束写法}
\label{sec:dmpc-three-in-mpc}

设期望编队几何由一组期望相对位移 $\{\mathbf{d}_{ij}\}$ 指定（$\mathbf{d}_{ij}$ $=$ agent $j$ 相对 agent $i$ 的期望偏移，如一字队形中 $\mathbf{d}_{12}=(1,0,0)^\top$）。

\textbf{(1) Displacement-based 编队代价（软约束形式）.} 在预测时域 $T$ 上，对每条通信边 $(i,j)\in\mathcal{E}$、每个时刻 $t$，惩罚实际相对位移偏离期望：
\begin{equation}
  J_{\text{form}}=\sum_{t=1}^{T}\sum_{(i,j)\in\mathcal{E}} w_{ij}\big\|\,(\mathbf{p}_j(t)-\mathbf{p}_i(t)) - \mathbf{d}_{ij}\,\big\|^2.
  \label{eq:dmpc-disp-soft}
\end{equation}
这是\textbf{二次代价}、关于状态线性，保持 MPC 的凸性（QP）。$w_{ij}$ 是编队权重——越大越“硬”地维持队形。注意 $J_{\text{form}}$ \emph{同时含 $\mathbf{p}_i$ 和 $\mathbf{p}_j$}，这就是耦合的来源：agent $i$ 的最优控制依赖 agent $j$ 的轨迹，反之亦然。

\textbf{(2) Displacement-based 硬约束形式.} 若队形必须严格满足（如刚性连接），写成等式约束：
\begin{equation}
  \mathbf{p}_j(t)-\mathbf{p}_i(t) = \mathbf{d}_{ij},\quad \forall (i,j)\in\mathcal{E},\ \forall t.
  \label{eq:dmpc-disp-hard}
\end{equation}
硬约束保证队形精确，但\emph{牺牲了避障/跟随的灵活性}，且让分布式求解更难（耦合从代价项变成耦合约束，必须用 \cref{sec:dmpc-dual,sec:dmpc-admm} 的对偶/ADMM 严格处理）。实践中常用软约束（大权重）$+$ 关键约束硬化的混合。

\textbf{(3) Distance-based 编队代价.} 只约束距离，保留旋转自由度：
\begin{equation}
  J_{\text{form}}=\sum_{t}\sum_{(i,j)\in\mathcal{E}} w_{ij}\big(\,\|\mathbf{p}_j(t)-\mathbf{p}_i(t)\|^2 - \|\mathbf{d}_{ij}\|^2\,\big)^2.
  \label{eq:dmpc-dist}
\end{equation}
注意这里是\textbf{距离平方的差的平方}（用平方距离避免开方的不可微）。这个代价关于状态是\emph{四次的、非凸的}——这就是 distance-based 更难的代价：MPC 不再是 QP，要用 SQP 或非线性求解器，且可能陷入局部最优（不同旋转构型都满足距离约束）。

\textbf{领航-跟随（leader-follower）作为锚定.} 纯 displacement/distance 约束保留了整体平移/旋转自由度——队形“漂浮”着，整体往哪走没人管。实践中常指定一个\textbf{领航者（leader）}，它跟踪一条全局参考轨迹（由高层任务规划给出），其余\textbf{跟随者（follower）}只维持相对约束。这样既保留分布式（跟随者只看邻居），又锚定整体运动。Kim--Fawcett--Hamed\cite{kim2023layered} 的双四足编队中，两台机器人通过共享负载耦合、整体运动由一条质心参考轨迹引导，是这种思想的体现。

\begin{insight}[三范式 $=$ “信息 vs 自由度 vs 非凸性”权衡的三个采样点]\label{ins:dmpc-three-tradeoff}
编队三范式的差异，本质是\textbf{“用越少的信息换越多的队形自由度，代价是越强的非凸性”}这个权衡的三个采样点。Position-based 用最多信息（全局坐标）换最严约束（零自由度），代价是线性、好解但不分布式；distance-based 用最少信息（仅距离）换最松约束（平移$+$旋转自由），代价是四次非凸、难解但最分布式。Displacement-based 卡在中间——线性凸、平移自由、需朝向一致，是分布式 MPC 的甜点区。选哪个范式，等于在“我能拿到什么信息”和“我能容忍多难的优化”之间定位。
\end{insight}

\begin{insight}[三范式像描述队列的三种口令（标边界）]\label{ins:dmpc-formation-commands}
三范式像描述一个队列的三种口令。Position-based 像“每个人站到地上画好的脚印上”（绝对、精确、需预先画好脚印 $=$ 全局坐标）；displacement-based 像“每个人保持在左边人的右后方 $1$ 米”（相对、需大家都知道哪是“前” $=$ 共同朝向）；distance-based 像“每个人和邻居保持 $1$ 米臂展”（最松、转个方向也成立 $=$ 旋转自由）。\emph{像}的地方：都在描述同一个队形。\emph{不像}的地方：约束强度和所需信息递减，而求解难度（非凸性）递增——这是类比对应不上的关键，口令的“难易”在人听来差不多，但在优化器看来是 QP 与非凸问题的天壤之别。
\end{insight}

\subsection{多足特有：编队约束作用在“质心”还是“机体”？}
\label{sec:dmpc-com-vs-body}

对四足机器人，一个容易被忽略的建模决策是：编队约束约束的是\textbf{质心（CoM）位置 $\mathbf{p}_i$} 还是某个机体参考点？SRB 模型的状态用质心，所以编队约束自然写在质心上。但这带来一个多足特有的微妙点：四足行走时质心相对地面\emph{周期性起伏和摆动}（随踏步节律左右晃、上下颠）。如果编队约束硬性要求两台机器人的瞬时质心相对位移恒等于 $\mathbf{d}_{ij}$，会和各自的踏步动力学\emph{打架}——为维持瞬时队形，机器人可能被迫做出别扭的质心运动，破坏自身平衡。

\begin{insight}[让“编队协调”与“自身平衡”在不同频段上解耦]\label{ins:dmpc-freq-sep-form}
正因为质心随踏步起伏，工程上编队约束通常（a）用\textbf{软约束}（大权重而非硬等式），给质心自然起伏留余地；或（b）约束\textbf{滤波后/平均的}质心位置（去掉踏步频率的高频分量），只维持低频队形趋势；或（c）在 SRB 这个\textbf{降阶模型}层面维持队形，把高频踏步细节留给底层 WBC（\cref{sec:dmpc-layered}）去消化。这不是随意的选择，而是为了让“编队协调”和“自身平衡”这两个目标在不同频段上解耦——队形是低频的群体行为，平衡是高频的个体行为。理解这个频段分离，是多足编队建模区别于轮式/无人机编队的核心。
\end{insight}

\subsection{编队代价的 Laplacian 结构：为什么编队就是“加权共识”}
\label{sec:dmpc-laplacian}

\cref{sec:cs-formation} 说“displacement-based 编队本质是共识”。这一小节把这句话\emph{写成矩阵}，让你看清楚为什么——这是本章把 Laplacian 和编队 MPC 焊接的关键一步，也是 \cref{sec:dmpc-admm} 练习的理论基础。先看最简情形：忽略期望偏移（$\mathbf{d}_{ij}=\mathbf{0}$，即所有机器人要聚到一起），$1$ 维、单时刻，编队代价
\begin{equation}
  J_{\text{form}}=\sum_{(i,j)\in\mathcal{E}}w_{ij}(p_j-p_i)^2.
  \label{eq:dmpc-form-scalar}
\end{equation}
把它对堆叠向量 $\mathbf{p}=[p_1,\dots,p_N]^\top$ 写成二次型。展开 $\sum_{(i,j)}w_{ij}(p_j-p_i)^2$，会发现它恰好等于
\begin{equation}
  \boxed{\ J_{\text{form}}=\mathbf{p}^\top \mathbf{L}_w\,\mathbf{p}\ }
  \label{eq:dmpc-form-laplacian}
\end{equation}
其中 $\mathbf{L}_w$ 是\textbf{加权 Laplacian}（$\mathbf{L}_w=\mathbf{D}_w-\mathbf{A}_w$，$\mathbf{A}_w$ 元素为 $w_{ij}$，$\mathbf{D}_w$ 为加权度对角阵）。这不是巧合——回顾 \cref{ch:mr-consensus}：Laplacian 的二次型 $\mathbf{p}^\top\mathbf{L}\mathbf{p}=\sum_{(i,j)\in\mathcal{E}}w_{ij}(p_j-p_i)^2$ 正是“边上差值的加权平方和”。\textbf{编队代价的 Hessian 就是通信图的加权 Laplacian。} 这立刻解释三件事：
\begin{enumerate}
  \item \textbf{梯度下降编队 $=$ 共识}。对 $J_{\text{form}}=\mathbf{p}^\top \mathbf{L}_w\mathbf{p}$ 做梯度下降 $\dot{\mathbf{p}}=-\nabla J_{\text{form}}=-2\mathbf{L}_w\mathbf{p}$——这\emph{就是} \cref{ch:mr-consensus} 的连续共识协议 $\dot{\mathbf{p}}=-\mathbf{L}\mathbf{p}$（差一个常数因子）。所以“最小化编队代价”和“跑共识”在数学上是同一件事。
  \item \textbf{收敛速率由 $\lambda_2(\mathbf{L}_w)$ 决定}。既然编队梯度流是 $\dot{\mathbf{p}}=-2\mathbf{L}_w\mathbf{p}$，它收敛到一致（$\mathbf{L}_w$ 零空间）的速率由第二小特征值 $\lambda_2$ 决定（\cref{ch:mr-consensus} 结论）——通信图越连通，编队收敛越快。
  \item \textbf{加偏移 $\mathbf{d}_{ij}$ 只是平移最优点}。有期望偏移时 $J_{\text{form}}=\sum w_{ij}((p_j-p_i)-d_{ij})^2$ 展开是 $\mathbf{p}^\top \mathbf{L}_w\mathbf{p}+$（线性项）$+$ 常数。Hessian 仍是 $\mathbf{L}_w$（偏移不改变二次项），只是把最优点从“全相等”平移到“满足期望相对位移”。所以\textbf{编队的“形状”由 $\mathbf{d}_{ij}$ 决定，编队的“收敛动力学”由 $\mathbf{L}_w$（即拓扑）决定}——形状和动力学解耦。
\end{enumerate}

\begin{insight}[编队控制 $=$ 共识 $+$ 偏移]\label{ins:dmpc-form-eq-consensus}
共识让大家“聚拢”（趋于一致），偏移 $\mathbf{d}_{ij}$ 让“一致”变成“保持期望相对位置”。所以编队不是一个独立于共识的新问题，而是\textbf{“带偏移的共识”}——把 \cref{ch:mr-consensus} 的共识协议 $\dot{\mathbf{p}}=-\mathbf{L}\mathbf{p}$ 的“目标”从“全相等”改成“相对位置 $=\mathbf{d}$”。这个洞察的威力在于：\emph{\cref{ch:mr-consensus} 关于共识的一切结论（收敛性、速率 $\lambda_2$、时延阈值 $\tau^\star$、双随机权重）全部直接适用于编队}——你不需要为编队重新发展一套理论。这也是为什么本章 \cref{sec:dmpc-topo} 能把共识理论无缝接到 DMPC 上。
\end{insight}

\begin{insight}[反事实：图不连通则编队裂成几块]\label{ins:dmpc-disconnect}
如果通信图不连通（$\lambda_2=0$，$\mathbf{L}_w$ 有多个零特征值）会怎样？编队梯度流 $\dot{\mathbf{p}}=-2\mathbf{L}_w\mathbf{p}$ 不会收敛到全局一致队形——每个连通分量各自聚拢，分量之间的相对位置无法协调（没有通信链路约束它们）。这对应一个具体故障：编队“裂成几块”，每块内部队形正确，块与块之间漂移。根因就是 $\lambda_2=0$。这把 \cref{ch:mr-consensus} “$\lambda_2>0\iff$ 图连通”的抽象结论，变成了“编队会不会裂开”的工程判据。
\end{insight}

在 MPC 里，这个 Laplacian 结构意味着：编队代价项的 Hessian 块是 $\mathbf{L}_w\otimes \mathbf{I}_3$（$3$ 维空间，每个时刻），稀疏且结构化——这个稀疏性正是 ADMM 能高效分解的原因之一（agent 只和 Laplacian 里相邻的 agent 耦合，对应只和通信邻居协调）。

\subsection{编队与避障的张力：软约束权重的本质}
\label{sec:dmpc-obstacle}

编队 MPC 的一个现实难点（Imran--Hamed 2025 障碍环境分布式 NMPC\cite{imran2025obstacle} 的核心）：\textbf{保持队形和避开障碍常常冲突}。一个障碍挡在编队路径上，要么编队变形绕过（破坏队形），要么硬保持队形撞上去（不可接受）。MPC 如何权衡？这正是\textbf{软约束 vs 硬约束}选择的用武之地：
\begin{itemize}
  \item \textbf{避障必须是硬约束}：$\text{dist}(\mathbf{p}_i,\text{obstacle})\ge r_{\text{safe}}$，绝不能违反（撞了就完了）。
  \item \textbf{编队应该是软约束}：编队代价 $J_{\text{form}}$ 是惩罚项，可以为了避障\emph{临时牺牲}队形。
\end{itemize}
于是 MPC 在避障硬约束下最小化（自身代价 $+$ 编队软代价）——障碍逼近时避障硬约束主导、编队软代价被“暂时放弃”（队形变形绕障）；障碍过后编队软代价重新主导、队形恢复。\textbf{软约束权重 $w_{ij}$ 的本质，就是“为了别的目标（避障、自身平衡）愿意牺牲多少队形”的量化。}

\begin{insight}[约束的软硬划分本身就是一个安全决策]\label{ins:dmpc-soft-hard}
编队 MPC 里“硬约束 vs 软约束”的划分，编码了\textbf{任务的优先级}——硬约束是“绝不能违反的”（避障、动力学、摩擦锥），软约束是“尽量满足、必要时可让步的”（队形、舒适度）。这不是技术细节，而是\textbf{安全与性能的边界}：把什么设成硬约束，等于声明“什么是底线”。多足编队的底线是“不撞 $+$ 不摔”（避障 $+$ 平衡硬约束），队形是“锦上添花”（软约束）。搞反了——把队形设成硬约束、避障设成软惩罚——机器人会为了保持队形撞上障碍，这是灾难性的。\emph{反事实}：若编队设成硬约束、避障设成软惩罚，障碍逼近时编队硬约束强制队形不变、避障软惩罚被压制，MPC 宁可付一点避障惩罚也要保队形，结果\emph{撞上障碍}。所以约束的软硬不是随意的建模偏好，而是直接决定系统在冲突情境下的行为，必须按“安全优先级”严格设定。
\end{insight}

\noindent Imran--Hamed 2025\cite{imran2025obstacle} 用控制屏障函数（CBF，见 \cref{sec:sm-cbf-game}）把避障做成硬约束嵌进分布式 NMPC，正是这个思想的前沿实现——CBF 提供“保证安全（避障）的同时尽量完成任务（编队）”的严格机制。这条“硬安全 $+$ 软任务”的设计原则，贯穿整个多足分布式 MPC 的工程实践。

\begin{pitfall}[概念误区：用绝对位置代价冒充分布式编队]
编队代价写成 $\sum_i\|\mathbf{p}_i-\mathbf{p}_i^{\text{ref}}\|^2$，以为每项只含一个 agent 就是分布式。\emph{后果}：表面上每个 agent 解自己的项，但 $\mathbf{p}_i^{\text{ref}}$ 的生成需要全局信息和中央规划，本质仍是中心化；且队形整体运动僵硬，无法自然避障。\emph{根因}：分布式的判据不是“代价项是否可分”，而是“求解所需信息是否仅来自邻居”。绝对参考点把全局信息藏在了 $\mathbf{p}_i^{\text{ref}}$ 里。\emph{正确做法}：用相对量（displacement/distance）建模，编队代价显含邻居状态 $\mathbf{p}_j$，求解时只需邻居信息；整体运动由领航者或高层参考锚定。
\end{pitfall}

\begin{pitfall}[编程陷阱：distance-based 代价直接用 $\|\mathbf{p}_j-\mathbf{p}_i\|$ 导致不可微]
代价写成 $(\|\mathbf{p}_j-\mathbf{p}_i\|-d_{ij})^2$，含欧氏范数（开方）。\emph{后果}：当 $\mathbf{p}_j\to\mathbf{p}_i$（两机器人重合）时 $\|\cdot\|$ 在原点不可微，梯度爆炸/NaN，求解器崩溃或给出垃圾解。\emph{根因}：欧氏范数 $\|\mathbf{v}\|=\sqrt{\mathbf{v}^\top\mathbf{v}}$ 在 $\mathbf{v}=\mathbf{0}$ 处梯度 $\mathbf{v}/\|\mathbf{v}\|$ 无定义。\emph{正确做法}：用\textbf{平方距离}写代价 $(\|\mathbf{p}_j-\mathbf{p}_i\|^2-d_{ij}^2)^2$（处处可微，多项式）；或对范数加小正则 $\sqrt{\mathbf{v}^\top\mathbf{v}+\epsilon}$。前者更常用。
\end{pitfall}

\begin{pitfall}[概念误区：对四足硬约束瞬时质心相对位置]
要求 $\mathbf{p}_j(t)-\mathbf{p}_i(t)=\mathbf{d}_{ij}$ 在每个 $t$ 严格成立，包括踏步起伏的高频分量。\emph{后果}：机器人为维持瞬时队形被迫做别扭的质心运动，与踏步动力学冲突，自身平衡变差甚至摔倒。\emph{根因}：四足质心随踏步周期性起伏，把高频起伏也纳入硬队形约束，会让“编队”和“平衡”两个目标在高频上互相打架。\emph{正确做法}：用软约束（大权重）给起伏留余地，或只约束低频/滤波后的质心趋势，把高频细节交给底层 WBC。队形是低频群体行为，平衡是高频个体行为，应在频段上分离。
\end{pitfall}

\begin{pitfall}[思维陷阱：忽略编队约束让联合问题“耦合”，以为各 agent 仍可独立解]
建好编队代价后，仍以为每个 agent 可以无视邻居独立求解。\emph{后果}：若真各自独立解（把邻居状态当固定值且从不更新），队形误差无法收敛，或收敛到次优。\emph{根因}：$J_{\text{form}}$ 显含多个 agent 状态，是耦合项；独立求解等于砍掉了协调机制。\emph{正确做法}：认识到编队建模的产物就是“耦合”，处理耦合正是 \cref{sec:dmpc-dual}（对偶分解）和 \cref{sec:dmpc-admm}（ADMM）的任务。建模和求解是两步，本节只负责正确地“制造”耦合，下两节负责“分解”耦合。
\end{pitfall}

\begin{exercise}
\begin{enumerate}
  \item \textbf{（三范式对比实现）} 对 $N=3$ 双积分器一字队形（间距 $1$\,m），分别用 displacement-based 软约束、distance-based 软约束写出编队代价，在一个 QP/NLP 求解器（cvxpy/scipy.optimize）里实现单步优化。给两台机器人初始扰动，观察：displacement-based 恢复到唯一队形；distance-based 可能恢复到\emph{旋转后的}队形（保留旋转自由度）。验证 distance-based 代价的非凸性（不同初值收敛到不同构型）。
  \item \textbf{（耦合结构分析）} 对一字队形 $1$-$2$-$3$（边集 $\mathcal{E}=\{(1,2),(2,3)\}$），写出 displacement-based 编队代价 $J_{\text{form}}$ 关于堆叠变量 $[\mathbf{p}_1,\mathbf{p}_2,\mathbf{p}_3]$ 的 Hessian 稀疏结构。说明为什么 agent $1$ 和 $3$ 不直接耦合（Hessian 对应块为零），但通过 agent $2$ 间接耦合？这个稀疏结构和通信图 Laplacian $\mathbf{L}$ 有什么关系？（提示：Hessian 正比于 $\mathbf{L}\otimes \mathbf{I}$，见 \cref{eq:dmpc-form-laplacian}）
  \item \textbf{（多足建模思考）} 查阅 Kim--Fawcett--Hamed\cite{kim2023layered}（源稿标 arXiv 2211.06913）SRB 网络的状态定义。说明：他们的编队/协同约束作用在质心层面，为何与“把高频踏步交给底层控制器”的分层思想一致？若反过来把编队约束直接加在足端层面会有什么问题？这道题为 \cref{sec:dmpc-srb,sec:dmpc-layered} 铺垫。
\end{enumerate}
\end{exercise}

% ════════════════════════════════════════════════════════════════
\section{对偶分解与价格协调：把耦合约束“卖”给市场}
\label{sec:dmpc-dual}

\paragraph{动机：耦合约束挡在“可分”的路上。}
回顾 \cref{sec:dmpc-formation} 的硬约束编队问题，把它整理成标准的“耦合约束”结构。设 agent $i$ 的本地决策变量是 $\mathbf{X}_i$（它在时域上的轨迹），本地代价 $J_i(\mathbf{X}_i)$，本地约束 $\mathbf{X}_i\in\mathcal{C}_i$（动力学、摩擦锥等，\emph{只含自己}）。编队耦合约束写成线性形式：
\begin{equation}
  \min_{\mathbf{X}_1,\dots,\mathbf{X}_N}\ \sum_{i=1}^N J_i(\mathbf{X}_i)\quad \text{s.t.}\ \mathbf{X}_i\in\mathcal{C}_i\ (\forall i),\quad \underbrace{\sum_{i=1}^N \mathbf{A}_i\mathbf{X}_i = \mathbf{b}}_{\text{耦合约束}}.
  \label{eq:dmpc-coupled}
\end{equation}
这里 $\sum_i \mathbf{A}_i\mathbf{X}_i=\mathbf{b}$ 是把所有编队等式约束（如 $\mathbf{p}_j-\mathbf{p}_i=\mathbf{d}_{ij}$）堆叠成的统一形式。\textbf{关键观察}：如果没有这个耦合约束，目标 $\sum_i J_i(\mathbf{X}_i)$ 和本地约束 $\mathbf{X}_i\in\mathcal{C}_i$ 都是\emph{完全可分}的——每个 agent 独立解 $\min_{\mathbf{X}_i\in\mathcal{C}_i}J_i(\mathbf{X}_i)$ 就行。\textbf{耦合约束是唯一挡在“可分”路上的障碍。} 于是问题变成：能不能把耦合约束“处理掉”，让问题恢复可分？对偶分解的答案是：把约束\emph{移进目标}（拉格朗日松弛），代价是引入一个“价格”。

\paragraph{反面：如果用罚函数硬塞会怎样。}
一个朴素想法是不用对偶，直接把耦合约束做成二次罚项加进目标（就像 \cref{sec:dmpc-formation} 的软约束）：
\begin{equation}
  \min_{\mathbf{X}_i\in\mathcal{C}_i}\ \sum_i J_i(\mathbf{X}_i) + \frac{\rho}{2}\Big\|\sum_i \mathbf{A}_i\mathbf{X}_i-\mathbf{b}\Big\|^2.
  \label{eq:dmpc-penalty}
\end{equation}
这看似可行，但有两个致命问题：
\begin{enumerate}
  \item \textbf{罚项本身不可分}。$\|\sum_i \mathbf{A}_i\mathbf{X}_i-\mathbf{b}\|^2$ 展开后有\emph{交叉项} $\mathbf{X}_i^\top \mathbf{A}_i^\top \mathbf{A}_j\mathbf{X}_j$（$i\ne j$），把所有 agent 又耦合回去了——加罚项非但没解耦，反而制造了新的稠密耦合。
  \item \textbf{罚函数法的老问题}：要让约束精确满足，$\rho$ 必须 $\to\infty$，但 $\rho$ 太大会让问题病态（Hessian 条件数爆炸）。有限的 $\rho$ 只能近似满足约束。
\end{enumerate}
这里先记一个伏笔：\textbf{ADMM 的全部巧妙，就在于它既要用增广（二次罚）项的“强制收敛”好处，又要绕开“罚项不可分”的坏处。} 它通过引入辅助变量 $\mathbf{z}$ 把交叉项拆开——这是 \cref{sec:dmpc-admm} 的核心。现在先看不带增广项的“纯对偶分解”如何实现可分。

\paragraph{历史：从拉格朗日对偶到 Dantzig-Wolfe 与价格协调。}
对偶分解的思想根源是\textbf{拉格朗日对偶}（$19$ 世纪的拉格朗日乘子法 $+$ $20$ 世纪的凸对偶理论）。$1960$ 年代，Dantzig--Wolfe 分解和 Benders 分解把大规模线性规划按结构拆成主问题$+$子问题，奠定了“分解协调（decomposition-coordination）”范式。Everett（$1963$）明确提出用拉格朗日乘子作为“资源价格”协调子系统。$1970$–$80$ 年代，这套方法在经济学、电力系统、网络优化中大放异彩，对偶变量被赋予“影子价格（shadow price）”的经济学解释。Bertsekas 的并行与分布式计算专著（$1989$）系统总结了对偶分解的分布式实现。\pz{Everett 1963、Dantzig-Wolfe/Benders、Bertsekas 1989 的归属系源稿转述，年份/作者未逐一核。}

\begin{insight}[对偶分解的灵魂：把“满足耦合约束”看成“市场出清”]\label{ins:dmpc-market}
对偶分解的灵魂是一个深刻的经济学类比——\textbf{把“满足耦合约束”看成“市场出清（market clearing）”}。耦合约束 $\sum_i \mathbf{A}_i\mathbf{X}_i=\mathbf{b}$ 就像“总供给 $=$ 总需求”；对偶变量 $\boldsymbol{\lambda}$ 就是商品的\textbf{价格}。每个 agent 是一个自利的市场参与者，它根据当前价格 $\boldsymbol{\lambda}$ 决定自己的“生产/消费”（解本地问题）；如果总供需不平衡（约束违反），就调整价格（对偶上升）——供过于求就降价，供不应求就涨价——直到市场出清（约束满足）。这不是修辞，而是数学上严格成立的对应：对偶子梯度法的每一步，就是一次“价格调整”。
\end{insight}

\subsection{拉格朗日松弛 \texorpdfstring{$\to$}{->} 可分的对偶函数 \texorpdfstring{$\to$}{->} 价格迭代}
\label{sec:dmpc-dual-derive}

现在严格推导。对耦合约束 $\sum_i \mathbf{A}_i\mathbf{X}_i=\mathbf{b}$ 引入拉格朗日乘子 $\boldsymbol{\lambda}$，构造拉格朗日函数（\textbf{注意：本地约束 $\mathbf{X}_i\in\mathcal{C}_i$ 不松弛，保留}）：
\begin{equation}
  \mathcal{L}(\mathbf{X}_1,\dots,\mathbf{X}_N,\boldsymbol{\lambda}) = \sum_i J_i(\mathbf{X}_i) + \boldsymbol{\lambda}^\top\Big(\sum_i \mathbf{A}_i\mathbf{X}_i-\mathbf{b}\Big).
  \label{eq:dmpc-lagrangian}
\end{equation}

\textbf{第一步：分组。} 把 $\boldsymbol{\lambda}^\top\sum_i \mathbf{A}_i\mathbf{X}_i=\sum_i\boldsymbol{\lambda}^\top \mathbf{A}_i\mathbf{X}_i$ 拆开，按 agent 重组：
\begin{equation}
  \mathcal{L} = \sum_i\big[\underbrace{J_i(\mathbf{X}_i)+\boldsymbol{\lambda}^\top \mathbf{A}_i\mathbf{X}_i}_{\text{只含 }\mathbf{X}_i}\big] - \boldsymbol{\lambda}^\top\mathbf{b}.
  \label{eq:dmpc-lagrangian-split}
\end{equation}
\textbf{这是对偶分解的关键一步}：松弛耦合约束后，拉格朗日函数关于 $\{\mathbf{X}_i\}$ \emph{完全可分}了！因为耦合约束被“价格化”了——它不再是一个把变量绑在一起的硬约束，而是变成了每个 agent 目标里的一个\textbf{线性价格项} $\boldsymbol{\lambda}^\top \mathbf{A}_i\mathbf{X}_i$。

\textbf{第二步：对偶函数 $=$ 各 agent 独立最小化。} 对偶函数 $g(\boldsymbol{\lambda})=\min_{\mathbf{X}_i\in\mathcal{C}_i}\mathcal{L}$。由可分性：
\begin{equation}
  g(\boldsymbol{\lambda}) = \sum_i\underbrace{\Big[\min_{\mathbf{X}_i\in\mathcal{C}_i}\ J_i(\mathbf{X}_i)+\boldsymbol{\lambda}^\top \mathbf{A}_i\mathbf{X}_i\Big]}_{\text{agent }i\text{ 的本地子问题 }g_i(\boldsymbol{\lambda})} - \boldsymbol{\lambda}^\top\mathbf{b}.
  \label{eq:dmpc-dualfunc}
\end{equation}
\textbf{每个 agent 独立解一个本地子问题} $g_i(\boldsymbol{\lambda})=\min_{\mathbf{X}_i\in\mathcal{C}_i}J_i(\mathbf{X}_i)+\boldsymbol{\lambda}^\top \mathbf{A}_i\mathbf{X}_i$——这就是“给定价格 $\boldsymbol{\lambda}$，我（agent $i$）的最优生产计划”。子问题只含本地变量、本地约束，规模是单体大小，可并行。

\textbf{第三步：价格更新（对偶上升 / 子梯度法）。} 对偶问题是 $\max_{\boldsymbol{\lambda}}g(\boldsymbol{\lambda})$（$g$ 凹，无约束）。$g$ 的（次）梯度有一个漂亮的形式：
\begin{equation}
  \nabla g(\boldsymbol{\lambda}) = \sum_i \mathbf{A}_i\mathbf{X}_i^*(\boldsymbol{\lambda}) - \mathbf{b} = \text{（当前价格下的）耦合约束残差},
  \label{eq:dmpc-dualgrad}
\end{equation}
其中 $\mathbf{X}_i^*(\boldsymbol{\lambda})$ 是 agent $i$ 在价格 $\boldsymbol{\lambda}$ 下的本地最优解。于是价格按子梯度上升更新：
\begin{equation}
  \boxed{\ \boldsymbol{\lambda}^{k+1} = \boldsymbol{\lambda}^k + \alpha^k\Big(\sum_i \mathbf{A}_i\mathbf{X}_i^*(\boldsymbol{\lambda}^k)-\mathbf{b}\Big)\ }
  \label{eq:dmpc-price-update}
\end{equation}
$\alpha^k$ 是步长。这个更新的经济学含义完美对应市场：\textbf{残差 $\sum_i \mathbf{A}_i\mathbf{X}_i^*-\mathbf{b}$ 就是“供需失衡”}——大于 $0$（供过于求）则价格上调，小于 $0$（供不应求）则价格下调，直到残差为 $0$（市场出清，约束满足）。完整算法见 \cref{alg:dmpc-dual}。

\begin{algo}[对偶分解（分布式）]\label{alg:dmpc-dual}
重复直到收敛：(1) 广播当前价格 $\boldsymbol{\lambda}^k$ 给所有 agent（或在通信图上传播）；(2) 每个 agent $i$ 并行解本地子问题 $\mathbf{X}_i^*=\argminop_{\mathbf{X}_i\in\mathcal{C}_i} J_i(\mathbf{X}_i)+(\boldsymbol{\lambda}^k)^\top \mathbf{A}_i\mathbf{X}_i$；(3) 收集残差 $\mathbf{r}=\sum_i \mathbf{A}_i\mathbf{X}_i^*-\mathbf{b}$（需某种聚合：中央求和或图上共识）；(4) 价格更新 $\boldsymbol{\lambda}^{k+1}=\boldsymbol{\lambda}^k+\alpha^k\mathbf{r}$。
\end{algo}

\noindent 注意第 $3$ 步的残差聚合 $\sum_i \mathbf{A}_i\mathbf{X}_i^*$ 需要把各 agent 的贡献加起来——这一步如何分布式实现（用 \cref{ch:mr-consensus} 的\textbf{共识平均}而非中央求和），正是 \cref{sec:dmpc-topo} 的话题，也是对偶分解/ADMM 与共识算法的接口。

\subsection{对偶分解为什么“慢”且“脆”——这正是 ADMM 的动机}
\label{sec:dmpc-dual-weak}

对偶分解优雅，但有三个实际痛点，理解它们就理解了为什么工程上几乎都用 ADMM 而非纯对偶分解：
\begin{enumerate}
  \item \textbf{子梯度法收敛慢（$O(1/\sqrt{k})$ 或 $O(1/k)$）}。对偶函数 $g$ 一般不可微（在最优解附近有“折点”），只能用子梯度法，收敛速率慢，且需精心调步长 $\alpha^k$（递减步长 $\sum\alpha^k=\infty,\sum(\alpha^k)^2<\infty$）。对需要每个控制周期（毫秒级）内收敛的 MPC，这个速度往往不够。
  \item \textbf{要求严格凸才有唯一本地解}。如果本地子问题 $J_i+\boldsymbol{\lambda}^\top \mathbf{A}_i\mathbf{X}_i$ 不严格凸（如线性 $J_i$），$\mathbf{X}_i^*(\boldsymbol{\lambda})$ 可能不唯一甚至“跳变”（价格微小变化导致最优解大幅跳动），对偶迭代剧烈震荡。
  \item \textbf{中间迭代不可行}。对偶分解只在收敛时满足耦合约束；中间任何一轮的 $\{\mathbf{X}_i^*\}$ \emph{不满足}编队约束。如果 MPC 被迫提前停（时间到了），拿到的是一个违反队形的解。
\end{enumerate}

\begin{insight}[ADMM 不是新方法，而是“对偶分解打了三个补丁”]\label{ins:dmpc-admm-patches}
ADMM 不是“另一种全新的分解方法”，\textbf{而是“对偶分解打了三个补丁”}——（补丁 1）加增广二次项，让对偶函数变得可微、收敛加速到线性率；（补丁 2）增广项提供正则，即使本地不严格凸也能稳定；（补丁 3）引入辅助变量 $\mathbf{z}$ 解决增广项的不可分。换句话说，\textbf{你必须先理解对偶分解的思想（价格协调）和它的三个弱点，才能理解 ADMM 的每一处设计都在精确地修复哪个弱点}。这就是本章先讲对偶分解再讲 ADMM 的原因——跳过这一节直接学 ADMM，你只会背三步公式，不懂它为什么是这三步。
\end{insight}

\begin{insight}[影子价格在多足里是可测的真实力]\label{ins:dmpc-shadow-force}
尽管纯对偶分解工程上少用，它的“价格”视角在多足编队里有直接的物理对应。当两台四足通过刚性杆连接，杆里的\textbf{约束力 $\boldsymbol{\lambda}$}（拉格朗日乘子，\cref{sec:dmpc-srb}）就是对偶变量——它正是“维持刚性连接这个耦合约束”的价格：杆受拉是正价格、受压是负价格。所以对偶分解里抽象的“影子价格”，在多足编队里是一个你能用力传感器测到的真实的力。这个对应不是巧合：约束力在物理上就是“维持约束所需的力”，在数学上就是“维持约束的拉格朗日乘子”，二者是同一个 $\boldsymbol{\lambda}$。
\end{insight}

\subsection{走一遍：两 agent 的价格协调过程（手算 \texorpdfstring{$+$}{+} 解读）}
\label{sec:dmpc-dual-example}

为了让“价格协调”从抽象变具体，把整个对偶迭代在一个最小例子上手算一遍。两个 agent，本地代价 $J_1(x_1)=\tfrac12(x_1-3)^2$、$J_2(x_2)=\tfrac12(x_2-1)^2$（agent $1$ 想去 $3$、agent $2$ 想去 $1$），耦合约束 $x_1=x_2$（要求一致，写成 $x_1-x_2=0$，即 $A_1=1,A_2=-1,b=0$），无本地约束。

\textbf{最优解（先用集中式对照）}：$\min\tfrac12(x_1-3)^2+\tfrac12(x_2-1)^2$ s.t.\ $x_1=x_2$，令 $x_1=x_2=x$，对 $x$ 求导 $(x-3)+(x-1)=0\Rightarrow x=2$。所以联合最优是两 agent 都到 $2$（两愿望的中点）。

\textbf{对偶分解过程}：拉格朗日 $\mathcal{L}=\tfrac12(x_1-3)^2+\tfrac12(x_2-1)^2+\lambda(x_1-x_2)$。
\begin{itemize}
  \item \textbf{本地子问题}（给定价格 $\lambda$）：agent $1$：$\min_{x_1}\tfrac12(x_1-3)^2+\lambda x_1\Rightarrow x_1^*(\lambda)=3-\lambda$（求导 $(x_1-3)+\lambda=0$）；agent $2$：$\min_{x_2}\tfrac12(x_2-1)^2-\lambda x_2\Rightarrow x_2^*(\lambda)=1+\lambda$。
  \item \textbf{残差（次梯度）}：$r(\lambda)=x_1^*-x_2^*=(3-\lambda)-(1+\lambda)=2-2\lambda$。
  \item \textbf{价格更新}：$\lambda^{k+1}=\lambda^k+\alpha(2-2\lambda^k)$。
\end{itemize}
手算迭代（步长 $\alpha=0.25$，从 $\lambda^0=0$）见 \cref{tab:dmpc-price-iter}。

\begin{table}[H]\centering\small
\caption{两 agent 价格协调的手算迭代（$\alpha=0.25$）}
\label{tab:dmpc-price-iter}
\begin{tabular}{cccccp{0.30\linewidth}}
\toprule
$k$ & $\lambda^k$ & $x_1^*=3-\lambda$ & $x_2^*=1+\lambda$ & 残差 $r=x_1-x_2$ & 解读 \\
\midrule
$0$ & $0.00$ & $3.00$ & $1.00$ & $2.00$ & 价格为 $0$，各去各的愿望，相差 $2$（违约最大） \\
$1$ & $0.50$ & $2.50$ & $1.50$ & $1.00$ & 价格上调，两者被拉近 \\
$2$ & $0.75$ & $2.25$ & $1.75$ & $0.50$ & 继续拉近，残差减半 \\
$3$ & $0.875$ & $2.125$ & $1.875$ & $0.25$ & 接近一致 \\
$\cdots$ & $\to1.00$ & $\to2.00$ & $\to2.00$ & $\to0$ & 收敛：价格 $\lambda^*=1$，两 agent 都到 $2$ \\
\bottomrule
\end{tabular}
\end{table}

\textbf{解读（价格协调的全过程）}：
\begin{enumerate}
  \item \textbf{初始 $\lambda=0$（无价格）}：耦合约束不起作用，agent $1$ 去 $3$、agent $2$ 去 $1$——这正是“完全去中心”的结果（\cref{sec:dmpc-why} 反面），相差 $2$。
  \item \textbf{价格逐步上调}：残差 $r=2>0$（agent $1$ 在 agent $2$ 右边，违反一致），价格 $\lambda$ 上调。$\lambda>0$ 对 agent $1$ 是“惩罚 $x_1$ 太大”（$+\lambda x_1$），把它往左拉；对 agent $2$ 是“奖励 $x_2$ 增大”（$-\lambda x_2$），把它往右拉。一拉一推，两者靠拢。
  \item \textbf{收敛到 $\lambda^*=1$}：此时 $x_1=x_2=2$，残差 $0$，市场出清。\textbf{$\lambda^*=1$ 的物理含义}：它是维持“$x_1=x_2$”这个约束所需的“价格”——也是该约束的影子价格（最优值对约束右端 $b$ 的灵敏度）。
\end{enumerate}

\begin{insight}[对偶迭代把“无协调”和“完美协调”放在同一条线的两端]\label{ins:dmpc-coordination-line}
\cref{tab:dmpc-price-iter} 把 \cref{sec:dmpc-why} 反面的“完全去中心”和本章的“分布式协调”放在了同一条迭代线的\textbf{两端}——$\lambda=0$（第 $0$ 行）就是完全去中心（各去愿望、队形崩），$\lambda=\lambda^*$（最后一行）就是协调到联合最优（队形成）。\textbf{对偶迭代的本质，就是从“无协调”通过逐步调价走到“完美协调”的过程。} 价格 $\lambda$ 是协调强度的旋钮：$\lambda=0$ 无协调，$\lambda$ 增大协调加强，$\lambda^*$ 恰好协调。这让你直观看到“协调”在数学上到底是什么——不是模糊的概念，而是一个被价格驱动的、可量化的、逐步收敛的过程。
\end{insight}

\begin{insight}[反事实：步长过大则价格“调过头”反复震荡]\label{ins:dmpc-stepsize}
如果步长 $\alpha$ 取太大（如 $\alpha=1.5$）会怎样？$\lambda^{k+1}=\lambda^k+1.5(2-2\lambda^k)$，从 $\lambda=0$：$\lambda^1=3$，$x_1=0,x_2=4$，残差 $-4$（\emph{反向越过了}，agent $1$ 跑到 agent $2$ 左边）；继续震荡发散。对照 $\alpha=0.25$ 的平稳收敛，你能亲手感受步长对子梯度法的敏感——也理解了为什么 ADMM 用固定 $\rho$ 却稳：增广项给了额外曲率，不靠精调步长。
\end{insight}

\begin{pitfall}[概念误区：松弛本地约束 $\mathbf{X}_i\in\mathcal{C}_i$]
做拉格朗日松弛时，把所有约束（包括本地动力学、摩擦锥）都松弛掉。\emph{后果}：本地子问题变成无约束最小化，解出的 $\mathbf{X}_i$ 违反动力学/摩擦锥，物理上无意义；且失去“本地子问题保持单体 MPC 结构”的好处。\emph{根因}：对偶分解只需松弛\textbf{耦合}约束（把变量绑在一起的那些），本地约束应保留在各自子问题里。松弛本地约束既无必要（它本来就可分）又有害（丢失结构、解不可行）。\emph{正确做法}：只松弛耦合约束 $\sum_i \mathbf{A}_i\mathbf{X}_i=\mathbf{b}$，本地约束 $\mathbf{X}_i\in\mathcal{C}_i$ 原样留在 agent $i$ 的子问题中。区分“耦合约束”与“本地约束”是对偶分解的第一要务。
\end{pitfall}

\begin{pitfall}[概念误区：把对偶变量 $\boldsymbol{\lambda}$ 当成可随便初始化、随便调的超参数]
以为 $\boldsymbol{\lambda}$ 像学习率一样是个调参旋钮，初值无所谓。\emph{后果}：$\boldsymbol{\lambda}$ 初值离最优价格太远时，对偶迭代要很多轮才收敛；在 MPC 里每周期冷启动 $\boldsymbol{\lambda}=\mathbf{0}$ 会浪费大量迭代。\emph{根因}：$\boldsymbol{\lambda}$ 是有物理含义的（影子价格/约束力），它的最优值由问题决定，不是自由超参；但它\textbf{可以 warm-start}——相邻控制周期的最优价格变化不大。\emph{正确做法}：MPC 中用上一周期收敛的 $\boldsymbol{\lambda}$ 作为本周期初值（warm-start），大幅减少迭代轮数。这是把对偶/ADMM 用于实时 MPC 的关键工程技巧。
\end{pitfall}

\begin{pitfall}[思维陷阱：用对偶分解时假设中间迭代的解可以直接执行]
MPC 时间到了就拿当前轮的 $\{\mathbf{X}_i^*\}$ 去执行，以为“差不多收敛了”。\emph{后果}：中间解不满足耦合约束，机器人队形违反（如刚性约束未满足导致对负载施力冲突）。\emph{根因}：对偶分解只保证收敛时可行，中间迭代的原始变量不满足耦合约束（这是对偶法的固有性质，区别于原始可行方法）。\emph{正确做法}：要么跑到足够收敛（残差 $<$ 阈值）再执行，要么改用对中间解更友好的方法（ADMM 的原始残差随轮数下降，中间解的约束违反更小、更可控）。这也是 ADMM 比纯对偶分解更适合实时 MPC 的原因之一。
\end{pitfall}

\begin{pitfall}[编程陷阱：子梯度法步长 $\alpha^k$ 用固定常数]
价格更新 $\boldsymbol{\lambda}^{k+1}=\boldsymbol{\lambda}^k+\alpha\cdot \mathbf{r}$ 用固定 $\alpha$。\emph{后果}：固定步长下子梯度法\emph{不保证收敛到最优}，只能收敛到最优值附近的一个邻域内震荡。\emph{根因}：子梯度（区别于梯度）方向不一定是下降方向，固定步长无法消除震荡。理论要求递减步长 $\alpha^k\to0$ 且 $\sum\alpha^k=\infty$。\emph{正确做法}：用递减步长（如 $\alpha^k=a/(b+k)$）保证收敛；或干脆换 ADMM——ADMM 的“步长”就是固定的 $\rho$ 却能线性收敛，因为增广项提供了额外的曲率信息。
\end{pitfall}

\begin{exercise}
\begin{enumerate}
  \item \textbf{（手推对偶分解）} 考虑两个 agent 的简单耦合 QP：$\min \tfrac12 x_1^2+\tfrac12 x_2^2$ s.t.\ $x_1+x_2=2$（$A_1=A_2=1,\ b=2$，无本地约束）。(a) 写出拉格朗日函数并验证它关于 $(x_1,x_2)$ 可分；(b) 解本地子问题得 $x_i^*(\lambda)$；(c) 写出对偶函数 $g(\lambda)$ 并解析求最优 $\lambda^*$；(d) 验证 $x_1^*(\lambda^*)=x_2^*(\lambda^*)=1$。
  \item \textbf{（价格迭代仿真）} 把上题用对偶子梯度法数值求解：从 $\lambda^0=0$ 开始，用步长 $\alpha^k=0.5/(1+k)$ 和固定步长 $\alpha=0.5$ 各跑 $50$ 轮，画出 $\lambda^k$ 和残差 $r^k=x_1^*+x_2^*-2$ 的收敛曲线。验证：固定步长在最优附近震荡不收敛，递减步长收敛。
  \item \textbf{（影子价格的物理意义）} 结合 \cref{sec:dmpc-srb} 预告：把上题的耦合约束 $x_1+x_2=2$ 想象成“两台机器人通过一根杆维持总位移为 $2$”。最优对偶变量 $\lambda^*$ 对应杆里的什么物理量？若把约束改成 $x_1+x_2=2.5$（杆变长），$\lambda^*$ 会怎么变？用对偶变量的灵敏度含义（$\lambda^*=\partial f^*/\partial b$）解释。
\end{enumerate}
\end{exercise}

% ════════════════════════════════════════════════════════════════
\section{ADMM 分布式 MPC：又快又稳的分解}
\label{sec:dmpc-admm}

\paragraph{动机：既要可分，又要快、稳、中间可用。}
\cref{sec:dmpc-dual} 给了对偶分解的思想和它的三个弱点（慢、脆、中间不可行）。这一节用 ADMM（交替方向乘子法）逐一修复它们，并把它落成多足分布式 MPC 的可执行算法。这是全章的\textbf{算法核心}——你会看到 \cref{sec:cs-admm} 学的 ADMM 三步迭代，在这里变成“每个 agent 解一个带二次惩罚的本地 MPC $+$ 邻居预测的共识平均 $+$ 对偶上升”。ADMM 的设计目标见 \cref{tab:dmpc-admm-fix}。

\begin{table}[H]\centering\small
\caption{对偶分解的弱点与 ADMM 的修复}
\label{tab:dmpc-admm-fix}
\begin{tabular}{p{0.44\linewidth}p{0.48\linewidth}}
\toprule
对偶分解的弱点 & ADMM 的修复 \\
\midrule
子梯度法慢（$O(1/\sqrt k)$），需调步长 & 增广二次项让对偶可微，固定 $\rho$ 即可，凸情形线性收敛 \\
本地不严格凸则震荡 & 增广项提供正则（强凸化本地子问题），稳定 \\
中间迭代不满足耦合约束 & 原始残差单调下降，中间解约束违反更小、更可控 \\
\bottomrule
\end{tabular}
\end{table}

\paragraph{反面：如果直接给联合 MPC 加增广项会怎样——回到“不可分”。}
\cref{sec:dmpc-dual} 的反面已指出：直接给耦合约束加二次罚 $\tfrac{\rho}{2}\|\sum_i \mathbf{A}_i\mathbf{X}_i-\mathbf{b}\|^2$，展开后的交叉项 $\mathbf{X}_i^\top \mathbf{A}_i^\top \mathbf{A}_j\mathbf{X}_j$ 把 agent 又耦合回去。所以\textbf{不能直接加}。ADMM 的破局点是：\textbf{引入一个辅助变量 $\mathbf{z}$，让每个 agent 的变量只和 $\mathbf{z}$ 比较，而不直接和别的 agent 比较}——交叉项被 $\mathbf{z}$“隔开”了。这就是 consensus-ADMM 的核心技巧。

\paragraph{历史：从 Glowinski/Gabay 到 Boyd 的 consensus-ADMM 复兴。}
ADMM 由 Glowinski--Marrocco 和 Gabay--Mercier 在 $1970$ 年代提出，最初用于偏微分方程的数值解（变分不等式）。它沉寂了几十年，直到 $2011$ 年 Boyd 等人的著名综述\cite{boyd2011admm} 把它重新带火——该文系统阐述了 \textbf{consensus-ADMM（一致性 ADMM）}：把“$N$ 个 agent 的变量要一致”建模为“每个 agent 有本地副本 $+$ 一个全局共识变量 $\mathbf{z}$，约束本地副本 $=\mathbf{z}$”。这个形式天然适配分布式机器学习和多 agent 优化，成为此后十年分布式优化的主力工具。Summers--Lygeros、Conte 等人随后把它用于分布式 MPC\cite{conte2016distributed}，Hamed 团队把它用于多足协同\cite{kim2023layered}。

\begin{insight}[用一个虚拟全局变量 $\mathbf{z}$ 把多边耦合化为单边耦合]\label{ins:dmpc-z-star}
consensus-ADMM 的核心创意是\textbf{“用一个虚拟的全局变量 $\mathbf{z}$ 当所有 agent 的‘公共参照物’”}。直接让 $N$ 个 agent 两两一致，是 $O(N^2)$ 条耦合关系（稠密、不可分）；引入 $\mathbf{z}$ 后，每个 agent 只需和 $\mathbf{z}$ 一致，是 $N$ 条关系（星形、可分）。$\mathbf{z}$ 把“多边耦合”转化为“每个 agent 对中心的单边耦合”，而这个“中心” $\mathbf{z}$ 的更新恰好是一个\textbf{平均}操作——于是在分布式实现里，$\mathbf{z}$ 不需要真有一个中央节点，用 \cref{ch:mr-consensus} 的共识平均就能算出来。这就是为什么 ADMM 协调步“本质是共识”。
\end{insight}

\subsection{consensus-ADMM 形式的分布式 MPC}
\label{sec:dmpc-admm-derive}

现在严格推导。用 consensus 形式重写编队 MPC，核心思想：\textbf{为耦合的“接口变量”建立一个全局共识副本}。

\textbf{第一步：识别接口变量并引入共识副本。} 编队耦合发生在机器人的某些“接口量”上——可能是相对位姿、共享负载状态、或边上的相对位移。把 agent $i$ 涉及耦合的那部分变量记为 $\mathbf{w}_i(\mathbf{X}_i)$（接口变量，如 agent $i$ 对共享负载状态的本地估计、或它在编队边上的位置贡献）。引入全局共识变量 $\mathbf{z}$，要求所有相关 agent 的接口副本与 $\mathbf{z}$ 一致：
\begin{equation}
  \min_{\{\mathbf{X}_i\},\mathbf{z}}\ \sum_i J_i(\mathbf{X}_i)\quad\text{s.t.}\ \mathbf{X}_i\in\mathcal{C}_i,\quad \underbrace{\mathbf{w}_i(\mathbf{X}_i)=\mathbf{z}_i\ (\forall i)}_{\text{一致性约束}}.
  \label{eq:dmpc-consensus-form}
\end{equation}
这里 $\mathbf{z}_i$ 是 $\mathbf{z}$ 中与 agent $i$ 相关的分量（如果是全局共享负载，所有 $\mathbf{z}_i$ 指向同一个 $\mathbf{z}$；如果是边耦合，$\mathbf{z}_i$ 是边变量）。\textbf{关键}：现在每个 agent 的变量只通过“$=\mathbf{z}_i$”和全局变量绑定，agent 之间不再直接耦合。

\textbf{第二步：增广拉格朗日。} 对一致性约束 $\mathbf{w}_i(\mathbf{X}_i)=\mathbf{z}_i$ 引入对偶变量 $\mathbf{y}_i$ 和增广项：
\begin{equation}
  \mathcal{L}_\rho = \sum_i\Big[J_i(\mathbf{X}_i) + \mathbf{y}_i^\top(\mathbf{w}_i(\mathbf{X}_i)-\mathbf{z}_i) + \tfrac{\rho}{2}\|\mathbf{w}_i(\mathbf{X}_i)-\mathbf{z}_i\|^2\Big].
  \label{eq:dmpc-aug-lag}
\end{equation}
注意：增广项 $\|\mathbf{w}_i-\mathbf{z}_i\|^2$ \emph{只含 agent $i$ 的变量和全局 $\mathbf{z}_i$，没有 agent 间的交叉项}——这正是引入 $\mathbf{z}$ 的回报。

\textbf{第三步：ADMM 三步交替。} ADMM 轮流优化 $\{\mathbf{X}_i\}$、$\mathbf{z}$、$\{\mathbf{y}_i\}$。用缩放对偶变量 $\mathbf{u}_i=\mathbf{y}_i/\rho$（缩放形式更简洁）：

\textbf{($\mathbf{X}$-更新) 每个 agent 并行解本地 MPC（带二次惩罚）}：
\begin{equation}
  \boxed{\ \mathbf{X}_i^{k+1}=\argminop_{\mathbf{X}_i\in\mathcal{C}_i}\ J_i(\mathbf{X}_i)+\tfrac{\rho}{2}\big\|\mathbf{w}_i(\mathbf{X}_i)-\mathbf{z}_i^k+\mathbf{u}_i^k\big\|^2\ }
  \label{eq:dmpc-x-update}
\end{equation}
这就是\textbf{单体 MPC 加一个把接口变量拉向共识值 $\mathbf{z}_i^k-\mathbf{u}_i^k$ 的二次惩罚}。规模是单体大小，可并行，且因增广项是二次的，本地问题仍是 QP（凸 SRB MPC 情形）。

\textbf{($\mathbf{z}$-更新) 共识平均}：
\begin{equation}
  \boxed{\ \mathbf{z}^{k+1}=\frac{1}{|\mathcal{I}_z|}\sum_{i\in\mathcal{I}_z}\big(\mathbf{w}_i(\mathbf{X}_i^{k+1})+\mathbf{u}_i^k\big)\ }
  \label{eq:dmpc-z-update}
\end{equation}
全局共识变量更新为相关 agent 的接口副本（加对偶偏移）的\textbf{平均}——$\mathcal{I}_z$ 是与该 $\mathbf{z}$ 分量相关的 agent 集合。\textbf{这一步就是 \cref{ch:mr-consensus} 的共识平均！} 分布式实现时不需要中央节点，用共识迭代算这个平均（\cref{sec:dmpc-topo}）。

\textbf{($\mathbf{y}$-更新) 对偶上升}：
\begin{equation}
  \boxed{\ \mathbf{u}_i^{k+1}=\mathbf{u}_i^k+\mathbf{w}_i(\mathbf{X}_i^{k+1})-\mathbf{z}_i^{k+1}\ }
  \label{eq:dmpc-y-update}
\end{equation}
对偶变量累积“本地副本与共识值的偏差”——和 \cref{sec:dmpc-dual} 的价格更新同理，但这里因增广项的存在，收敛是线性的。

\begin{insight}[算法骨架不变，规模和物理含义变了]\label{ins:dmpc-same-skeleton}
到这里我们得到了完整的 consensus-ADMM 分布式 MPC：$\mathbf{X}$-更新（本地 MPC $+$ 二次惩罚，并行）、$\mathbf{z}$-更新（共识平均，邻居协调）、$\mathbf{y}$-更新（对偶上升，逼迫一致）。把这三步和 \cref{sec:cs-admm} 的标准 ADMM 对照，你会发现\textbf{结构完全一样}，只是 $x\to\mathbf{X}_i$（轨迹）、$z\to\mathbf{z}$（共识接口）、惩罚项作用在接口变量 $\mathbf{w}_i$ 上。把这套三步嵌进 MPC 的滚动框架，就是分布式 MPC 的完整循环：每个控制周期，各 agent 跑若干轮 ADMM（通常 warm-start $+$ 固定 $3$–$5$ 轮，而非跑到收敛），取第一步控制执行，滚动到下一周期。
\end{insight}

\subsection{ADMM 分布式 MPC 协调器的实现要点}
\label{sec:dmpc-admm-impl}

把三步落成代码时，最重要的是把“本地 MPC 求解”和“共识/对偶更新”分成独立步骤——这正对应分布式系统的两种操作：本地 MPC 求解是纯计算（CPU 密集，可完全并行），共识平均 $+$ 对偶更新是通信 $+$ 轻量计算（需邻居数据）。分开后：(a) 计算和通信的边界清晰，便于在真机上把“算”放本地、“传”走网络；(b) 可独立优化各步——本地 MPC 用 warm-start 加速，通信步用共识压缩轮数；(c) 固定迭代轮数（不跑到收敛）成为可能——每轮就是一次“算$+$传”，时间可预算。算法骨架见 \cref{alg:dmpc-admm-loop}，C++ 协调器的核心循环见下面的代码框。

\begin{algo}[consensus-ADMM 分布式 MPC（方式 C：固定 $K$ 轮 $+$ warm-start）]\label{alg:dmpc-admm-loop}
每控制周期：用上周期解 warm-start $\mathbf{z},\mathbf{u}$ 及各本地解；for $k=1\ldots K$（$K=3$–$5$）：\;(1) $\mathbf{X}$-更新（本地、并行、可 warm-start）解 \cref{eq:dmpc-x-update}；(2) $\mathbf{z}$-更新（通信、共识平均）\cref{eq:dmpc-z-update}；(3) $\mathbf{y}$-更新（本地、对偶累积）\cref{eq:dmpc-y-update}；(4) 算残差 $r^k,s^k$，若均 $<$ 阈值则提前停。取第一步控制执行，滚动。
\end{algo}

\begin{codebox}[C++]{consensus-ADMM 协调器核心循环（接口变量为全局共享量，如负载状态）}
// 跑 K 轮 ADMM（实时 MPC 中 K 固定，如 3~5;不跑到收敛）
void run(std::array<AgentData, N>& ag, double rho, int K) {
  for (int k = 0; k < K; ++k) {
    // X-更新:每个 agent 并行解本地 MPC（带二次惩罚）
    #pragma omp parallel for
    for (int i = 0; i < N; ++i)         // min J_i + (rho/2)||w_i - z_i + u_i||^2
      solve_local_mpc(i, ag[i], rho);   // warm-start 用上一轮解初始化
    // Z-更新:共识平均 average(w_i + u_i)  ← 关键:含 u_i！
    for (int t = 0; t < horizon_; ++t) {
      State avg = State::Zero();
      for (int i = 0; i < N; ++i) avg += ag[i].w_local[t] + ag[i].u_dual[t];
      avg /= N;
      for (int i = 0; i < N; ++i) ag[i].z_cons[t] = avg;   // 全局共享 z
    }
    // Y-更新:对偶上升 u_i += w_i - z
    for (int i = 0; i < N; ++i)
      for (int t = 0; t < horizon_; ++t)
        ag[i].u_dual[t] += ag[i].w_local[t] - ag[i].z_cons[t];
    auto [r, s] = compute_residuals(ag, rho);   // 残差监控 / 自适应 rho
    if (r < eps_pri_ && s < eps_dual_) break;    // 收敛则提前停
  }
}
\end{codebox}

\noindent 三个最常见的错误写法值得点名（对应 \cref{sec:dmpc-admm-impl} 末尾的陷阱）：(1) $\mathbf{z}$-更新里漏掉对偶项、只平均 $\mathbf{w}_i$——consensus-ADMM 的 $\mathbf{z}$-更新是 \verb@average(w_i + u_i)@、不是 \verb@average(w_i)@，漏掉 $\mathbf{u}_i$ 破坏对偶配平，残差降不到 $0$、编队有稳态偏差（极难发现，因为“看起来在平均”）；(2) 每控制周期把 $\mathbf{z},\mathbf{u}$ 清零冷启动——固定 $K=3$–$5$ 根本不够收敛、编队跟踪滞后；(3) $\rho$ 取固定且过大（如 $\rho=10^4$）期望“收敛更快”——二次惩罚主导、本地 MPC 几乎只顾贴近 $\mathbf{z}$ 而忽略自身代价 $J_i$，原始残差掉得快但解质量差、甚至本地 QP 病态。

\textbf{三种 ADMM 调度方式}（\cref{tab:dmpc-schedule}）：同步 ADMM（所有 agent 每轮严格对齐，等价标准 ADMM、凸情形收敛有保证，但被最慢 agent 拖累、要全局时钟）；异步 ADMM（agent 算完就更新、不等别人，无需全局同步，但收敛分析更难、需有界延迟条件）；固定轮数 $+$ warm-start（实时 MPC 标准做法，每周期跑固定 $K$ 轮同步 ADMM $+$ warm-start，时间可预算、硬实时友好，但不保证收敛到最优、有小 gap）。

\begin{table}[H]\centering\small
\caption{三种 ADMM 调度方式对比}
\label{tab:dmpc-schedule}
\begin{tabular}{ccccc}
\toprule
方式 & 收敛保证 & 实时性 & 同步要求 & 适用 \\
\midrule
A 同步 & 凸有保证 & 一般 & 全局时钟 & 仿真/有同步 \\
B 异步 & 需条件（有界延迟） & 好 & 无 & 大规模/异构 \\
C 固定轮数$+$warm-start & 有 gap & 硬实时 & 周期内同步 & 真机多足 MPC（推荐） \\
\bottomrule
\end{tabular}
\end{table}

\subsection{\texorpdfstring{$\rho$}{rho} 调节与收敛性：理论与残差平衡}
\label{sec:dmpc-rho}

ADMM 的收敛速度对 $\rho$ 敏感。给出三条要点：

\textbf{(1) 凸情形的收敛率。} 对凸问题，ADMM 的目标值与约束违反以 $O(1/k)$ 收敛；若 $f$ 或 $g$ 之一强凸且光滑，可达 $O(1/k^2)$，甚至在适当条件下线性收敛。\textbf{这比对偶子梯度法的 $O(1/\sqrt k)$ 快得多}——这是 ADMM 优于纯对偶分解的关键定量优势\cite{boyd2011admm}。

\textbf{(2) 残差与停止准则。} 定义两个残差：\textbf{原始残差} $r^k=\mathbf{w}_i(\mathbf{X}_i^k)-\mathbf{z}_i^k$（一致性违反，衡量“本地副本与共识差多少”）；\textbf{对偶残差} $s^k=\rho(\mathbf{z}^k-\mathbf{z}^{k-1})$（共识变量的变化，衡量“还在不在动”）。当 $\|r^k\|$ 和 $\|s^k\|$ 都小于阈值时停止。$\rho$ 的作用直观：$\rho$ 大 $\to$ 重罚一致性违反 $\to r^k$ 降得快但 $s^k$ 可能大；$\rho$ 小则反之。

\textbf{(3) 残差平衡（residual balancing）自适应 $\rho$。} Boyd 综述\cite{boyd2011admm} 给的经典自适应规则——让两个残差量级相当：
\begin{equation}
  \rho^{k+1}=\begin{cases}\tau^{\text{incr}}\rho^k & \text{if } \|r^k\|>\nu\|s^k\|\ (\text{原始残差太大，加大 }\rho),\\[2pt] \rho^k/\tau^{\text{decr}} & \text{if } \|s^k\|>\nu\|r^k\|\ (\text{对偶残差太大，减小 }\rho),\\[2pt] \rho^k & \text{otherwise},\end{cases}
  \label{eq:dmpc-residual-balance}
\end{equation}
典型取 $\nu=10$，$\tau^{\text{incr}}=\tau^{\text{decr}}=2$。直觉：哪个残差大就调 $\rho$ 去压它，使原始/对偶残差同步下降。

\begin{insight}[非凸接触约束下 ADMM 无收敛保证]\label{ins:dmpc-nonconvex}
必须重申：上面的收敛率（$O(1/k)$ 等）建立在\textbf{凸性}上。多足的接触约束（单边接触 $f_z\ge0$ 的“开/关”、姿态在 $\SO(3)$ 上的非线性）使本地子问题非凸，连 ADMM 的\emph{收敛性}都无理论保证。工程上靠：固定接触序列后用摩擦锥金字塔\textbf{凸化}本地 QP $+$ warm-start $+$ 固定轮数“够用就停”，并接受有 gap。把非凸 ADMM 当“启发式但实践有效”的工具，不依赖它的理论收敛——这是本章贯穿始终的清醒认知（\cref{sec:dmpc-topo} 再展开）。
\end{insight}

\begin{pitfall}[编程陷阱：$\mathbf{z}$-更新漏掉对偶项 $\mathbf{u}_i$，只平均 $\mathbf{w}_i$]
$\mathbf{z}$-更新写成 \verb@average(w_i)@ 而非 \verb@average(w_i + u_i)@。\emph{后果}：破坏 ADMM 的对偶配平，收敛到错误的一致点或不收敛——残差降不到 $0$、编队有稳态偏差，且极难发现（“看起来在平均”）。\emph{根因}：consensus-ADMM 的 $\mathbf{z}$-更新（\cref{eq:dmpc-z-update}）是 $\mathbf{w}_i+\mathbf{u}_i$ 的平均，对偶偏移 $\mathbf{u}_i$ 是配平的关键。\emph{正确做法}：严格按 \cref{eq:dmpc-z-update} 把 $\mathbf{w}_i+\mathbf{u}_i$ 一起平均。这是 consensus-ADMM 最隐蔽的 bug。
\end{pitfall}

\begin{pitfall}[编程陷阱：每个控制周期冷启动（$\mathbf{u}$ 和 $\mathbf{z}$ 清零）]
每周期把 $\mathbf{z},\mathbf{u}$ 清零重新跑。\emph{后果}：固定少轮根本不够收敛，编队跟踪滞后、抖动。\emph{根因}：相邻控制周期的最优解高度相似，冷启动丢弃了这个先验，浪费迭代。\emph{正确做法}：warm-start——保留上周期的 $\mathbf{z},\mathbf{u}$ 和本地解作为初值。这是实时 ADMM-MPC 的命脉，没有它固定轮数方案不成立。
\end{pitfall}

\begin{pitfall}[概念误区：$\rho$ 越大收敛越快]
把 $\rho$ 调到很大期望加速。\emph{后果}：原始残差掉得快但对偶残差大、本地 QP 病态、机器人运动僵硬、解质量差。\emph{根因}：$\rho$ 过大让二次惩罚淹没本地代价 $J_i$，agent 只顾贴近共识而不顾自身最优；$\rho$ 还影响本地 QP 的条件数。\emph{正确做法}：$\rho$ 适中（如 $1$–$100$，问题相关）$+$ 残差平衡自适应。把 $\rho$ 理解为“全局一致 vs 本地自主”的权重旋钮（呼应 \cref{sec:cs-admm}），而非“速度旋钮”。
\end{pitfall}

\begin{pitfall}[概念误区：以为 ADMM 在非凸接触约束下也保证收敛]
把凸 ADMM 的收敛保证直接套到含单边接触/姿态非凸的多足 MPC。\emph{后果}：某些步态切换或大姿态变化时 ADMM 震荡或发散，却以为是代码 bug。\emph{根因}：ADMM 的收敛理论建立在凸性上；多足接触约束非凸，理论保证失效。\emph{正确做法}：凸化（固定接触序列后用摩擦锥金字塔）$+$ warm-start $+$ 固定轮数“够用就停”，并接受有 gap。
\end{pitfall}

\begin{pitfall}[编程陷阱：用单一固定 $\rho$ 且不监控残差]
设个 $\rho$ 就不管了，也不算 $r^k,s^k$。\emph{后果}：无法判断是否收敛、无法触发提前停、无法自适应调 $\rho$，调参全靠盲试。\emph{根因}：缺少残差这个“仪表盘”，ADMM 的运行状态不可观测。\emph{正确做法}：每轮计算原始/对偶残差，用于停止准则、自适应 $\rho$ 和调试（残差曲线一眼看出收敛健康度）。
\end{pitfall}

\begin{exercise}
\begin{enumerate}
  \item \textbf{（实现 $+$ 验证 consensus-ADMM）} 用 Python 实现两 agent 共识：$\min \tfrac12(x_1-1)^2+\tfrac12(x_2-3)^2$ s.t.\ $x_1=x_2$。写出三步（接口 $w_i=x_i$，共识 $z$，对偶 $u_i$），$\rho=1$ 跑 $30$ 轮，验证收敛到 $x_1=x_2=2$。然后\textbf{故意漏掉 $\mathbf{z}$-更新的 $u_i$ 项}，观察收敛到错误值——亲手制造并诊断这个最隐蔽的 bug。
  \item \textbf{（$\rho$ 扫描与残差平衡）} 对 $\rho\in\{0.1,1,10,100\}$ 各跑 ADMM，画原始残差 $r^k$ 和对偶残差 $s^k$ 曲线。观察小 $\rho$ 原始残差降得慢、大 $\rho$ 对偶残差大。再实现残差平衡自适应 $\rho$（$\nu=10,\tau=2$），对比固定 $\rho$ 的收敛轮数。
  \item \textbf{（双积分器 ADMM 编队，承接 \cref{sec:dmpc-why,sec:dmpc-formation}）} 把 $N=3$ 双积分器一字队形 MPC 用 consensus-ADMM 分布式求解：接口变量取相邻 agent 在编队边上的相对位移。实现“方式 C”（固定 $K=5$ 轮 $+$ warm-start）。对比集中式联合 MPC：(a) 队形跟踪误差；(b) 单周期计算时间（分布式取 max over agents）；(c) ADMM 轮数对 gap 的影响（$K=1,3,5,10$）。
  \item \textbf{（跨章综合题）} 结合 \cref{ch:mr-consensus}（收敛率由 $\lambda_2$ 决定、Laplacian）与本节（ADMM 的 $\mathbf{z}$-更新是共识平均）：在一字链 $1$-$2$-$3$-$4$-$5$ 和星形（中心连四个）下，分别估算 ADMM 的 $\mathbf{z}$-更新若用分布式共识迭代实现，需多少轮共识才能算准平均。说明为什么通信拓扑既影响 ADMM 外层迭代轮数、又影响每轮 $\mathbf{z}$-更新内层共识轮数？（预告 \cref{sec:dmpc-topo}）
\end{enumerate}
\end{exercise}

% ════════════════════════════════════════════════════════════════
\section{互联 SRB 网络与多足编队：CoM 协调 \texorpdfstring{$+$}{+} 接触约束 \texorpdfstring{$+$}{+} 耦合力}
\label{sec:dmpc-srb}

\paragraph{动机：把“会走路的机器人”装进分布式 MPC。}
回顾单体四足 MPC（前置）：一台四足被近似成单刚体（SRB），状态 $12$ 维 $[\mathbf{p},\boldsymbol{\Theta},\dot{\mathbf{p}},\boldsymbol{\omega}]$，控制是四条腿的接触力 $\mathbf{f}_{i1},\dots,\mathbf{f}_{i4}$，动力学是牛顿-欧拉方程，约束是摩擦锥 $+$ 单边接触 $+$ 踏步时序。现在把 $N$ 台这样的四足放进编队。最直接的想法是：每台机器人的本地 MPC 就是它的单体 SRB MPC，编队耦合（\cref{sec:dmpc-formation}）加在质心上，用 ADMM（\cref{sec:dmpc-admm}）协调。\textbf{对“编队只靠相对位置耦合”的场景（如两台机器人保持间距并排走、不共享负载），这确实够了}——本地 SRB MPC $+$ 质心 displacement 约束 $+$ ADMM，直接组装。

但一旦机器人\textbf{共享一个物理负载}（一起抬一根杆、一块板、一个箱子），事情变复杂：它们通过负载\emph{互相施力}。agent $i$ 推负载，负载就推 agent $j$——这个力既不是 agent 自己的接触力（来自地面），也不是重力，而是\textbf{通过刚性连接传递的约束力}。它必须进入每台机器人的 SRB 动力学，否则动力学方程从根上就是错的。这就是 SRB \textbf{网络}模型要解决的：不是 $N$ 个孤立的 SRB，而是通过约束力耦合的\textbf{互联 SRB 网络（interconnected SRB network）}——Kim--Fawcett--Hamed T-RO 2023\cite{kim2023layered} 的核心建模。

\paragraph{反面：如果忽略耦合力直接堆 $N$ 个单体 SRB 会怎样。}
假设建模时偷懒，把双四足抬杆当成两个独立的单体 SRB，各自动力学只含自己的接触力和重力，编队约束加在质心上：
\begin{enumerate}
  \item \textbf{动力学方程缺项，物理错误}。负载通过杆对每台机器人施加的力没出现在牛顿-欧拉方程里，MPC 预测的质心运动会偏离真实——预测错，控制就错。
  \item \textbf{负载的运动无人负责}。负载自身也是一个刚体，有自己的动力学（受两台机器人施的力 $+$ 重力）。不建模负载，就无法保证负载被托住、被搬到目标。
  \item \textbf{内力失控}。两台机器人对负载施的力如果不协调，会产生\textbf{内力}（互相对抗、挤压或拉扯负载）——轻则浪费能量，重则把负载弄坏或把机器人自己拽倒。这正是协同搬运章的核心，而它的根在本节的耦合力建模。
\end{enumerate}

\begin{insight}[SRB 网络 $=$ 子系统 $+$ 缝合约束]\label{ins:dmpc-srb-stitch}
互联 SRB 网络不是“$N$ 个单体 SRB 的简单并列”，\textbf{而是“$N$ 个单体 SRB $+$ 共享负载 $+$ 把它们绑在一起的刚性耦合约束（及其产生的约束力 $\boldsymbol{\lambda}$）”}。这个“子系统 $+$ 缝合约束”的结构，正是 \cref{sec:dmpc-dual,sec:dmpc-admm} 分布式优化最擅长处理的形式：子系统 $\to$ 本地子问题，缝合约束 $\to$ 要协调的耦合，约束力 $\boldsymbol{\lambda}\to$ 对偶变量/接口变量。所以建模选择 SRB 网络不是偶然——它是为了让“多足协同”恰好落进“分布式优化”的标准框架里。建模和求解在这里是共生设计的。
\end{insight}

\paragraph{历史：从单体 SRB 到互联 SRB 网络。}
单刚体（SRB）模型在足式机器人里的流行，源于 MIT Cheetah（Di Carlo, Wensing, Kim 等 $2018$\cite{dicarlo2018cheetah}）——他们发现对腿轻的四足，把全身近似成单刚体、用凸 MPC 优化接触力，既简单又准确，成为四足 MPC 的标准模板。SRB 之所以好，是因为它\emph{保留了接触力与质心运动的本质耦合，又避免了全身高维动力学的复杂}。把单体 SRB 推广到多机协同，关键一步是 Kim--Hamed 等人（ACC 2022 监督式预测控制\cite{kim2022cooperative} $\to$ T-RO 2023 分层控制\cite{kim2023layered}）提出的\textbf{互联 SRB 网络}：每个 agent 一个 SRB，通过共享负载的刚性约束（holonomic constraint，完整约束）耦合，约束产生拉格朗日乘子（约束力）。它让“多足协同”既保持单体 SRB 的简洁（每个 agent 仍是 $12$ 维 SRB），又显式刻画协同的物理本质（约束力耦合）——从而可直接套用 \cref{sec:dmpc-admm} 的 ADMM 分解（约束力就是要在 agent 间达成一致的接口变量）。

\subsection{互联 SRB 网络的动力学}
\label{sec:dmpc-srb-dyn}

先单个 agent，再耦合。

\textbf{(1) 单 agent SRB 动力学。} agent $i$ 质心位置 $\mathbf{p}_i\in\mathbb{R}^3$、姿态 $\mathbf{R}_i\in\SO(3)$（或欧拉角 $\boldsymbol{\Theta}_i$）、质量 $m_i$、机体惯量 $\mathbf{I}_i$。它受：四条腿的地面接触力 $\mathbf{f}_{ij}$（$j=1\dots4$，作用点 $\mathbf{r}_{ij}$ 相对质心）、负载施加的耦合力 $\boldsymbol{\lambda}_i$（作用点为附着点 $\mathbf{p}_i^{\text{att}}$）、重力 $\mathbf{g}$。牛顿-欧拉方程：
\begin{align}
  m_i\ddot{\mathbf{p}}_i &= \sum_{j=1}^{4}\mathbf{f}_{ij} + \boldsymbol{\lambda}_i + m_i\mathbf{g}, \label{eq:dmpc-srb-newton}\\
  \tfrac{d}{dt}(\mathbf{I}_i\boldsymbol{\omega}_i) &= \sum_{j=1}^{4}\mathbf{r}_{ij}\times\mathbf{f}_{ij} + (\mathbf{p}_i^{\text{att}}-\mathbf{p}_i)\times\boldsymbol{\lambda}_i. \label{eq:dmpc-srb-euler}
\end{align}
与单体 SRB 唯一的区别就是多了 $\boldsymbol{\lambda}_i$ 项（耦合力及其力矩）。\textbf{这一项就是协同的全部物理内容。}

\textbf{(2) 共享负载动力学。} 负载（质量 $m_L$、惯量 $\mathbf{I}_L$、质心 $\mathbf{p}_L$）受各 agent 通过附着点施加的反作用力 $-\boldsymbol{\lambda}_i$（牛顿第三定律）$+$ 重力：
\begin{equation}
  m_L\ddot{\mathbf{p}}_L = -\sum_{i=1}^{N}\boldsymbol{\lambda}_i + m_L\mathbf{g},\qquad \tfrac{d}{dt}(\mathbf{I}_L\boldsymbol{\omega}_L)=-\sum_i(\mathbf{p}_i^{\text{att}}-\mathbf{p}_L)\times\boldsymbol{\lambda}_i.
  \label{eq:dmpc-load-dyn}
\end{equation}

\textbf{(3) 刚性耦合约束（holonomic constraint）。} “刚性连接”意味着 agent $i$ 的附着点在负载坐标系里的位置固定。等价地，附着点的绝对位置由负载位姿决定：
\begin{equation}
  \mathbf{p}_i^{\text{att}} = \mathbf{p}_L + \mathbf{R}_L\,\mathbf{a}_i \quad(\mathbf{a}_i=\text{附着点在负载系的固定坐标}),
  \label{eq:dmpc-holonomic}
\end{equation}
或写成两个 agent 间的相对形式（如 Kim--Fawcett--Hamed\cite{kim2023layered} 的 holonomic constraint）：附着点相对位置固定 $\mathbf{p}_i^{\text{att}}-\mathbf{p}_j^{\text{att}}=\mathbf{d}_{ij}^{\text{fixed}}$。

\textbf{(4) 约束力 $\boldsymbol{\lambda}$ 从哪来——把约束微分两次。} 这是初学者最困惑的地方：$\boldsymbol{\lambda}_i$ 不是独立给定的，它是\textbf{维持约束所需的力}，由约束方程解出。机理（拉格朗日乘子法的标准流程）：把位置层约束 $\boldsymbol{\Phi}(\mathbf{q})=\mathbf{0}$（$\mathbf{q}$ 为广义坐标）对时间求两次导，得到加速度层约束：
\begin{equation}
  \boldsymbol{\Phi}(\mathbf{q})=\mathbf{0} \xrightarrow{\frac{d}{dt}} \mathbf{J}_\Phi\dot{\mathbf{q}}=\mathbf{0} \xrightarrow{\frac{d}{dt}} \mathbf{J}_\Phi\ddot{\mathbf{q}}+\dot{\mathbf{J}}_\Phi\dot{\mathbf{q}}=\mathbf{0},
  \label{eq:dmpc-constraint-diff}
\end{equation}
其中 $\mathbf{J}_\Phi=\partial\boldsymbol{\Phi}/\partial\mathbf{q}$ 是约束 Jacobian。约束力沿约束法向，写成 $\mathbf{J}_\Phi^\top\boldsymbol{\lambda}$。把它代入各刚体的动力学 $\mathbf{M}\ddot{\mathbf{q}}=\mathbf{Q}+\mathbf{J}_\Phi^\top\boldsymbol{\lambda}$，与加速度层约束联立，解出：
\begin{equation}
  \boxed{\ \boldsymbol{\lambda}=-(\mathbf{J}_\Phi \mathbf{M}^{-1}\mathbf{J}_\Phi^\top)^{-1}\big(\mathbf{J}_\Phi \mathbf{M}^{-1}\mathbf{Q}+\dot{\mathbf{J}}_\Phi\dot{\mathbf{q}}\big)\ }
  \label{eq:dmpc-constraint-force}
\end{equation}
这就是约束力的解析表达式。\textbf{物理含义}：$\boldsymbol{\lambda}$ 恰好是让所有刚体的加速度满足“附着点保持刚性连接”所需的那个力。它正是 \cref{sec:dmpc-dual} 预告的“影子价格”在物理上的化身——维持耦合约束的拉格朗日乘子，等于约束力。\pz{\cref{eq:dmpc-constraint-force} 为约束动力学（约束微分两次 $+$ 拉格朗日乘子）的标准结果，源稿转录；其中 $\mathbf{M},\mathbf{Q},\mathbf{J}_\Phi$ 的具体堆叠（含 Baumgarte 镇定项与否）随实现而异，未逐项核。}

\begin{note}
\textbf{阶段小结.}\;
到这里我们建立了完整的 SRB 网络：(1) 每个 agent 的 SRB 动力学（多了 $\boldsymbol{\lambda}_i$ 项）；(2) 负载动力学；(3) 刚性耦合约束；(4) 约束力 $\boldsymbol{\lambda}$ 的解析解。接下来把它装进分布式 MPC——每个 agent 的本地 MPC 含自己的 SRB $+$ 接触约束，$\boldsymbol{\lambda}_i$ 是要在 agent 间协调的接口变量。
\end{note}

\subsection{接触约束：摩擦锥、单边接触、踏步时序}
\label{sec:dmpc-contact}

多足 MPC 的本地约束（每个 agent 自己的 $\mathcal{C}_i$）比双积分器复杂得多，三类必须进 MPC：

\textbf{(1) 摩擦锥（friction cone）。} 接触力不能让脚打滑：切向力 $\le$ 法向力 $\times$ 摩擦系数 $\mu$。精确形式是二阶锥 $\sqrt{f_x^2+f_y^2}\le\mu f_z$（SOCP，慢）。MPC 里几乎都用\textbf{金字塔近似}（$4$ 个线性不等式，QP 友好）：
\begin{equation}
  -\mu f_z\le f_x\le\mu f_z,\quad -\mu f_z\le f_y\le\mu f_z \quad(\text{或更精确的 }\sqrt2\text{ 缩放}).
  \label{eq:dmpc-pyramid}
\end{equation}
金字塔近似把 SOCP 降为 QP，让 OSQP 这类 QP 求解器能在毫秒内解完——这是多足 MPC 实时性的前提。

\textbf{(2) 单边接触（unilateral contact）。} 脚只能“推”地不能“拉”地：支撑腿法向力 $f_z\ge0$；摆动腿（不接触）力为零 $\mathbf{f}_{ij}=\mathbf{0}$。

\textbf{(3) 踏步时序（contact schedule / gait）。} 哪条腿在哪个时刻支撑、哪条摆动，由步态调度器给定（\cref{sec:dmpc-layered} 中层）。MPC 在预测时域上\textbf{按给定的接触序列}施加上面的约束——支撑相用摩擦锥 $+$ $f_z\ge0$，摆动相强制 $\mathbf{f}=\mathbf{0}$。\textbf{关键}：固定接触序列后，本地 MPC 是凸 QP（这是 \cref{sec:dmpc-admm} 凸化的来源）；如果让 MPC 自己决定接触序列（接触时机优化），问题变成混合整数/非凸，难解得多。把这些装进 agent $i$ 的本地 MPC（即 \cref{sec:dmpc-admm} $\mathbf{X}$-更新里的 $\min_{\mathbf{X}_i\in\mathcal{C}_i}$），结构如 \cref{alg:dmpc-local-qp}。

\begin{algo}[agent $i$ 本地 MPC（凸 QP，固定接触序列）]\label{alg:dmpc-local-qp}
$\min_{\mathbf{X}_i,\mathbf{f}_i,\boldsymbol{\lambda}_i}\ \sum_t\big[\|\mathbf{x}_i(t)-\mathbf{x}_i^{\text{ref}}(t)\|_{\mathbf{Q}}^2+\|\mathbf{f}_i(t)\|_{\mathbf{R}}^2\big]+\tfrac{\rho}{2}\|\boldsymbol{\lambda}_i-\mathbf{z}_\lambda+\mathbf{u}_i\|^2$（末项为 ADMM 惩罚，$\boldsymbol{\lambda}_i$ 是接口变量）；s.t.\ SRB 动力学 \cref{eq:dmpc-srb-newton,eq:dmpc-srb-euler}、摩擦锥金字塔 \cref{eq:dmpc-pyramid}（支撑腿）、单边接触（$f_{ij,z}\ge0$ 支撑、$\mathbf{f}_{ij}=\mathbf{0}$ 摆动）、按 $\text{contact\_schedule}(t)$ 施加。
\end{algo}

\noindent 注意 ADMM 惩罚项作用在\textbf{耦合力 $\boldsymbol{\lambda}_i$} 上——这就是多足编队里的“接口变量”：两台机器人对负载施的力必须协调一致（满足负载动力学与刚性约束），$\boldsymbol{\lambda}$ 就是要在 ADMM 协调步达成一致的量。这把 \cref{sec:dmpc-admm} 抽象的 $\mathbf{w}_i$ 具体化为物理的约束力。

\subsection{CoM 协调：踏步起伏与编队的频段分离}
\label{sec:dmpc-com-freq}

\cref{sec:dmpc-formation} 末尾的多足特有点在这里展开。SRB 网络的状态用质心，编队/协同约束作用在质心层面。但四足质心随踏步周期性起伏摆动——如果硬性协调瞬时质心，会和踏步动力学打架。SRB 网络如何处理？\textbf{答案是分层 $+$ 降阶天然实现的频段分离}：SRB MPC 工作在\emph{降阶模型}层面，优化的是质心的\emph{低频趋势}（往哪走、整体协调），预测时域 $0.3$–$0.5$ 秒覆盖几个踏步周期，自然平滑掉单步细节；高频的踏步起伏、足端轨迹由\emph{底层 WBC}（\cref{sec:dmpc-layered}）处理；耦合力 $\boldsymbol{\lambda}$ 在 SRB 层协调——两台机器人保证对负载的\emph{净力/净力矩}协调（托住负载、搬向目标），而各自如何踏步是本地自由。

\begin{insight}[把协同放降阶层、把平衡放全阶层 $=$ 按频率分解控制任务]\label{ins:dmpc-freq-decompose}
这就是为什么 Kim--Fawcett--Hamed\cite{kim2023layered} 用“SRB 网络（高层 MPC）$+$ 全阶非线性控制器（低层）”的分层——SRB 网络在低频协调质心与耦合力（群体行为），全阶控制器在高频跟踪足端与平衡（个体行为）。这个频段分离不是工程技巧，而是多足协同得以分布式实现的\textbf{结构性原因}——慢的群体行为信息量小、可低频通信协调（ADMM 几轮够用），快的个体行为留本地（WBC 无需通信）。
\end{insight}

\begin{insight}[SRB 网络像一队人合抬大桌子走路（标边界）]\label{ins:dmpc-table}
SRB 网络协调多足，像\textbf{一队人合抬一张大桌子走路}。每个人主要管自己走稳、迈步（本地 SRB $+$ 踏步 $=$ 个体高频平衡）；同时大家要让桌子保持水平、一起朝门口移动（耦合力 $\boldsymbol{\lambda}$ 协调 $=$ 群体低频协同）。\emph{像}的地方：个体动作（迈步）和群体目标（搬桌子）在不同层面，各人迈步节奏不必一致，只要对桌子的合力协调。\emph{不像}的地方：人靠手感和默契协调桌子受力，SRB 网络靠 ADMM 严格优化耦合力 $\boldsymbol{\lambda}$ 并保证满足负载动力学；且人能实时感知桌子晃动即时调整，SRB-MPC 是预测式的。\emph{边界}：这个类比对“刚性连接”成立，若是绳索（柔性、只能拉不能推），耦合约束就不是刚性 holonomic 而是单边的，建模要变（见前沿 PACC 被动臂方案\cite{turrisi2024pacc}）。
\end{insight}

\subsection{质心动力学的离散化与 MPC 预测方程}
\label{sec:dmpc-discretize}

上面给的 SRB 网络是连续时间的微分方程。MPC 是离散优化，必须把它\textbf{离散化}成预测方程 $\mathbf{x}_i(t{+}1)=\mathbf{A}_i\mathbf{x}_i(t)+\mathbf{B}_i\mathbf{u}_i(t)+\mathbf{B}_i^\lambda\boldsymbol{\lambda}_i(t)+\mathbf{c}_i$。这一步有几个多足特有的微妙处：

\textbf{(1) 姿态的处理：欧拉角线性化 vs $\SO(3)$。} SRB 的姿态在旋转群 $\SO(3)$ 上，是非线性流形，两种处理：\emph{欧拉角 $+$ 小角度线性化}（MIT Cheetah 凸 MPC\cite{dicarlo2018cheetah}）假设机体姿态变化不大，把旋转动力学在当前姿态附近线性化、姿态用欧拉角表示，得到\textbf{线性时变（LTV）}预测方程、问题是凸 QP，缺点是大姿态变化时线性化误差大；\emph{$\SO(3)$ 上的几何积分}（De Vincenti--Coros RSS 2023\cite{devincenti2023centralized} 用四元数 $+$ 李群积分）在流形上精确积分、不做小角度近似，大姿态精确但问题变非凸（SQP）、慢。

\begin{insight}[凸化是贯穿始终的主题]\label{ins:dmpc-convexify}
为什么多足编队 MPC 大多选欧拉角线性化而非 $\SO(3)$ 精确？因为编队/搬运任务中机体姿态通常变化平缓（平地搬运、保持负载水平），小角度假设成立，换来凸 QP 的实时性和 ADMM 的可分解性。$\SO(3)$ 精确积分留给大机动场景。这又是一个“精度 vs 实时/可解”的工程权衡——和摩擦锥金字塔、固定接触序列、长时域替代终端集（\cref{sec:dmpc-topo}）是同一类决策：在多足分布式 MPC 里，\textbf{凸化是贯穿始终的主题，因为只有凸 QP 才能在毫秒级 $+$ ADMM 分解下跑起来}。（与全书右扰动 $\SO(3)$ 约定的关系：源稿在此用欧拉角小角度线性化，属降阶 MPC 的常见近似；本书右扰动几何处理见 \nameref{ch:lie}，二者在“大姿态精确”需求下应改用 $\SO(3)$ 几何积分。）
\end{insight}

\textbf{(2) 接触力作为控制输入的特殊性。} 多足 MPC 的控制 $\mathbf{u}_i=[\mathbf{f}_{i1},\dots,\mathbf{f}_{i4}]$ 是接触力，不是常见的“加速度/力矩指令”。这带来两个特点：控制直接进入动力学的\emph{右端}（$m_i\ddot{\mathbf{p}}_i=\sum_j\mathbf{f}_{ij}+\dots$），$\mathbf{B}_i$ 矩阵结构简单（力的求和 $+$ 力矩的叉乘）；接触力受\emph{单边 $+$ 摩擦锥}约束（只能在锥内），不是无约束的——这正是 MPC 必须显式建约束的原因，也是它区别于“无约束 LQR”的地方。

\textbf{(3) 离散化方法的选择。} 把连续 $\dot{\mathbf{x}}=f(\mathbf{x},\mathbf{u})$ 离散，常用：前向欧拉 $\mathbf{x}(t{+}1)=\mathbf{x}(t)+\Delta t\cdot f$（最简单，大 $\Delta t$ 不稳定；多足 MPC 常用，$\Delta t$ 小如 $0.02$\,s）、零阶保持 ZOH（对 LTI 精确离散 $\mathbf{A}_d=e^{\mathbf{A}\Delta t}$，线性化后用、更精确）、半隐式/RK（更高阶、更准更重，非线性 MPC 用）。

\textbf{(4) 完整的离散预测方程（单 agent，线性化后）。} 综合上述，agent $i$ 在固定接触序列下的离散 SRB 预测方程是 LTV 形式：
\begin{equation}
  \mathbf{x}_i(t{+}1)=\mathbf{A}_i(t)\mathbf{x}_i(t)+\sum_{j\in\text{支撑腿}(t)}\mathbf{B}_{ij}(t)\mathbf{f}_{ij}(t)+\mathbf{B}_i^\lambda(t)\boldsymbol{\lambda}_i(t)+\mathbf{g}_i,
  \label{eq:dmpc-ltv}
\end{equation}
其中 $\mathbf{A}_i(t)$ 含线性化的旋转动力学，$\mathbf{B}_{ij}(t)$ 依赖当前足端相对质心位置 $\mathbf{r}_{ij}(t)$（决定力矩臂），$\mathbf{g}_i$ 含重力。注意 $\mathbf{B}_{ij}$ 随踏步\emph{时变}（哪条腿支撑、足端在哪都在变），这是“线性时变”而非“线性时不变”的原因。把这个方程在时域 $T$ 上展开，就是 \cref{alg:dmpc-local-qp} 本地 MPC 的动力学约束。

\subsection{一个具体数值例子：双足抬杆的耦合力}
\label{sec:dmpc-srb-example}

为了让 $\boldsymbol{\lambda}$ 不停留在抽象，看一个简化数值例子（一维竖直方向）。两台机器人各在杆的一端，杆水平、质量 $m_L=5$\,kg、重力 $g=9.8$。杆刚性，要求两端竖直位置满足约束（杆不变形、不旋转，两端等高加速度）：$\ddot z_1^{\text{att}}=\ddot z_2^{\text{att}}=\ddot z_L$。设两台机器人通过腿对地的净竖直力分别为 $F_1,F_2$（它们能控制的量），杆给机器人的竖直耦合力为 $\lambda_1,\lambda_2$（这里取 $\lambda_i$ 为机器人托杆的力，作用在杆上）：
\begin{itemize}
  \item 杆动力学：$m_L\ddot z_L=\lambda_1+\lambda_2-m_L g$（两端托力 $-$ 重力）。
  \item 若要杆匀速（$\ddot z_L=0$）：$\lambda_1+\lambda_2=m_L g=49$\,N。这是\textbf{运动力约束}——两端托力之和必须等于杆重才能托住。
  \item 但 $\lambda_1,\lambda_2$ 的\textbf{分配}有一个自由度：$\lambda_1=49/2+\delta,\ \lambda_2=49/2-\delta$。$\delta$ 就是\textbf{内力}——$\delta\ne0$ 意味着一端多托、一端少托，杆受一个净力矩（若杆能转）或内部剪力。
\end{itemize}

\begin{insight}[耦合力 $=$ 运动力（和，唯一确定）$+$ 内力（差，必须受控）]\label{ins:dmpc-motion-internal}
这个例子赤裸裸地展示了耦合力 $\boldsymbol{\lambda}$ 分两部分：\textbf{和（$\lambda_1+\lambda_2$）由负载动力学唯一确定（运动力，托住杆）}，\textbf{差（$\delta$）是自由的（内力，必须由协调控制约束住）}。ADMM 协调 $\boldsymbol{\lambda}$ 时，运动力部分被负载动力学钉死（两端必须合力托杆），内力部分 $\delta$ 必须被显式控制到 $0$ 或安全范围（否则两台机器人一个使劲往上托、一个往下压，把杆内部应力拉爆，或互相把对方拽失衡）。这就是为什么“协调耦合力”不只是“让合力对”，还要“让内力受控”——前者是本章 \cref{sec:dmpc-srb}，后者通向协同搬运章的 grasp matrix。这个 $49$\,N、$\delta$ 的小例子，就是源稿所述 Kim--Fawcett--Hamed 双 A1 搬 $5$\,kg 负载在做的事情的一维投影。\pz{“双 A1 搬 $5$\,kg”为源稿给出的实验设定，未核。}
\end{insight}

\begin{pitfall}[概念误区：把耦合力 $\boldsymbol{\lambda}$ 当成可独立优化的控制输入]
在 MPC 里把 $\boldsymbol{\lambda}_i$ 当成和接触力 $\mathbf{f}_{ij}$ 一样的自由决策变量随便优化。\emph{后果}：解出的 $\boldsymbol{\lambda}$ 不满足刚性约束（附着点位置漂移），双足对负载施力不自洽，仿真里负载乱飞或机器人被拽飞。\emph{根因}：$\boldsymbol{\lambda}$ 是约束力，由刚性约束 $+$ 动力学\emph{解出}（\cref{eq:dmpc-constraint-force}），不是独立自由量；它在 agent 间必须满足负载动力学与第三定律的一致性。\emph{正确做法}：把 $\boldsymbol{\lambda}_i$ 作为\textbf{接口/对偶变量}，通过 ADMM 协调使其在 agent 间一致且满足负载动力学；或显式用约束力公式代入。它的自由度受刚性约束严格限制。
\end{pitfall}

\begin{pitfall}[编程陷阱：摩擦锥用精确二阶锥导致 MPC 解不动]
本地 MPC 里直接写 $\sqrt{f_x^2+f_y^2}\le\mu f_z$。\emph{后果}：问题变 SOCP，QP 求解器（OSQP）用不了，换 SOCP 求解器慢一个数量级，毫秒级控制循环超时。\emph{根因}：二阶锥是非线性约束，破坏了 QP 结构。\emph{正确做法}：用金字塔线性近似（\cref{eq:dmpc-pyramid}，$4$ 个线性不等式），保持 QP，OSQP 毫秒内解完。代价是可行域略保守（金字塔内切于锥），实践可接受。
\end{pitfall}

\begin{pitfall}[概念误区：忽略摆动腿力为零的约束]
所有腿都加摩擦锥 $+$ $f_z\ge0$，不区分支撑/摆动。\emph{后果}：MPC 可能给摆动腿（悬空、不接触地面）分配非零力——物理上不可能，导致控制指令错误、机器人动作诡异。\emph{根因}：摆动腿不接触地面，接触力必须为零；只对支撑腿施加摩擦锥。\emph{正确做法}：按踏步时序 $\text{contact\_schedule}(t)$，摆动相强制 $\mathbf{f}_{ij}(t)=\mathbf{0}$，支撑相施加摩擦锥 $+$ 单边约束。
\end{pitfall}

\begin{pitfall}[思维陷阱：把编队约束加在足端而非质心]
为了“精确”，把编队/协同约束加在脚的位置上。\emph{后果}：足端是高频运动（每步抬落），把编队约束加在足端会和踏步剧烈冲突，机器人无法正常迈步。\emph{根因}：违反频段分离——编队是低频群体行为，应在质心（降阶）层；足端是高频个体行为，属底层 WBC。\emph{正确做法}：编队/协同约束加在质心（SRB 层），足端细节交给底层。
\end{pitfall}

\begin{pitfall}[概念误区：用刚性 holonomic 约束建模绳索/柔性连接]
不管连接是刚性杆还是绳索，一律用 $\mathbf{p}_i^{\text{att}}-\mathbf{p}_j^{\text{att}}=\text{const}$。\emph{后果}：绳索只能拉不能推（单边）、还能松弛，用刚性等式约束会建出错误的力（如绳索“推”负载）。\emph{根因}：刚性连接是双边等式约束（杆能拉能推），绳索是单边不等式约束（只能拉、可松弛）。\emph{正确做法}：刚性杆用 holonomic 等式约束（本节）；绳索用单边约束（$\boldsymbol{\lambda}$ 沿绳方向且 $\ge0$、长度 $\le$ 绳长）；被动臂方案（PACC\cite{turrisi2024pacc}）则用另一套建模。连接类型决定约束类型。
\end{pitfall}

\begin{exercise}
\begin{enumerate}
  \item \textbf{（推导约束力）} 最简一维：两个质量 $m_1=m_2=1$ 的滑块由刚性杆连接，约束 $x_1-x_2=\ell$，各自受外力 $F_1,F_2$。(a) 写出约束 Jacobian $\mathbf{J}_\Phi$；(b) 把约束微分两次得加速度层约束；(c) 用 \cref{eq:dmpc-constraint-force} 解出杆中力 $\lambda$；(d) 验证 $\lambda=(F_2-F_1)/2$ 的物理意义（杆传递两端外力差的一半）。
  \item \textbf{（SRB 网络动力学搭建）} 为“双 Go2 $+$ 刚性杆”写出完整的 SRB 网络动力学方程组（两个 agent SRB $+$ 杆作为负载 $+$ 刚性附着约束）。明确：状态向量维度（两个 $12$ 维 SRB $+$ 杆 $6$ 维 $=30$ 维）、每个方程里 $\boldsymbol{\lambda}$ 出现的位置、约束 Jacobian 的结构。
  \item \textbf{（接触约束实现）} 用 cvxpy/OSQP 实现单台四足的一步 SRB MPC（含摩擦锥金字塔 $+$ 单边接触 $+$ 给定 trot 踏步时序），预测 $10$ 步。给定期望质心速度求接触力。然后\textbf{故意把摩擦锥换成精确二阶锥}（SOCP），对比求解时间——亲手验证金字塔 vs 二阶锥的实时性差异。
  \item \textbf{（跨章综合题）} 结合协同搬运章预告（grasp matrix 与内力）$+$ 本节耦合力 $\boldsymbol{\lambda}$：双四足抬杆时，用 $\boldsymbol{\lambda}_1,\boldsymbol{\lambda}_2$ 表达：哪部分组合是运动力（搬动杆）、哪部分是内力（挤压/拉扯杆）？为什么 ADMM 协调 $\boldsymbol{\lambda}$ 时必须保证内力受控？
\end{enumerate}
\end{exercise}

% ════════════════════════════════════════════════════════════════
\section{分层控制：从 MPC 命令到关节力矩}
\label{sec:dmpc-layered}

\paragraph{动机：MPC 给的是“宏观命令”，电机要的是“微观力矩”。}
\cref{sec:dmpc-srb} 的分布式 MPC 输出的是 SRB 级命令——质心运动目标、四条腿的接触力 $\mathbf{f}_{ij}$、耦合力 $\boldsymbol{\lambda}_i$。但这些都是“宏观”的——它们描述的是“质心该怎么动、脚该使多大力”，而不是“髋关节、膝关节该输出多少力矩”。电机只能控制关节力矩。从“足端接触力”到“关节力矩”，要经过腿的运动学/动力学映射（雅可比转置 $\boldsymbol{\tau}=\mathbf{J}^\top\mathbf{f}$ 等）。而且 SRB 是个近似——它假设腿无质量、机器人是单刚体，真实机器人有 $12+$ 个关节、腿有质量、姿态在 $\SO(3)$ 上。需要一个\textbf{全阶（full-order）}控制器把 SRB 级命令“翻译”成满足全身动力学的关节力矩，同时处理 SRB 忽略的细节。这个全阶控制器就是 \textbf{WBC（Whole-Body Control，全身控制）}。分层的理由（呼应 \cref{sec:dmpc-com-freq} 的频段分离）：MPC 算得“慢但远”（预测 $0.4$ 秒、计算几毫秒、可较低频如 $100$–$500$\,Hz），WBC 算得“快但近”（只看当前、计算极快、高频如 $1$\,kHz）。两层各司其职——MPC 管协调与规划（含通信），WBC 管执行与平衡（纯本地）。

\paragraph{反面：如果不分层、让一个控制器全包会怎样。}
假设写一个“全阶分布式 MPC”——直接在关节力矩层面、用全身动力学、跨预测时域、还带编队耦合，一把求解。会怎样？\emph{维度爆炸 $\times$ 非线性 $\times$ 耦合，三重暴击}（全身动力学每台机器人 $18+$ 维，$N$ 台 $+$ 负载 $+$ 预测时域 $+$ 编队耦合，规模和非线性远超实时可解）；\emph{通信负担巨大}（全阶量都要在 agent 间协调，关节级轨迹比质心级大一个数量级）；\emph{数值条件差}（全身动力学含快慢混合的动态，统一在一个优化里数值病态）。

\begin{insight}[分层 $=$ 降阶 MPC 管低频协调（含通信）$+$ 全阶 WBC 管高频执行（纯本地）]\label{ins:dmpc-layered-not}
多足协同控制不是“一个无所不包的大优化”，\textbf{而是“分层——降阶 MPC 管低频协调（含通信），全阶 WBC 管高频执行（纯本地）”}。分层不是为了偷懒，而是因为“协调”和“执行”在频率、维度、通信需求上根本不同，强行合并会让两边都做不好。
\end{insight}

\paragraph{历史：从 Whole-Body Control 到分层 loco-manipulation。}
WBC 的思想源于 Khatib 的操作空间控制（operational space control，$1987$\cite{khatib1987operational}）——用任务优先级和零空间投影协调多个任务。足式机器人里，WBC 发展为基于 QP 的全身控制（Feng, Kuindersma 等 $2010$s）：在一个 QP 里同时满足全身动力学、接触约束、跟踪 MPC 给的质心/足端目标，输出关节力矩。把 MPC $+$ WBC 组成分层架构成为现代四足控制的标准（MIT Cheetah、ETH ANYmal）。Kim--Fawcett--Hamed T-RO 2023\cite{kim2023layered} 把它推广到多机协同：高层\textbf{分布式} MPC（SRB 网络 $+$ ADMM）$+$ 低层\textbf{分布式}非线性控制器。“分布式 WBC”的关键洞察是：\textbf{WBC 可以纯本地、无通信}——因为协调已经在 MPC 层通过耦合力 $\boldsymbol{\lambda}$ 完成了，WBC 只需让本机跟踪 MPC 给的（已协调好的）命令。\pz{Feng/Kuindersma 的 QP-WBC、Khatib 1987 操作空间控制归属系源稿转述；Khatib 1987 键已在库。}

\subsection{三层架构的数据流}
\label{sec:dmpc-dataflow}

完整的分层数据流（每台机器人 $i$ 本地运行，仅高层 MPC 需通信）见 \cref{fig:dmpc-hierarchy}。

\begin{figure}[ht]
\centering
\begin{tikzpicture}[
  >=Stealth, font=\small,
  lay/.style={draw, rounded corners, align=left, inner sep=5pt, text width=0.82\linewidth},
  hi/.style={lay, fill=main!8, draw=main!60!black},
  mid/.style={lay, fill=second!8, draw=second!60!black},
  lo/.style={lay, fill=third!8, draw=third!60!black},
  ar/.style={->, very thick, gray!55!black}]
\node[hi] (high) {\textbf{高层：分布式 MPC（SRB 网络 $+$ ADMM）\quad $\sim100$–$500$\,Hz}\\
  输入：自身状态 $\mathbf{x}_i$ $+$ 邻居预测（通信）$+$ 编队/任务参考\\
  计算：consensus-ADMM（\cref{sec:dmpc-admm}）$+$ 本地 SRB MPC（\cref{sec:dmpc-srb}）\\
  输出：质心轨迹参考 $+$ 足端接触力 $\mathbf{f}_{ij}$ $+$ 耦合力 $\boldsymbol{\lambda}_i$\\
  $\updownarrow$ 通信：仅交换接口变量（邻居预测/耦合力一致性，$\sim$KB 级）};
\node[mid, below=8mm of high] (midl) {\textbf{中层：步态调度（gait scheduler）\quad $=$ 控制频率}\\
  输入：步态类型（trot/walk）、相位\\
  输出：$\text{contact\_schedule}(t)$（哪条腿何时支撑/摆动）$+$ 摆动腿足端轨迹（落脚点）\\
  （通常本地、可预设；多机可共享步态相位以同步）};
\node[lo, below=8mm of midl] (low) {\textbf{低层：分布式 WBC（全身 QP）\quad $\sim1$\,kHz}\\
  输入：MPC 的足端力 $\mathbf{f}_{ij}$ $+$ 耦合力 $\boldsymbol{\lambda}_i$ $+$ 步态 $+$ 全身状态\\
  计算：本地全身 QP（满足全身动力学 $+$ 接触约束）\\
  输出：关节力矩 $\boldsymbol{\tau}_i\to$ 电机\\
  \textbf{纯本地、无通信——协调已由 MPC 的 $\boldsymbol{\lambda}_i$ 完成}};
\draw[ar] (high) -- node[right, font=\scriptsize] {SRB 级命令（已跨 agent 协调）} (midl);
\draw[ar] (midl) -- node[right, font=\scriptsize] {接触序列 $+$ 足端目标} (low);
\end{tikzpicture}
\caption{多足分布式 MPC 的三层控制数据流。通信只发生在高层 MPC、且只传 KB 级接口变量；中层步态调度通常本地、相位可选共享；低层 WBC 纯本地、可 $1$\,kHz 高频。仿 Kim--Fawcett--Hamed\cite{kim2023layered} 的分层架构重绘。}
\label{fig:dmpc-hierarchy}
\end{figure}

\textbf{低层分布式 WBC 的本地 QP}（每台机器人独立解，无通信）：
\begin{align}
  \min_{\boldsymbol{\tau}_i,\,\mathbf{f}_i,\,\ddot{\mathbf{q}}_i}\ & \big\|\mathbf{A}_i\ddot{\mathbf{q}}_i - \ddot{\mathbf{q}}_i^{\text{des}}\big\|^2 + w_f\big\|\mathbf{f}_i-\mathbf{f}_i^{\text{MPC}}\big\|^2 + w_\tau\|\boldsymbol{\tau}_i\|^2 \label{eq:dmpc-wbc-obj}\\
  \text{s.t.}\quad & \mathbf{M}_i\ddot{\mathbf{q}}_i+\mathbf{h}_i = \mathbf{S}_i^\top\boldsymbol{\tau}_i + \mathbf{J}_{c,i}^\top\mathbf{f}_i + \mathbf{J}_{\text{att},i}^\top\boldsymbol{\lambda}_i^{\text{MPC}}\quad(\text{全身动力学}), \label{eq:dmpc-wbc-dyn}\\
  & \text{friction\_cone}(\mathbf{f}_i),\quad \text{joint\_limits}(\boldsymbol{\tau}_i),\quad \text{contact constraints}. \notag
\end{align}
\textbf{关键}：耦合力 $\boldsymbol{\lambda}_i^{\text{MPC}}$ 由高层 MPC 的 ADMM 协调步给出（已和邻居一致），WBC 直接拿来用作全身动力学里的外力项 $\mathbf{J}_{\text{att},i}^\top\boldsymbol{\lambda}_i^{\text{MPC}}$，\emph{不需要再和邻居通信}。这就是“分布式 WBC 纯本地计算”的全部秘密——协调的“账”在 MPC 层已经算清，WBC 只管本机执行。

\begin{insight}[分层本质是一套“通信预算最小化”的设计]\label{ins:dmpc-comm-budget}
分层架构把“多机协调”这个需要通信的难题，\textbf{完全限制在了高层 MPC 这一层}——而且只在 SRB（降阶、低频、小数据量）层协调。一旦 MPC 通过耦合力 $\boldsymbol{\lambda}$ 达成一致，下游的步态和 WBC 就退化成 $N$ 个完全独立的本地问题，零通信、可任意高频。这是分层最深刻的价值：\textbf{它把“必须通信的部分”压缩到最小（低频、降阶、KB 级），把“可以本地的部分”最大化（高频、全阶、零通信）}。通信是分布式系统最贵、最不可靠的资源，分层架构本质是一套“通信预算最小化”的设计。这也回答了为什么 Kim--Fawcett--Hamed 能在 WiFi 上跑双四足协同——因为真正过网络的只有 MPC 层那几 KB 的接口变量。
\end{insight}

\begin{insight}[分层像公司“战略-战术-执行”三层（标边界）]\label{ins:dmpc-company}
分层架构像\textbf{一家公司的“战略层-战术层-执行层”}。高层 MPC 像高管层开会定方向（要协调、要沟通、频率低）；中层步态像部门经理把战略拆成任务（调度）；低层 WBC 像一线员工执行具体工作（按已定好的任务干、不用再开公司级大会）。\emph{像}的地方：协调成本集中在高层，执行层各自高效干活。\emph{不像}的地方：公司各层是人、靠语言沟通且会扯皮，分层控制各层是数学映射、信息流严格单向且有明确频率分工。\emph{边界}：不要延伸到“下层能否影响上层”——真实系统里 WBC 的状态会反馈给 MPC（状态估计），不是纯单向，但\emph{协调决策}是单向下传的。
\end{insight}

\subsection{WBC 为什么是 QP——多任务的优先级与折中}
\label{sec:dmpc-wbc-qp}

\cref{eq:dmpc-wbc-obj} 给的 WBC 是一个 QP（最小化加权任务误差）。但为什么 WBC 天然是优化问题、而不是一个简单的代数映射？\textbf{WBC 要同时满足多个相互竞争的任务}：跟踪 MPC 给的质心加速度、跟踪足端力 $\mathbf{f}^{\text{MPC}}$、维持机体姿态、避免关节超限、满足全身动力学。这些任务\emph{可能冲突}——比如完美跟踪质心加速度所需的接触力，可能超出摩擦锥或关节力矩限制。当任务冲突时必须\textbf{折中}，而折中天然是一个优化问题（最小化加权误差 $+$ 满足硬约束）。两种 WBC 范式：\emph{加权 WBC（weighted）}（所有任务放进一个 QP 目标、用权重 $w$ 平衡；简单、单 QP，但权重难调、软优先级）；\emph{严格优先级 WBC（hierarchical/null-space）}（高优先级任务先满足、低优先级在高优先级的零空间里做；严格优先级如“平衡 $>$ 跟踪”，但多级 QP、复杂）。严格优先级 WBC 用\textbf{零空间投影}——次要任务的解被投影到主任务的零空间，保证不破坏主任务（回顾 Khatib 操作空间控制\cite{khatib1987operational}）。

\begin{insight}[MPC 优化“未来轨迹”，WBC 优化“当前瞬时的力矩分配”]\label{ins:dmpc-wbc-underdet}
WBC 是 QP 而非代数映射，根源在于“从 SRB 降阶命令到全阶执行”是一个\textbf{欠定 $+$ 受约束}的问题——欠定（关节冗余，多种关节力矩都能实现同一质心运动）需要优化挑一个最优，受约束（摩擦锥、关节限位）需要在可行域内挑。MPC 在 SRB 层“宏观规划”，WBC 在全阶层“微观分配 $+$ 折中”。两者都是优化，但优化的对象不同：MPC 优化“未来轨迹”（跨时域），WBC 优化“当前瞬时的力矩分配”（单时刻）。这个“跨时域规划（MPC）$+$ 瞬时分配（WBC）”的分工，是分层架构的另一种解读——按“时间尺度”分。\emph{反事实}：如果 WBC 不做优化、直接用 $\boldsymbol{\tau}=\mathbf{J}^\top\mathbf{f}^{\text{MPC}}$（纯雅可比转置映射）会怎样？这忽略了全身动力学（腿部惯量、科氏力）、不处理关节限位、不能在任务冲突时折中——静态或慢速时勉强能用，动态行走时质心跟踪差、可能违反关节限制损坏电机。所以 WBC 必须是带约束的优化。
\end{insight}

\subsection{一步通信延迟、一致性协议与多机步态同步}
\label{sec:dmpc-delay-gait}

Kim--Fawcett--Hamed\cite{kim2023layered} 明确把分布式 MPC 建模为带“one-step communication delay（一步通信延迟）$+$ agreement protocol（一致性协议）”。落到分层架构里：\textbf{一步延迟}——高层 MPC 的 ADMM 协调步，agent $i$ 用的是邻居\emph{上一控制周期}算出的接口变量（耦合力/预测），不是本周期实时值；\textbf{一致性协议}——ADMM 的 $\mathbf{z}$-更新（共识平均）$+$ $\mathbf{y}$-更新（对偶上升）就是 agreement protocol 的具体实现，它逼迫各 agent 的耦合力副本在迭代中趋于一致（满足负载动力学与第三定律）。这个延迟是稳定性分析的关键隐患（\cref{sec:dmpc-topo}）——它相当于在反馈回路里插入了一个滞后，可能降低稳定裕度。

\textbf{多机步态同步——中层调度的协同维度。} \cref{sec:dmpc-dataflow} 的中层步态调度，在单体四足里只是个本地时钟（按相位决定支撑/摆动）。但在多足\emph{协同搬运}里，它多出一个被初学者忽略的协同维度：\textbf{多台机器人的步态相位要不要同步？} 考虑双足抬一根刚性杆，两台机器人各自 trot，每个周期里都有“双支撑$\to$单对角支撑”的切换、切换时质心和对地力会波动。如果两台机器人\emph{步态完全异步}（相位随机），它们对杆的力波动也不同步——杆会被两端不协调的力波动“摇晃”，内力剧烈起伏。三种步态协同策略：\emph{完全异步}（各跑各的时钟，简单、无需同步通信，但力波动不协调、杆易晃）；\emph{相位同步（in-phase）}（步态相位锁定一致，力波动同步、杆平稳，需共享相位、轻量通信）；\emph{反相同步（anti-phase）}（相位差半周期，一台双支撑时另一台单支撑、对杆的支撑可靠度互补，可能更稳）。

\begin{insight}[步态同步是多足协同特有的协调维度，且可用共识实现]\label{ins:dmpc-gait-sync}
步态同步是多足协同\textbf{特有}的协调维度——轮式/无人机编队没有“步态”，它们的执行是连续的；多足执行是\emph{离散切换}的（支撑腿在变），所以多了“切换时机要不要对齐”这个问题。这把协同从“协调连续的力/位置”扩展到“协调离散的接触事件”。从信息论看，步态相位是一个标量（相位角），共享它的通信开销极小（每周期一个浮点数），但对搬运平稳性影响巨大——“小通信、大收益”。\emph{理论-工程桥接}：步态相位同步可以用 \cref{ch:mr-consensus} 的\textbf{共识}实现！把每台机器人的步态相位 $\phi_i$ 看作要达成一致的标量，跑相位共识 $\dot\phi_i=\omega+k\sum_{j\in\mathcal{N}_i}\sin(\phi_j-\phi_i)$（Kuramoto 振子同步）就能让相位锁定。所以多足步态同步是共识理论在“离散接触事件”上的又一应用——凡是“$N$ 个量要只靠邻居达成一致”，都是共识。
\end{insight}

\begin{pitfall}[概念误区：以为分布式 WBC 也要通信]
既然 MPC 是分布式（要通信），就以为 WBC 也得和邻居交换信息。\emph{后果}：在 WBC 层加无谓的通信，徒增 $1$\,kHz 高频通信负担（网络扛不住），还引入延迟。\emph{根因}：协调已在 MPC 层通过耦合力 $\boldsymbol{\lambda}$ 完成，WBC 只需本机跟踪已协调好的命令。\emph{正确做法}：WBC 纯本地——用 MPC 传下来的 $\boldsymbol{\lambda}_i^{\text{MPC}}$ 作为已知外力，独立解本地全身 QP，零通信、可 $1$\,kHz。
\end{pitfall}

\begin{pitfall}[思维陷阱：把 MPC 的足端力直接当关节力矩发给电机]
跳过 WBC，把 MPC 输出的接触力 $\mathbf{f}_{ij}$ 直接当作控制下发。\emph{后果}：电机吃的是关节力矩不是足端力，直接发足端力是量纲/物理错误；且忽略了腿部动力学、关节限位。\emph{根因}：MPC 在 SRB 降阶层输出 SRB 级命令，必须经 WBC 映射到关节力矩并补上 SRB 忽略的全阶细节。\emph{正确做法}：MPC 足端力 $\mathbf{f}_{ij}\to$ WBC（含 $\boldsymbol{\tau}=\mathbf{J}^\top\mathbf{f}$ $+$ 全身动力学）$\to$ 关节力矩 $\boldsymbol{\tau}_i\to$ 电机。
\end{pitfall}

\begin{pitfall}[概念误区：忽略 MPC 与 WBC 的频率差异，让它们同频]
MPC 和 WBC 都跑 $1$\,kHz（或都跑 $100$\,Hz）。\emph{后果}：MPC 跑 $1$\,kHz 计算扛不住（ADMM 多轮 $+$ 通信来不及）；WBC 跑 $100$\,Hz 则平衡反应太慢、易摔。\emph{根因}：MPC 算得慢但管远（低频够用），WBC 算得快管近（必须高频）。\emph{正确做法}：MPC $\sim100$–$500$\,Hz（含通信预算），WBC $\sim1$\,kHz（纯本地高频）。MPC 的命令在 WBC 周期间保持或插值。
\end{pitfall}

\begin{exercise}
\begin{enumerate}
  \item \textbf{（分层数据流梳理）} 画出双 Go2 抬杆的完整分层数据流图，标注每一层的输入、输出、计算内容、频率、是否通信。特别标出“通信只发生在高层 MPC、且只传接口变量”和“WBC 纯本地”。
  \item \textbf{（WBC 本地 QP 实现）} 用 pinocchio $+$ OSQP 为单台四足实现一个简化 WBC：输入 MPC 给的足端力参考 $\mathbf{f}^{\text{MPC}}$ 和质心加速度目标，在全身动力学约束下求关节力矩 $\boldsymbol{\tau}$。验证：输出的 $\boldsymbol{\tau}$ 经正动力学仿真后质心加速度接近目标。然后加入耦合力 $\boldsymbol{\lambda}^{\text{MPC}}$ 项，观察它如何改变关节力矩分配。
  \item \textbf{（开放思考）} Kim--Fawcett--Hamed\cite{kim2023layered} 的低层用“全阶非线性控制器”而非 QP-WBC。查阅论文中 low-level controller 部分，对比“非线性反馈控制器”和“QP-based WBC”两种低层方案的优劣（计算量、约束处理能力、对模型误差的鲁棒性）。说明为什么两者都能纳入“低层、本地、跟踪 MPC 命令”的分层框架。
\end{enumerate}
\end{exercise}

% ════════════════════════════════════════════════════════════════
\section{通信拓扑、一致性与稳定性：把共识接到 DMPC}
\label{sec:dmpc-topo}

\paragraph{动机：ADMM 的“共识平均”到底怎么分布式算。}
回顾 \cref{sec:dmpc-admm} 的 $\mathbf{z}$-更新 $\mathbf{z}^{k+1}=\tfrac{1}{N}\sum_i(\mathbf{w}_i+\mathbf{u}_i)$——这是一个\textbf{平均}。代码里我们用了一个 for 循环求和再除以 $N$，但这\emph{隐含了一个中央节点}（谁来收集所有 $\mathbf{w}_i$ 求和？）。如果真有中央节点，那不就回到集中式的单点故障了吗？这里正是 \cref{ch:mr-consensus} 登场的地方。\cref{ch:mr-consensus} 证明了：\textbf{只靠邻居通信的共识迭代 $\mathbf{x}(k{+}1)=\mathbf{W}\mathbf{x}(k)$（$\mathbf{W}$ 双随机）能收敛到所有节点的平均值}。所以 ADMM 的 $\mathbf{z}$-更新（求平均）可以用 \cref{ch:mr-consensus} 的\textbf{分布式共识}实现——不需要中央节点，每个 agent 只和邻居交换，跑若干轮共识就能算出（近似）平均。这就把“分布式 MPC”真正变成了“无中心”。本节直接用两个结论：连续共识 $\dot{\mathbf{x}}=-\mathbf{L}\mathbf{x}$ 收敛到平均、速率由 $\lambda_2$ 决定（$e^{-\lambda_2 t}$）；离散共识要求 $\mathbf{W}$ 双随机、速率由第二大特征值模 $|\rho_2(\mathbf{W})|$ 决定。

\paragraph{反面：如果通信拓扑/质量不管会怎样。}
\emph{拓扑差 $\to$ ADMM 收敛慢}（一字链 $\lambda_2$ 小、直径大，$\mathbf{z}$-更新的分布式共识要很多轮才算准平均，外层 ADMM 也因信息传播慢而收敛轮数多）；\emph{时延 $\to$ 失稳}（\cref{ch:mr-consensus} 给了临界时延 $\tau^\star=\pi/(2\lambda_N)$，超过它共识本身就发散，叠加一步通信延迟，反馈回路滞后增大，DMPC 闭环可能失稳）；\emph{丢包 $\to$ 残差不收敛}（某轮邻居数据丢了还硬用旧值，ADMM 的对偶配平被打乱，残差降不下去）。

\begin{insight}[分布式 MPC 性能 $=$ 算法 $\times$ 通信拓扑的匹配]\label{ins:dmpc-algo-times-topo}
分布式 MPC 的性能不只取决于“算法多聪明（ADMM、$\rho$ 调得多好）”，\textbf{而是取决于“算法 $\times$ 通信拓扑的匹配”}。同一套 ADMM，在 $\lambda_2$ 大的拓扑（全连接、星形）上几轮收敛，在 $\lambda_2$ 小的拓扑（长链）上几十轮还在挣扎。通信拓扑是分布式 MPC 的“地基”，地基不好楼盖不高。
\end{insight}

\subsection{ADMM 与共识：同一枚硬币的两面}
\label{sec:dmpc-admm-consensus}

\cref{ch:mr-consensus} 讲过共识（Olfati-Saber, Murray 2004 奠基\cite{olfatisaber2004consensus}）和 ADMM（Boyd 2011 复兴\cite{boyd2011admm}）。这两条线在 $2010$ 年代汇合——人们意识到 \textbf{ADMM 的协调步本质是共识，而共识是分布式优化的特例}。“ADMM vs 平均共识谁更快”本身是个研究问题（取决于问题和拓扑，ADMM 在通信轮数上常更省但每轮计算更重）。\pz{“ADMM vs 共识谁更快取决于问题/拓扑”及源稿提及的 Teixeira 等工作，系源稿转述，未核。}

\begin{insight}[共识和 ADMM 是同一枚硬币的两面]\label{ins:dmpc-two-sides}
共识和 ADMM 不是两个独立的工具，而是\textbf{同一枚硬币的两面}——共识解决“$N$ 个标量/向量如何只靠邻居达成一致”，ADMM 解决“$N$ 个优化问题如何只靠邻居协调到联合最优”，而 ADMM 的协调步\emph{就是}一次共识。所以 \cref{ch:mr-consensus} 不是本章的“前置知识”那么简单，它是本章分布式实现的\textbf{内核}：拆掉 ADMM 的 $\mathbf{z}$-更新，里面跑的就是 \cref{ch:mr-consensus} 的 $\mathbf{x}(k{+}1)=\mathbf{W}\mathbf{x}(k)$。理解这个嵌套（外层 ADMM 优化、内层共识平均），你就理解了整个分布式多足 MPC 的运行机理。
\end{insight}

\subsection{ADMM \texorpdfstring{$\mathbf{z}$}{z}-更新的分布式共识实现}
\label{sec:dmpc-z-consensus}

把 \cref{sec:dmpc-admm} 的 $\mathbf{z}$-更新用 \cref{ch:mr-consensus} 的离散共识实现。要算的是平均 $\bar{\mathbf{v}}=\tfrac1N\sum_i\mathbf{v}_i$（$\mathbf{v}_i=\mathbf{w}_i+\mathbf{u}_i$）。每个 agent 初始化 $\mathbf{v}_i^{(0)}=\mathbf{v}_i$，跑共识迭代：
\begin{equation}
  \mathbf{v}_i^{(\ell+1)}=\sum_{j\in\mathcal{N}_i\cup\{i\}}W_{ij}\,\mathbf{v}_j^{(\ell)}\quad(\ell=0,1,\dots,L_{\text{cons}}),
  \label{eq:dmpc-inner-consensus}
\end{equation}
$\mathbf{W}$ 是双随机权重矩阵（Metropolis--Hastings 权重是常用构造，保证双随机 $+$ 连通，见 \cref{ch:mr-consensus}）。由 \cref{ch:mr-consensus} 定理，$\mathbf{v}_i^{(\ell)}\to\bar{\mathbf{v}}$，误差以 $|\rho_2(\mathbf{W})|^\ell$ 衰减。跑 $L_{\text{cons}}$ 轮后，每个 agent 都有了（近似）平均，即 $\mathbf{z}$。\textbf{两层迭代结构}：外层 ADMM（\cref{sec:dmpc-admm}），每轮 $k$ 做 $\mathbf{X}$-更新（本地 MPC，并行，无通信）、$\mathbf{z}$-更新 $=$ 内层共识（跑 $L_{\text{cons}}$ 轮 $\mathbf{v}_i\leftarrow\mathbf{W}\mathbf{v}_i$，邻居通信，得到平均 $\mathbf{z}$）、$\mathbf{y}$-更新（本地对偶上升，无通信）；内层共识轮数 $L_{\text{cons}}$ 取决于拓扑（$\lambda_2$ / $|\rho_2(\mathbf{W})|$）：$\lambda_2$ 大则 $L_{\text{cons}}$ 小。

\textbf{通信量分析}：每轮内层共识，每个 agent 向邻居发送 $\mathbf{v}_i$（接口变量，如 SRB $12$ 维 $\times$ 时域 $T=20=240$ 维 $\approx1.92$\,KB）。外层 ADMM $K$ 轮 $\times$ 内层共识 $L_{\text{cons}}$ 轮 $\times1.92$\,KB $\approx$ 几十 KB/agent/控制周期。WiFi（数十 Mbps）下完全可行。\pz{$1.92$\,KB/$240$ 维/几十 KB 的具体数值为源稿估算，未核。}

\begin{insight}[“不精确但够用”：内层 $L_{\text{cons}}$ 常取 $1$（inexact ADMM）]\label{ins:dmpc-inexact}
为什么实践中 $L_{\text{cons}}$ 常取 $1$（每个 ADMM 外层只通信一次）而非跑到内层共识完全收敛？因为外层 ADMM 本身是迭代的——即使每轮 $\mathbf{z}$-更新只做一次不精确的共识（“inexact ADMM”），外层多轮下来仍能收敛（在有界误差条件下）。这是“不精确但够用”哲学的又一体现（呼应 \cref{sec:dmpc-admm} 的固定轮数）：与其每轮把内层平均算到极致，不如内外层都“够用就停”，总通信轮数更少。Kim--Fawcett--Hamed\cite{kim2023layered} 的“一步通信延迟 $+$ agreement protocol”正是这种 $L_{\text{cons}}=1$ 的实现——每控制周期交换一次邻居预测，靠滚动 $+$ ADMM 对偶逐步配平。
\end{insight}

\subsection{通信拓扑对编队的影响}
\label{sec:dmpc-topo-effect}

把图论量对应到 DMPC 性能，见 \cref{tab:dmpc-topo-effect}。

\begin{table}[H]\centering\small
\caption{通信拓扑量对 DMPC 的影响}
\label{tab:dmpc-topo-effect}
\begin{tabular}{p{0.26\linewidth}p{0.36\linewidth}p{0.30\linewidth}}
\toprule
拓扑量（\cref{ch:mr-consensus}） & 对 DMPC 的影响 & 工程含义 \\
\midrule
代数连通性 $\lambda_2$ & 越大，内层共识与外层 ADMM 收敛越快 & 编队拓扑尽量连通（多加几条通信边） \\
直径（diameter） & 越大，信息传遍全网越慢 & 长链拓扑协调迟钝，避免 \\
临界时延 $\tau^\star=\pi/(2\lambda_N)$ & 通信时延超过它，共识/协调发散 & 高频控制下时延预算紧张，需评估 \\
丢包率 & 破坏双随机性/对偶配平，残差不收敛 & 用对丢包鲁棒的 ADMM 变体（异步、确认重传） \\
\bottomrule
\end{tabular}
\end{table}

\begin{insight}[通信图与物理队形是两张图]\label{ins:dmpc-two-graphs}
编队的\textbf{通信拓扑}和编队的\textbf{物理队形}是两张不同的图——物理上两台机器人挨得近，不代表它们通信链路好。分布式 MPC 的协调质量由\emph{通信图}的 $\lambda_2$ 决定，而不是物理队形的紧密度。所以设计编队时，通信拓扑要单独规划：宁可让物理上远的机器人之间也保持一条可靠通信边（提升 $\lambda_2$），也不要让通信图退化成长链。这是多机系统“物理层”与“信息层”解耦的又一例证。
\end{insight}

\subsection{收敛率的定量刻画与“够用就停”的本钱}
\label{sec:dmpc-quant}

\cref{sec:dmpc-admm} 说 ADMM 凸情形 $O(1/k)$、强凸 $O(1/k^2)$，\cref{sec:dmpc-z-consensus} 说内层共识误差 $\propto|\rho_2(\mathbf{W})|^\ell$。把这些拼起来，回答工程上最实际的问题：\textbf{到底要跑多少轮（外层 $K$、内层 $L_{\text{cons}}$）才“够用”？}
\begin{enumerate}
  \item \textbf{外层 ADMM 轮数 $K$}。凸情形原始残差 $\|r^k\|\sim O(1/k)$，把残差降到阈值 $\epsilon$ 需 $K\sim O(1/\epsilon)$ 轮（强凸时 $O(1/\sqrt\epsilon)$）。但实时 MPC 的关键武器是 \textbf{warm-start}——相邻控制周期的解高度相似，初始残差就很小，从 warm-start 起只需补几轮（$K=3$–$5$）。这就是 \cref{sec:dmpc-admm} 方式 C 成立的原因：不是 ADMM 本身收敛快，而是 warm-start 把“起点”挪到离最优解很近的地方。
  \item \textbf{内层共识轮数 $L_{\text{cons}}$}。误差 $\propto|\rho_2(\mathbf{W})|^{L_{\text{cons}}}$，降到 $\epsilon$ 需 $L_{\text{cons}}\sim\ln\epsilon/\ln|\rho_2(\mathbf{W})|$。$|\rho_2|$ 越接近 $1$（拓扑越差）这个数越大。但实践常取 $L_{\text{cons}}=1$，靠外层 ADMM 多轮逐步配平（inexact ADMM）。
  \item \textbf{两层轮数的乘积 $=$ 总通信成本}。真实代价是\textbf{外层 $\times$ 内层 $=K\cdot L_{\text{cons}}$ 次邻居通信}。优化目标不是单独最小化 $K$ 或 $L_{\text{cons}}$，而是最小化乘积（\cref{tab:dmpc-rounds}）。
\end{enumerate}

\begin{table}[H]\centering\small
\caption{两层轮数策略的通信权衡}
\label{tab:dmpc-rounds}
\begin{tabular}{p{0.24\linewidth}cccp{0.30\linewidth}}
\toprule
策略 & 外层 $K$ & 内层 $L_{\text{cons}}$ & 总通信 & 适用 \\
\midrule
内层精确 & 小（$3$–$5$） & 大（$\ge10$） & 中-高 & 拓扑好（$\lambda_2$ 大）、$L_{\text{cons}}$ 本就小 \\
内层不精确（inexact） & 中（$5$–$10$） & $1$ & 低 & 实时多足 MPC 推荐 \\
直接交换预测（一步延迟） & 滚动（每周期 $1$ 次） & — & 最低 & Kim--Fawcett--Hamed\cite{kim2023layered} 的做法 \\
\bottomrule
\end{tabular}
\end{table}

\begin{insight}[滚动机制是分布式 MPC 容忍双重不精确的根本]\label{ins:dmpc-rolling}
实时分布式 MPC 的全部工程智慧浓缩成一句话——\textbf{“够用就停”在两个层面同时发生：外层 ADMM 不跑到收敛（warm-start $+$ 固定 $K$），内层共识不算到精确（inexact $+$ $L_{\text{cons}}=1$）}。两个“不精确”叠加，理论上似乎应该一塌糊涂，但实践却 work——因为 MPC 是\emph{滚动}的，每个周期只取第一步控制执行，剩下的“不精确”会在下个周期被新的测量 $+$ 新的优化纠正。\emph{反事实}：如果不滚动（开环执行整条 ADMM 算出的轨迹），双重不精确的误差无人纠正，会随时间累积——编队逐渐偏离、内力逐渐失控、最终失稳。所以分布式 MPC 必须是闭环滚动的，开环用 ADMM 解一条长轨迹然后照着走是行不通的。
\end{insight}

\subsection{稳定性：递归可行性与终端代价}
\label{sec:dmpc-stability}

最后回答“分布式 MPC 闭环稳不稳”。MPC 的稳定性靠两个支柱：

\textbf{(1) 递归可行性（recursive feasibility）。} 如果当前周期 MPC 有可行解，下一周期是否仍有可行解？分布式情形更难——每个 agent 基于邻居的（滞后的）假设轨迹求解，假设与实际的偏差可能累积导致不可行。经典处理（Dunbar\cite{dunbar2006distributed}、Conte\cite{conte2016distributed} 等）：\emph{兼容性约束（compatibility constraint）}——限制 agent 的实际轨迹与它广播给邻居的假设轨迹的偏差（$\|\mathbf{x}_i^{\text{actual}}-\hat{\mathbf{x}}_i\|\le\delta$）；\emph{约束收紧（constraint tightening）}——留安全裕度，吸收分布式求解的不精确。

\textbf{(2) 终端代价/终端约束（terminal cost / terminal set）。} MPC 的有限时域要逼近无限时域的最优，靠在时域末端加终端代价 $V_f$（近似 cost-to-go）$+$ 终端约束集（保证末态落在不变集内）。分布式情形需要\emph{分布式的终端不变集计算}——这本身是难题（联合不变集难分布式算）。一些方法（如无终端代价的 DMPC，用松弛动态规划给出显式稳定性约束）绕开这个难点。

\begin{insight}[工程用长时域 $+$ warm-start $+$ 凸化替代严格的终端不变集]\label{ins:dmpc-practical-stability}
为什么很多实际的多足分布式 MPC（包括 Kim--Fawcett--Hamed\cite{kim2023layered}）不严格构造终端不变集，却仍能稳定运行？因为它们靠：(a) 足够长的预测时域（覆盖几个踏步周期，末端动态已“稳定”）；(b) warm-start 提供的好初值让每周期解接近上周期（隐式的连续性）；(c) 在 SRB 降阶层的凸化使本地 QP 总可行。严格的递归可行性 $+$ 终端不变集是\textbf{充分条件}（理论保证），工程上常用这些\emph{实践充分}的替代。这是理论与工程的典型张力：理论给最强保证但代价高（分布式不变集），工程取够用的近似。对研究者，弥合这个差距（给分布式多足 MPC 严格的稳定性证明）正是开放问题。
\end{insight}

\begin{insight}[非凸的坦白：连 ADMM 收敛性都无保证]\label{ins:dmpc-nonconvex-honest}
必须重申 \cref{sec:dmpc-admm} 的警告——多足接触约束非凸，使得连 ADMM 的\textbf{收敛性}都无理论保证，更别说闭环稳定性。当前（$2024$–$2026$）的前沿（CBF-ADMM\cite{zeng2026admmcbf}、数据驱动 DMPC\cite{fawcett2024datadriven}）仍在攻关。本章给你的是\textbf{当前最佳实践 $+$ 清醒的局限认知}：方法在仿真和受控实验中有效，但理论保证尚不完整。把这个认知带进研究，比假装一切都已证明更诚实、更有用。
\end{insight}

\subsection{兼容性约束为什么是分布式 MPC 递归可行的关键}
\label{sec:dmpc-compat}

上面提到“兼容性约束”保递归可行，但它为什么必要、怎么起作用，值得展开——这是分布式 MPC 区别于集中式 MPC 稳定性分析的\textbf{核心增量}。\emph{问题的根源：假设与实际的偏差}。在分布式 MPC 里，agent $i$ 求解时把邻居 $j$ 的轨迹当作\emph{已知的假设轨迹} $\hat{\mathbf{x}}_j$（\cref{sec:dmpc-prediction-exchange} 预测交换）。但 agent $j$ \emph{实际}执行的轨迹 $\mathbf{x}_j^{\text{actual}}$ 可能和它之前广播的 $\hat{\mathbf{x}}_j$ 不同——因为 $j$ 在新周期里基于新信息重新优化了。于是：\textbf{agent $i$ 基于“$j$ 会走 $\hat{\mathbf{x}}_j$”做的规划，可能因为“$j$ 实际走了别的”而失效}（比如 $i$ 留的避让空间不够、或编队约束因 $j$ 的偏移而违反）。\emph{集中式没有这个问题}——集中式一次解出所有 agent 的轨迹，不存在“假设别人会怎么走”。这个偏差是分布式\emph{特有}的稳定性威胁。\emph{兼容性约束的解法：限制偏差}。给每个 agent 加一条约束——它的实际轨迹不能偏离它广播的假设轨迹太多：
\begin{equation}
  \|\mathbf{x}_i^{\text{actual}}(t)-\hat{\mathbf{x}}_i(t)\|\le\delta,\quad\forall t.
  \label{eq:dmpc-compat}
\end{equation}
这样，邻居基于 $\hat{\mathbf{x}}_i$ 做的规划至多面对 $\delta$ 的偏差——只要规划时留了 $\delta$ 的裕度（约束收紧），就不会失效。$\delta$ 太大则偏差大、邻居规划要留大裕度（保守、性能差）但 agent 自己灵活；$\delta$ 太小则 agent 被迫贴着假设走（不灵活、可能不可行）但邻居几乎不用留裕度。

\begin{insight}[兼容性约束是分布式协调“信任但验证”的数学化]\label{ins:dmpc-contract}
兼容性约束的精髓是\textbf{“我承诺我不会偏离我告诉你的计划太远，于是你可以放心地基于我的计划来规划”}——它把分布式系统里“各自独立重规划”的不确定性，用一个显式的偏差上界 $\delta$ 兜住。这本质是一种\textbf{契约（contract）}：每个 agent 牺牲一点自己的灵活性（被 $\delta$ 约束），换取整个系统的可预测性（邻居能信任它的广播）。没有这个契约，分布式 MPC 的递归可行性就建立在“假设别人不变”的沙地上；有了它，可行性建立在“别人变化有界”的实地上。这是分布式协调“信任但验证（trust but verify）”哲学的数学化。
\end{insight}

\begin{insight}[兼容性约束像合租室友的“作息公约”（标边界）]\label{ins:dmpc-roommate}
兼容性约束像\textbf{合租室友之间的“作息公约”}。每个人承诺“我作息不会偏离公布的时间表超过半小时”（$\delta=$半小时），于是大家可以基于彼此的时间表安排共用厨房/卫生间（基于假设规划）。\emph{像}的地方：用一个偏差承诺让彼此的计划可互信、可协调。\emph{不像}的地方：室友靠自觉守约，分布式 MPC 把 $\delta$ 写成硬约束强制满足；且室友违约顶多尴尬，MPC 违反兼容性约束会直接导致递归不可行（系统失稳）。\emph{边界}：这个类比说明“为什么需要承诺”，但不要延伸到“承诺的代价”——室友守约几乎无成本，而 MPC 的 $\delta$ 约束实实在在牺牲了控制性能，是有代价的权衡。
\end{insight}

把兼容性约束 $+$ 约束收紧 $+$ 终端代价三者配齐，分布式 MPC 才有完整的递归可行性 $+$ 稳定性保证（Dunbar 2006\cite{dunbar2006distributed}、Conte 2016\cite{conte2016distributed} 的理论框架）。但如前所述，多足的非凸性让这套理论无法直接套用，工程上退而求其次用长时域 $+$ warm-start $+$ 凸化。理解兼容性约束的\emph{思想}（即使工程上没严格实现），能帮你判断一个分布式 MPC 在什么情况下会失稳：当某个 agent 的实际轨迹大幅偏离它的广播（如突遇大扰动急转向）、而邻居没留够裕度时。

\begin{pitfall}[概念误区：以为 ADMM 的 $\mathbf{z}$-更新（求平均）必须有中央节点]
看到 $\mathbf{z}$-更新是 $\tfrac1N\sum_i$，就以为需要一个中心收集求和。\emph{后果}：实现时真加了中央节点，退回单点故障，分布式名存实亡。\emph{根因}：平均可以用 \cref{ch:mr-consensus} 的分布式共识（只邻居通信）实现，不需中心。\emph{正确做法}：$\mathbf{z}$-更新用共识迭代 $\mathbf{v}_i\leftarrow\mathbf{W}\mathbf{v}_i$（$\mathbf{W}$ 双随机）分布式算平均，$L_{\text{cons}}$ 轮后各 agent 都有近似平均、无中心。
\end{pitfall}

\begin{pitfall}[思维陷阱：把物理队形的紧密度当作通信质量]
以为物理上挨得近的机器人通信就好、协调就快。\emph{后果}：用了物理近但通信差（如遮挡）的拓扑，ADMM 收敛慢、编队迟钝，却找不到原因。\emph{根因}：通信图和物理队形是两张图，协调质量由通信图 $\lambda_2$ 决定，与物理距离无必然关系。\emph{正确做法}：单独规划通信拓扑、保证 $\lambda_2$ 大（必要时让物理远的机器人也保持通信边），不依赖物理邻近。
\end{pitfall}

\begin{pitfall}[编程陷阱：丢包时硬用旧值且不补偿]
某轮没收到邻居数据，就用上一轮的旧值继续 ADMM，不做任何标记或补偿。\emph{后果}：对偶配平被打乱，残差降不到 $0$，编队有持续偏差，且时好时坏（取决于丢包随机性）。\emph{根因}：ADMM 的收敛依赖各 agent 数据的一致更新，无补偿的丢包破坏这个一致性。\emph{正确做法}：用对丢包鲁棒的变体（异步 ADMM 在有界延迟下仍收敛、确认重传、或丢包时跳过该邻居并重归一化权重），监控残差以发现异常。
\end{pitfall}

\begin{pitfall}[概念误区：以为长预测时域 $+$ warm-start 就等于严格稳定性保证]
因为实践中长时域 $+$ warm-start 跑得稳，就当成有了稳定性证明。\emph{后果}：在未测试的工况（大扰动、新步态切换）下可能失稳，因为没有理论保证兜底。\emph{根因}：长时域 $+$ warm-start 是实践充分条件，不是理论充分条件；非凸下尤其无保证。\emph{正确做法}：清醒区分“实践有效”和“理论保证”。关键应用要做充分的工况测试 $+$ 安全监控（如 CBF 安全层，\cref{sec:sm-cbf-game}），不把实践稳定误当数学证明。
\end{pitfall}

\begin{exercise}
\begin{enumerate}
  \item \textbf{（拓扑 $\to$ 收敛轮数）} 对 $N=5$ 的三种拓扑（全连接、星形、一字链），用 Metropolis 权重构造双随机 $\mathbf{W}$，计算各自的 $|\rho_2(\mathbf{W})|$。估算把平均算到 $1\%$ 精度各需多少轮共识 $L_{\text{cons}}$（$|\rho_2|^L\le0.01$）。验证：全连接/星形几轮够，长链要十几轮。
  \item \textbf{（嵌套迭代实现）} 在双积分器 ADMM 编队基础上，把 $\mathbf{z}$-更新从“中央求平均”改成“分布式共识迭代”（$L_{\text{cons}}=1,3,10$）。对比：$L_{\text{cons}}$ 越小，外层 ADMM 需越多轮才收敛（不精确 ADMM）；$L_{\text{cons}}$ 大则内层通信开销大。找到总通信轮数（外层 $\times$ 内层）最优的组合。
  \item \textbf{（时延稳定性）} 用临界时延 $\tau^\star=\pi/(2\lambda_N)$，对一字链拓扑（$N=5$）计算 $\lambda_N$ 和 $\tau^\star$。说明：如果叠加一步通信延迟（$=$ 一个控制周期），在什么控制频率下这个延迟会接近 $\tau^\star$ 而威胁稳定？
  \item \textbf{（跨章综合题 $+$ 开放）} 综合 \cref{ch:mr-consensus}（拓扑/$\lambda_2$/共识/ADMM/时延）与本章（DMPC）：为一个真实的 $4$ 台 Go2 抬板编队设计完整的通信与协调方案——选什么通信拓扑（理由）、ADMM 外层轮数 $K$ 与内层共识轮数 $L_{\text{cons}}$ 如何权衡、如何处理一步延迟与可能的丢包、稳定性靠什么保障。写成一段 $200$–$300$ 字的设计文档。
\end{enumerate}
\end{exercise}

% ════════════════════════════════════════════════════════════════
\section*{本章常见误解汇总}
\label{sec:dmpc-misconceptions}
\addcontentsline{toc}{section}{本章常见误解汇总}

\begin{table}[H]\centering\small
\caption{分布式 MPC 与多足编队常见误解辨正}
\label{tab:dmpc-misconceptions}
\begin{tabular}{p{0.34\linewidth}p{0.50\linewidth}p{0.08\linewidth}}
\toprule
误解 & 正确理解 & 出处 \\
\midrule
分布式 MPC 总比集中式好 & 用（较小的）最优性 gap 换可扩展性 $+$ 鲁棒性 $+$ 无单点；$N$ 小且中央可靠时集中式更优\rebuilt{原文“$<3\%$ gap”系杜撰，\cite{kim2023layered} 仅称“性能相近”、未给数值，已删此具体数。} & \cref{sec:dmpc-why} \\
“动力学解耦”等于“问题可分” & 多足编队动力学解耦但约束/代价耦合；耦合恰来自编队约束，必须协调 & \cref{sec:dmpc-why} \\
预测交换里能拿到邻居实时值 & 真实系统有一步通信延迟，用的是邻居上一轮的预测 & \cref{sec:dmpc-why,sec:dmpc-layered} \\
编队就是给每个机器人一个绝对目标点 & 编队是相对几何约束，整体位姿是自由（或由领航者定）的自由度 & \cref{sec:dmpc-formation} \\
distance-based 代价直接用欧氏范数 & 范数在重合点不可微，必须用平方距离 & \cref{sec:dmpc-formation} \\
四足编队硬约束瞬时质心相对位置 & 质心随踏步起伏，硬约束会与平衡打架；用软约束 $+$ 频段分离 & \cref{sec:dmpc-formation,sec:dmpc-srb} \\
对偶分解要松弛所有约束 & 只松弛耦合约束，本地约束保留在各子问题里 & \cref{sec:dmpc-dual} \\
对偶变量 $\boldsymbol{\lambda}$ 是自由超参 & $\boldsymbol{\lambda}$ 是有物理含义的影子价格/约束力，可 warm-start & \cref{sec:dmpc-dual} \\
ADMM 是全新方法 & ADMM $=$ 对偶分解 $+$ 增广项 $+$ 辅助变量，修复对偶分解的三个弱点 & \cref{sec:dmpc-admm} \\
$\mathbf{z}$-更新漏对偶项无所谓 & 必须平均 $\mathbf{w}_i+\mathbf{u}_i$，漏 $\mathbf{u}_i$ 收敛到错误点 & \cref{sec:dmpc-admm} \\
$\rho$ 越大收敛越快 & $\rho$ 是“一致 vs 自主”权重，过大让本地代价被淹没、QP 病态 & \cref{sec:dmpc-admm} \\
ADMM 在非凸下也保证收敛 & 多足接触非凸，ADMM 无收敛保证；靠凸化 $+$ warm-start $+$ 固定轮数 & \cref{sec:dmpc-admm,sec:dmpc-topo} \\
耦合力 $\boldsymbol{\lambda}$ 是独立控制输入 & $\boldsymbol{\lambda}$ 是约束力，由刚性约束 $+$ 动力学解出，受严格限制 & \cref{sec:dmpc-srb} \\
摩擦锥用精确二阶锥 & 二阶锥使问题变 SOCP（慢），用金字塔线性近似保持 QP & \cref{sec:dmpc-srb} \\
分布式 WBC 也要通信 & 协调已在 MPC 层由 $\boldsymbol{\lambda}$ 完成，WBC 纯本地、零通信 & \cref{sec:dmpc-layered} \\
MPC 足端力可直接发给电机 & 电机吃关节力矩，须经 WBC 映射 $+$ 补全阶细节 & \cref{sec:dmpc-layered} \\
ADMM 的平均必须有中央节点 & 用分布式共识算平均，无需中心 & \cref{sec:dmpc-topo} \\
物理队形紧密 $=$ 通信质量好 & 通信图与物理队形是两张图，协调质量由通信图 $\lambda_2$ 决定 & \cref{sec:dmpc-topo} \\
长时域 $+$ warm-start $=$ 稳定性证明 & 那是实践充分条件，非理论保证；非凸下尤其无保证 & \cref{sec:dmpc-topo} \\
\bottomrule
\end{tabular}
\end{table}

% ════════════════════════════════════════════════════════════════
\section*{本章小结}
\label{sec:dmpc-summary}
\addcontentsline{toc}{section}{本章小结}

本章把“一群会走路的机器人如何只靠彼此通信、在毫秒级控制周期内联合保持队形协同前进”这个问题，从动机到算法到多足落地，完整走了一遍。主线是：\textbf{集中式不可扩展（\cref{sec:dmpc-why}）$\to$ 编队建模制造耦合（\cref{sec:dmpc-formation}）$\to$ 对偶分解拆耦合但慢（\cref{sec:dmpc-dual}）$\to$ ADMM 又快又稳地拆（\cref{sec:dmpc-admm}）$\to$ SRB 网络把多足装进框架（\cref{sec:dmpc-srb}）$\to$ 分层控制接到关节力矩（\cref{sec:dmpc-layered}）$\to$ 通信拓扑与稳定性收口（\cref{sec:dmpc-topo}）}。三条贯穿全章的认知主线：(1) \textbf{分布式是权衡，不是免费午餐}——用最优性 gap 换可扩展性、鲁棒性、模块化；(2) \textbf{ADMM 的协调步本质是共识}——\cref{ch:mr-consensus} 不是前置而是内核，外层 ADMM 优化分解、内层共识算平均，两者吃同一张通信图的红利（$\lambda_2$）；(3) \textbf{多足的难点在频段分离与非凸}——编队/协同是低频群体行为（SRB 层 $+$ 通信协调），平衡是高频个体行为（WBC 层 $+$ 本地）；接触约束非凸使理论保证失效，工程靠凸化 $+$ warm-start $+$ 够用就停。集中式与分布式的完整对照见 \cref{tab:dmpc-cent-vs-dist}（源稿对比，数值未独立复算）。

\rebuilt{下表的数值版（$15$\,ms/$\sim$$6$\,ms$\times2$/$<3\%$ gap/$O(N^3)$）经核与所引 Kim--Fawcett--Hamed (T-RO 2023, arXiv:2211.06913) 源文\emph{矛盾}，系源稿杜撰：源文实测集中式 MPC $1.38$\,ms、分布式 $0.41$\,ms（“solve time reduced by 70\%”），并无“$<3\%$ gap”（仅定性“performance similar”），决策变量随 $N$\emph{线性}（集中 $245$、分布 $125$），从未出现 $O(N^3)$。故原数值行已注释停用，下面给按源文校正的方向版（仅供修订参考、仍需复核）。}
% —— 错因：以下原表数值与 kim2023layered 源文矛盾（详见上方 \rebuilt 标注），整表注释停用、待修，不就地重建为“正确版”——
% \begin{table}[H]\centering\small
% \caption{集中式 MPC vs 分布式 ADMM-MPC（源稿对比，待核）}
% \label{tab:dmpc-cent-vs-dist}
% \begin{tabular}{p{0.18\linewidth}p{0.24\linewidth}p{0.28\linewidth}p{0.20\linewidth}}
% \toprule
% 指标 & 集中式 MPC & 分布式 ADMM-MPC & 解读 \\
% \midrule
% 轨迹最优性 & 最优（基准） & $<3\%$ gap & 分布式近似但够好 \\           % <3% gap：源文无此数，仅“similar”
% 计算时间 & $\sim15$\,ms（单核全包） & $\sim6$\,ms $\times2$（并行） & 分布式并行更快 \\  % 源文：集中 1.38ms、分布 0.41ms
% 通信 & 集中式上下行 $\propto N$ & 几 KB/agent/周期 ADMM & 分布式通信轻、无单点 \\          % “几KB”源文未给具体数
% 鲁棒性 & 单点故障 & 无中心、抗节点失效 & 分布式胜 \\
% $N$ 扩展性 & $O(N^3)$ 爆炸 & 近线性（并行） & 分布式可扩展 \\                 % O(N^3) 源文未出现，变量随 N 线性
% \bottomrule
% \end{tabular}
% \end{table}
\begin{table}[H]\centering\small
\caption{集中式 MPC vs 分布式 ADMM-MPC（按 \cite{kim2023layered} 源文校正，仍待复核）}
\label{tab:dmpc-cent-vs-dist}
\begin{tabularx}{\linewidth}{p{0.18\linewidth}LLL}
\toprule
指标 & 集中式 MPC & 分布式 ADMM-MPC & 解读（源文依据） \\
\midrule
轨迹最优性 & 最优（基准） & 与集中式“性能相近”（源文未给具体 gap） & 仅定性表述，勿引“$<3\%$” \\
计算时间 & $1.38$\,ms（源文实测） & $0.41$\,ms（源文实测，求解时间降约 $70\%$） & 分布式更快源自降维 \\
通信 & 集中式上下行 $\propto N$ & 仅交换乘子 $\boldsymbol{\lambda}$、一步延迟（源文未给字节数） & 分布式通信轻、无单点 \\
鲁棒性 & 单点故障 & 无中心、抗节点失效 & 分布式胜 \\
$N$ 扩展性 & 决策变量随 $N$ 线性（源文 $245$@双机） & 决策变量随 $N$ 线性（源文 $125$@单机） & 源文未出现 $O(N^3)$ \\
\bottomrule
\end{tabularx}
\end{table}

\pz{上表为按 \cite{kim2023layered} 源文 PDF 校正后的方向版（集中 $1.38$\,ms、分布 $0.41$\,ms、降约 $70\%$、变量随 $N$ 线性）；原源稿数值版（$15$/$6$\,ms、$<3\%$ gap、$O(N^3)$）见紧邻上方注释块，已停用待修。两版均建议下游再次联网/复算确认。}

\subsection*{术语速查表}
\begin{table}[H]\centering\small
\caption{本章术语速查}
\label{tab:dmpc-glossary}
\begin{tabular}{p{0.28\linewidth}p{0.66\linewidth}}
\toprule
术语 & 一句话定义 \\
\midrule
分布式 MPC（DMPC） & 每个 agent 解本地 MPC $+$ 邻居交换信息协调，替代集中式联合 MPC \\
预测交换 & agent 广播自己的预测轨迹，邻居用它构造本地编队项 \\
对偶分解 & 拉格朗日松弛耦合约束使问题可分，用对偶变量（价格）协调 \\
影子价格 & 对偶变量的经济学解释——维持约束的边际代价 \\
ADMM & 交替方向乘子法 $=$ 对偶分解 $+$ 增广二次项 $+$ 辅助变量 \\
consensus-ADMM & 用全局共识变量 $\mathbf{z}$ 把多边耦合化为对中心的单边耦合 \\
罚参数 $\rho$ & ADMM 中“全局一致 vs 本地自主”的权重 \\
残差平衡 & 让原始/对偶残差量级相当的自适应 $\rho$ 规则 \\
SRB 网络 & $N$ 个单刚体 $+$ 共享负载 $+$ 刚性耦合约束的互联模型 \\
耦合力 $\boldsymbol{\lambda}$ & 维持刚性连接的约束力（拉格朗日乘子），协同的物理本质 \\
摩擦锥金字塔 & 用线性不等式近似摩擦锥，保持 QP 结构 \\
单边接触 & 脚只能推不能拉地面（$f_z\ge0$）；摆动腿力为零 \\
分层控制 & 高层 MPC（协调、通信）$\to$ 中层步态 $\to$ 低层 WBC（执行、本地） \\
分布式 WBC & 低层全身控制纯本地、无通信（协调已由 MPC 的 $\boldsymbol{\lambda}$ 完成） \\
兼容性约束 & 限制实际轨迹与广播假设轨迹的偏差，保递归可行 \\
递归可行性 & 当前周期可行则下周期仍可行 \\
终端代价/终端集 & MPC 末端加的代价/约束，保有限时域逼近无限时域稳定 \\
\bottomrule
\end{tabular}
\end{table}

\subsection*{定理与公式速查表}
\begin{table}[H]\centering\small
\caption{本章关键公式速查}
\label{tab:dmpc-formulas}
\begin{tabular}{p{0.30\linewidth}p{0.40\linewidth}p{0.20\linewidth}}
\toprule
名称 & 核心式 & 出处 \\
\midrule
集中式联合 MPC & $\min\sum_i J_i+J_{\text{form}}$ s.t.\ 动力学/编队/接触 & \cref{eq:dmpc-central-obj} \\
编队代价 $=$ 加权 Laplacian & $J_{\text{form}}=\mathbf{p}^\top \mathbf{L}_w\mathbf{p}$ & \cref{eq:dmpc-form-laplacian} \\
对偶分解价格更新 & $\boldsymbol{\lambda}^{k+1}=\boldsymbol{\lambda}^k+\alpha^k(\sum_i \mathbf{A}_i\mathbf{X}_i^*-\mathbf{b})$ & \cref{eq:dmpc-price-update} \\
ADMM $\mathbf{X}$-更新 & $\mathbf{X}_i^{k+1}=\argminop J_i+\tfrac{\rho}{2}\|\mathbf{w}_i-\mathbf{z}_i^k+\mathbf{u}_i^k\|^2$ & \cref{eq:dmpc-x-update} \\
ADMM $\mathbf{z}$-更新 & $\mathbf{z}^{k+1}=\tfrac{1}{|\mathcal{I}_z|}\sum(\mathbf{w}_i+\mathbf{u}_i^k)$ & \cref{eq:dmpc-z-update} \\
ADMM $\mathbf{y}$-更新 & $\mathbf{u}_i^{k+1}=\mathbf{u}_i^k+\mathbf{w}_i-\mathbf{z}_i^{k+1}$ & \cref{eq:dmpc-y-update} \\
残差平衡自适应 $\rho$ & 按 $\|r^k\|\gtrless\nu\|s^k\|$ 增/减 $\rho$ & \cref{eq:dmpc-residual-balance} \\
SRB 网络牛顿方程 & $m_i\ddot{\mathbf{p}}_i=\sum_j\mathbf{f}_{ij}+\boldsymbol{\lambda}_i+m_i\mathbf{g}$ & \cref{eq:dmpc-srb-newton} \\
约束力解析解 & $\boldsymbol{\lambda}=-(\mathbf{J}_\Phi \mathbf{M}^{-1}\mathbf{J}_\Phi^\top)^{-1}(\mathbf{J}_\Phi \mathbf{M}^{-1}\mathbf{Q}+\dot{\mathbf{J}}_\Phi\dot{\mathbf{q}})$ & \cref{eq:dmpc-constraint-force} \\
摩擦锥金字塔 & $|f_x|\le\mu f_z,\ |f_y|\le\mu f_z$ & \cref{eq:dmpc-pyramid} \\
分布式 WBC 全身动力学 & $\mathbf{M}_i\ddot{\mathbf{q}}_i+\mathbf{h}_i=\mathbf{S}_i^\top\boldsymbol{\tau}_i+\mathbf{J}_{c,i}^\top\mathbf{f}_i+\mathbf{J}_{\text{att},i}^\top\boldsymbol{\lambda}_i^{\text{MPC}}$ & \cref{eq:dmpc-wbc-dyn} \\
内层共识平均 & $\mathbf{v}_i^{(\ell+1)}=\sum_j W_{ij}\mathbf{v}_j^{(\ell)}$，误差 $\propto|\rho_2(\mathbf{W})|^\ell$ & \cref{eq:dmpc-inner-consensus} \\
兼容性约束 & $\|\mathbf{x}_i^{\text{actual}}-\hat{\mathbf{x}}_i\|\le\delta$ & \cref{eq:dmpc-compat} \\
\bottomrule
\end{tabular}
\end{table}

\subsection*{符号表}
\begin{table}[H]\centering\small
\caption{本章主要符号}
\label{tab:dmpc-symbols}
\begin{tabular}{p{0.16\linewidth}p{0.78\linewidth}}
\toprule
符号 & 含义 \\
\midrule
$N$, $T$ & agent 数；MPC 预测时域步数 \\
$i,j$；$\mathcal{N}_i$ & agent 索引；agent $i$ 的通信邻居集 \\
$\mathbf{x}_i,\mathbf{u}_i$ & agent $i$ 的状态（SRB $12$ 维）/控制（足端力堆叠） \\
$\mathbf{X}_i,\mathbf{U}_i$；$J_i$ & agent $i$ 在时域上的状态/控制轨迹；本地代价 \\
$\mathbf{p}_i$；$\mathbf{d}_{ij}$ & agent $i$ 质心位置；期望 $i$-$j$ 相对位移 \\
$\mathbf{A}_i,\mathbf{b}$ & 耦合约束 $\sum_i \mathbf{A}_i\mathbf{X}_i=\mathbf{b}$ 的系数 \\
$\boldsymbol{\lambda}$ & 拉格朗日乘子（\cref{sec:dmpc-dual} 对偶变量、\cref{sec:dmpc-srb} 刚性耦合约束力） \\
$g(\boldsymbol{\lambda})$ & 对偶函数 \\
$\mathbf{z}$；$\mathbf{w}_i$ & ADMM 全局共识变量；agent $i$ 的接口变量 \\
$\mathbf{y}_i$；$\mathbf{u}_i$ & ADMM 对偶变量；缩放对偶变量 $\mathbf{u}_i=\mathbf{y}_i/\rho$ \\
$\rho$；$r^k,s^k$ & ADMM 罚参数；第 $k$ 轮原始/对偶残差 \\
$\mathbf{f}_{ij}$；$\mu$ & agent $i$ 第 $j$ 条腿的接触力；摩擦系数 \\
$\mathbf{p}_i^{\text{att}}$；$\mathbf{J}_\Phi$ & agent $i$ 在共享负载上的附着点；约束 Jacobian \\
$\mathbf{L}$, $\lambda_2$；$\mathbf{W}$, $\rho_2(\mathbf{W})$ & 通信图 Laplacian、代数连通性；离散共识权重、第二大特征值模 \\
$K$；$L_{\text{cons}}$ & 外层 ADMM 轮数；内层共识轮数 \\
\bottomrule
\end{tabular}
\end{table}

\subsection*{知识点总表}
\begin{table}[H]\centering\small
\caption{本章知识点总表}
\label{tab:dmpc-knowledge}
\begin{tabularx}{\linewidth}{p{0.05\linewidth}p{0.22\linewidth}Lp{0.10\linewidth}p{0.04\linewidth}}
\toprule
编号 & 知识点 & 核心要点 & 对应节 & 难度 \\
\midrule
K1 & 集中式三重瓶颈 & $O(N^3)$ 爆炸 $+$ 单点故障 $+$ 不可模块化 & \cref{sec:dmpc-why} & $\star\star$ \\
K2 & 预测交换范式 & 交换预测轨迹，编队耦合退化为本地项；软约束、有一步延迟 & \cref{sec:dmpc-why} & $\star\star\star$ \\
K3 & 编队三范式 & position/displacement/distance：信息越少、自由度越多、越非凸 & \cref{sec:dmpc-formation} & $\star\star\star$ \\
K4 & 编队 $=$ 加权共识 & $J_{\text{form}}=\mathbf{p}^\top \mathbf{L}_w\mathbf{p}$；编队是“带偏移的共识” & \cref{sec:dmpc-formation} & $\star\star\star$ \\
K5 & 对偶分解 & 松弛耦合约束 $\to$ 可分 $\to$ 价格迭代；慢、脆、中间不可行 & \cref{sec:dmpc-dual} & $\star\star\star\star$ \\
K6 & 影子价格 $=$ 约束力 & 对偶变量在多足里是物理可测的耦合力 $\boldsymbol{\lambda}$ & \cref{sec:dmpc-dual,sec:dmpc-srb} & $\star\star\star$ \\
K7 & consensus-ADMM 三步 & $\mathbf{X}$（本地 MPC$+$惩罚）/$\mathbf{z}$（共识平均）/$\mathbf{y}$（对偶上升） & \cref{sec:dmpc-admm} & $\star\star\star\star$ \\
K8 & $\rho$ 与残差平衡 & $\rho$ 是一致-自主权重；残差平衡自适应；凸 $O(1/k)$ & \cref{sec:dmpc-admm} & $\star\star\star$ \\
K9 & 非凸无保证 & 接触约束非凸 $\to$ ADMM 无收敛保证 $\to$ 凸化$+$warm-start & \cref{sec:dmpc-admm,sec:dmpc-topo} & $\star\star\star\star$ \\
K10 & SRB 网络模型 & $N$ SRB $+$ 负载 $+$ 刚性约束 $\to$ 耦合力 $\boldsymbol{\lambda}$ & \cref{sec:dmpc-srb} & $\star\star\star\star$ \\
K11 & 约束力求解 & $\boldsymbol{\lambda}=-(\mathbf{J}\mathbf{M}^{-1}\mathbf{J}^\top)^{-1}(\mathbf{J}\mathbf{M}^{-1}\mathbf{Q}+\dot{\mathbf{J}}\dot{\mathbf{q}})$，约束微分两次 & \cref{sec:dmpc-srb} & $\star\star\star\star$ \\
K12 & 接触约束三件套 & 摩擦锥金字塔 $+$ 单边接触 $+$ 踏步时序；固定序列$\to$凸 QP & \cref{sec:dmpc-srb} & $\star\star\star$ \\
K13 & 频段分离 & 编队低频（质心层）$+$ 平衡高频（WBC 层） & \cref{sec:dmpc-formation,sec:dmpc-srb} & $\star\star\star$ \\
K14 & 分层数据流 & MPC（通信）$\to$步态$\to$WBC（本地）；通信预算最小化 & \cref{sec:dmpc-layered} & $\star\star\star$ \\
K15 & 分布式 WBC 纯本地 & 协调已由 MPC 的 $\boldsymbol{\lambda}$ 完成，WBC 零通信 & \cref{sec:dmpc-layered} & $\star\star\star$ \\
K16 & ADMM 协调步 $=$ 共识 & $\mathbf{z}$-更新用分布式共识实现，无需中心 & \cref{sec:dmpc-topo} & $\star\star\star\star$ \\
K17 & 两层迭代 & 外层 ADMM 优化、内层共识平均；$L_{\text{cons}}$ vs $K$ 权衡 & \cref{sec:dmpc-topo} & $\star\star\star$ \\
K18 & 拓扑影响 & $\lambda_2$/直径/时延/丢包影响收敛与稳定；通信图 $\ne$ 物理队形 & \cref{sec:dmpc-topo} & $\star\star\star$ \\
K19 & DMPC 稳定性支柱 & 递归可行性（兼容性约束）$+$ 终端代价；分布式不变集难算 & \cref{sec:dmpc-topo} & $\star\star\star\star$ \\
\bottomrule
\end{tabularx}
\end{table}

% ════════════════════════════════════════════════════════════════
\section*{延伸阅读与后续关系}
\label{sec:dmpc-further}
\addcontentsline{toc}{section}{延伸阅读与后续关系}

\textbf{核心论文（必读）}：Kim, Fawcett, Kamidi, Ames, Hamed, ``Layered Control for Cooperative Locomotion of Two Quadrupedal Robots: Centralized and Distributed Approaches''（T-RO 2023\cite{kim2023layered}，源稿标 arXiv:2211.06913）——本章主线，互联 SRB 网络 $+$ 集中/分布 MPC $+$ 分层控制的完整范本；Boyd, Parikh, Chu, Peleato, Eckstein, ``Distributed Optimization and Statistical Learning via ADMM''（Found.\ Trends ML 2011\cite{boyd2011admm}）——ADMM 与 consensus-ADMM 的权威教程，\cref{sec:dmpc-dual,sec:dmpc-admm} 的理论源头。

\textbf{分布式 MPC 基础}：Dunbar \& Murray（Automatica 2006\cite{dunbar2006distributed}）——预测交换范式奠基；Keviczky, Borrelli, Balas（Automatica 2008\cite{keviczky2008decentralized}）——DMPC 编队系统化\pz{此处 venue/作者两误：本键实为 IEEE T-CST 16(1):19--33, 2008（非 Automatica），作者应作 Keviczky, Borrelli, Fregene, Godbole, Balas。标题含“dynamically decoupled”者为同作者 Automatica 2006 那篇，勿混。}；Conte, Jones, Morari, Zeilinger（Automatica 2016\cite{conte2016distributed}）——DMPC 稳定性综合（终端代价、递归可行）。

\textbf{多足协同前沿}：De Vincenti \& Coros（RSS 2023\cite{devincenti2023centralized}）——集中式 SQP $+$ $\SO(3)$ 李群积分，对照集中式方案；Turrisi 等（IROS 2024\cite{turrisi2024pacc}）——被动臂方案，免去刚性耦合建模；Fawcett, Amanzadeh, Hamed 等（2024\cite{fawcett2024datadriven}）——数据驱动分布式预测控制；Imran, Hamed 等（2025\cite{imran2025obstacle}）与 ADMM-CBF（arXiv\cite{zeng2026admmcbf}）——障碍环境/安全的多足分布式 MPC。\pz{Fawcett 这条年份误：实为 Proc.\ IEEE ICRA 2023, pp.\,9924--9930（arXiv:2211.06917），非 2024，键宜改 fawcett2023distributed。其余 De Vincenti--Coros (RSS 2023)、Turrisi (IROS 2024)、Imran (ASME J.\ Dyn.\ Sys.\ 2025)、zeng2026admmcbf (arXiv:2603.19170) 经核 venue/年份属实。}

\textbf{后续关系}：本章管“一起走到哪”（运动协调），\cref{sec:dmpc-srb} 的耦合力 $\boldsymbol{\lambda}$ 是通往\emph{协同搬运}章（grasp matrix、内力分配，管“一起使多大劲”）的桥梁；本章的分布式协调骨架（共识接口、ADMM）继承自 \cref{ch:mr-consensus}，并与 \cref{sec:sm-cbf-game} 的安全证书（CBF）在“硬安全 $+$ 软任务”处汇合（Imran--Hamed\cite{imran2025obstacle} 的 CBF-DMPC）；$\SO(3)$ 姿态的精确几何处理见 \nameref{ch:lie}（本章降阶 MPC 用欧拉角小角度线性化，大机动时应改用李群积分）。
